% !TeX program = lualatex
% D033 editable mathematical derivative. Acceptance is established only by the separate corpus gate.
\documentclass[11pt]{article}
\usepackage{fontspec,amsmath,unicode-math}
\setmainfont{STIXTwoText}[Extension=.otf,UprightFont=*-Regular,ItalicFont=*-Italic,BoldFont=*-Bold,BoldItalicFont=*-BoldItalic]
\setmathfont{STIXTwoMath-Regular.otf}
\usepackage[main=english]{babel}
\usepackage[a4paper,left=25mm,right=25mm,top=22mm,bottom=24mm,headheight=15pt]{geometry}
\usepackage{graphicx,adjustbox,needspace,parskip,fancyhdr,microtype,hyperref,tikz-cd}
\hypersetup{hidelinks,pdftitle={Values of Abelian L-functions at Negative Integers over Totally Real Fields},pdfauthor={Pierre Deligne and Kenneth A. Ribet},pdfsubject={D033 bilingual scholarly edition}}
\pdfvariable suppressoptionalinfo 767\relax
\setlength{\parindent}{0pt}\setlength{\parskip}{.45\baselineskip}\setlength{\emergencystretch}{3em}\raggedbottom
\pagestyle{fancy}\fancyhf{}\fancyhead[L]{\small D033: Deligne--Ribet}\fancyhead[R]{\small English source edition}\fancyfoot[C]{\thepage}
\newcommand{\DAnchor}[2]{\clearpage\phantomsection\pdfbookmark[1]{Authority #1 / printed #2}{d033p#1}{\small\textsc{Authority #1 / printed #2}\par}\medskip}
\newcommand{\DAppAnchor}[2]{\Needspace{5\baselineskip}\phantomsection\pdfbookmark[1]{Authority #1 / printed #2}{d033p#1}\subsection*{Authority #1 / printed #2}}
\newcommand{\DFormula}[1]{\par\smallskip\noindent\begin{adjustbox}{max width=\linewidth}$\displaystyle #1$\end{adjustbox}\par\smallskip}
\newcommand{\DFallback}[2]{\par\medskip\begin{center}\includegraphics[width=.75\linewidth,height=.55\textheight,keepaspectratio]{../assets/#1}\par{\small #2}\end{center}\medskip}
\begin{document}

\DAnchor{1}{227}
Inventiones math. 59, 227--286 (1980)\par
\par
Values of Abelian L-functions at Negative Integers\par
over Totally Real Fields\par
\par
Pierre Deligne and Kenneth \(A\). Ribet\textsuperscript{*}\par
\par
Institut des Hautes Etudes Scientifiques, 35, route de Chartres,\par
F-91440 Bures-sur-Yvette, France\par
\par
\subsection*{§0}
\par
The purpose of this paper is to interpolate p-adically the values at negative integers of abelian L-functions for totally real fields, by using the method of Hilbert modular forms. We recall that the \(p\)-adic study of L-values by means of \(p\)-adic modular forms was initiated by Serre in his Antwerp paper [33] and was suggested by Siegel’s use of Eisenstein series to prove the rationality of L-values [38]. Serre used modular forms of one variable in his theory, and constructed \(p\)-adic zeta functions, as well as \(p\)-adic L-functions for powers of the Teichmüller character, over totally real fields. Soon after, Katz (unpublished) and Queen [27] observed that Serre’s theory applied more generally to the construction of \(p\)-adic L-functions attached to abelian characters over totally real fields. Here the point was to introduce \(p\)-adic modular forms with level, i.e., on \(\Gamma _{1}(N)\). Already in Serre’s paper, it was clear that the use of forms of one variable was insufficient for a complete theory. For example, one could not rule out the possibility of a pole at \(s=1\) for the \(p\)-adic L-functions attached to certain non-trivial powers of the Teichmüller character. At the close of [33], Serre suggested that one develop a theory of \(p\)-adic Hilbert modular forms to overcome these difficulties.\par
\par
This problem was taken up by the first author of this paper, who described in letters to Serre [5] how the required congruences among L-values would follow from a (conjectural) theory of \(p\)-adic Hilbert forms. This theory was in turn seen as a consequence of a construction over \(\mathbb{Z}\) of certain Hilbert-Blumenthal moduli schemes, whose fibres in characteristic \(p\) were further conjectured to be geometrically irreducible. The geometric irreducibility amounted to the surjectivity of a certain Galois representation, and was already known in the case of one-variable modular forms (i.e., with the totally real field taken to be \(\mathbb{Q})\) [12].\par
\par
The second author entered the picture by proving (in 1974) the conjectured irreducibility from the Galois-representation point of view [29]. Strictly speaking, the schemes being proved irreducible did not yet exist.\par
\par
\(^{\ast}\) Permanent address: Berkeley Math. Department, Berkeley, CA 94720, USA\par
\par
\DAnchor{2}{228}
In 1976, however, Rapoport constructed in his thesis [28] the schemes in question. The work of Rapoport contains the foundational material necessary to the theory of \(p\)-adic Hilbert forms, and in particular the all-important \(q\)-expansion principle. Thus, among other things, Rapoport gave algebraic meaning to the \(q\)-expansion of a Hilbert modular form, so that one could define the \(q\)-expansion of a Hilbert modular form over an arbitrary ring \(R\), at an arbitrary unramified cusp over \(R\). The \(q\)-expansion principle then states that each \(q\)-expansion of a Hilbert modular form \(F\) over a ring \(R\) completely determines that form, and furthermore that \(F\) is a form over a subring \(R_{0}\) of \(R\) if and only if the \(q\)-expansion coefficients of \(F\) lie in \(R_{0}\). Rapoport’s thesis completed the program begun by Serre and Deligne and thus removed the last tangible obstacle to the writing of the present paper.\par
\par
Before summarizing the contents of this paper, we wish to mention related work by other authors. First of all, the “moduli” approach to modular forms and its application to the construction of \(p\)-adic L-functions has been developed in a series of papers by Katz [15--19]. Especially, in [19], Katz uses the theory of \(p\)-adic Hilbert modular forms in a construction of “2n-variable” \(p\)-adic L-functions for CM fields. His paper begins with an excellent summary of the theory, on which our §5 is modeled.\par
\par
Secondly, Barsky [1] and Cassou-Noguès [2] have given alternate approaches to \(p\)-adic L-functions over totally real fields, based on the explicit formulas of Shintani [37]. The results of Barsky/Cassou-Noguès and those of the present paper were compared in [30]. The two sets of results are nearly identical; the only difference between them is that we obtain certain “extra” 2-adic divisibilities in the presence of functions with parity, which have not yet been proven by the method which uses Shintani’s formulas.\par
\par
As we have already suggested, our method for studying L-values is based on a general theorem about \(p\)-adic Hilbert modular forms. It states, roughly speaking, that such a form which at one cusp has a p-adically integral \(q\)-expansion will have a p-adically integral \(q\)-expansion at every cusp. We apply this theorem to certain linear combinations of Eisenstein series and thus obtain integrality statements for corresponding linear combinations of L-values. These integrality statements include, for example, the “axioms” introduced by Coates in [3] and therefore lead in particular to the construction of \(p\)-adic L-functions over totally real fields.\par
\par
For the convenience of the reader, we now briefly state our main results concerning modular forms and integrality of L-values. We then explain how these relate to \(p\)-adic L-functions.\par
\par
q-Expansions of Hilbert Modular Forms\par
\par
Let \(K\) be a totally real field, which for simplicity we choose different from the rational field. Let \(r=[K:\mathbb{Q}]\) be its degree. Let\par
\par
\DFormula{\mathfrak{H}\,=\,\lbrace \tau \in K\otimes \mathbb{C}\,\vert \,\operatorname{Im}\,\tau \gg 0\rbrace }
\DAnchor{3}{229}
be the usual Hilbert upper half plane associated to \(K\). Let \(k\) be a non-negative integer. We define an action \(\vert _{k}\) of \(\operatorname{SL}(2,K\otimes \mathbb{R})\) on the space of \(\mathbb{C}\)-valued functions \(F(\tau )\) on \(\mathfrak{H}\) by the usual formula\par
\par
\DFormula{(F\vert _{k}\,\begin{pmatrix}a&b\\c&d\end{pmatrix})(\tau )\,=\,\mathcal{N}(c\tau +d)^{-k}\,F((a\tau +b)/(c\tau +d)),}
\par
in which \(\mathcal{N}\): \(K\otimes \mathbb{C}\to \mathbb{C}\) denotes the norm.\par
\par
We now let \(N\) denote a positive integer. Let \(\Gamma _{00}(N)\) denote the subgroup of \(\operatorname{SL}(2,K)\) consisting of those matrices \(\begin{pmatrix}a&b\\c&d\end{pmatrix}\) which satisfy the conditions:\par
\par
\DFormula{a,d\in 1+N\mathcal{O},}
\par
\DFormula{b\in 𝔇^{-1},\,\,\,\,\,c\in N𝔇,}
\par
where \(\mathcal{O}\) and \(𝔇\) denote respectively the integer ring and the different of \(K\). \(A\) (Hilbert) modular form of weight \(k\) on \(\Gamma _{00}(N)\) is a holomorphic function\par
\par
\DFormula{F:\,\mathfrak{H}\to \mathbb{C}}
\par
such that\par
\par
\DFormula{F\vert _{k}\,M=F}
\par
for each \(M\in \Gamma _{00}(N)\).\par
\par
The condition \(K\neq \mathbb{Q}\) implies that \(F\) is “holomorphic at infinity,” as is well known. Since \(F\) is invariant under translations \(\tau \mapsto \tau +b\) \((b\in 𝔇^{-1})\), we may expand \(F\) as a Fourier series\par
\par
\DFormula{c(0)\,+\,\sum _{\mu \in \mathcal{O},\,\mu \gg 0}\,c(\mu )q^{\mu },}
\par
where\par
\par
\DFormula{q^{\mu }\,=\,e^{2\pi i\cdot \operatorname{tr}(\mu \cdot \tau )}.}
\par
This series is the standard \(q\)-expansion of \(F\). To obtain the others, it is convenient to introduce the action on \(F\) of the group \(\operatorname{SL}(2,\widehat{K})\), where \(\widehat{K}\) is the ring of finite adeles of \(K\). For this, let \(\widehat{\Gamma }\) be the closure in \(\operatorname{SL}(2,\widehat{K})\) of the group \(\Gamma =\Gamma _{00}(N)\). Using the strong approximation theorem for \(\operatorname{SL}(2)\), we write\par
\par
\DFormula{\operatorname{SL}(2,\widehat{K})=\widehat{\Gamma }\cdot \operatorname{SL}(2,K).}
\par
Given \(M\in \operatorname{SL}(2,\widehat{K})\), we write \(M=M_{1}M_{2}\), with \(M_{1}\in \widehat{\Gamma }\) and \(M_{2}\in \operatorname{SL}(2,K)\). We then define \(F\vert _{k}\) \(M\) to be \(F\vert _{k}\) \(M_{2}\). For each invertible element \(\alpha \) of \(\widehat{K}\) (i.e., for each “finite idele”) we let\par
\par
\DFormula{F_{\alpha }\,=\,F\vert _{k}\,\begin{pmatrix}\alpha&0\\0&\alpha^{-1}\end{pmatrix}.}
\par
We may write \(F_{\alpha }\) as a Fourier series\par
\par
\DFormula{c(0,\alpha )\,+\,\sum \,c(\mu ,\alpha )q^{\mu },}
\DAnchor{4}{230}
where the sum is restricted to totally positive elements of \(K\) which lie in the square of the ideal \((\alpha )\) of \(K\) “generated” by \(\alpha \). We call this series the \(q\)-expansion of \(F\) at the cusp determined by \(\alpha \). We are interested in comparing the various \(q\)-expansions of a form, or of a collection of forms. For example, we recall the following result of Rapoport [28]:\par
\par
(0.1) Theorem. Suppose that the \(q\)-expansion coefficients of \(F\) at one cusp are rational numbers, i.e., that there is an \(\alpha \in \widehat{K}^{\ast}\) such that \(c(\mu ,\alpha )\in \mathbb{Q}\) for all \(\mu \). Then the \(q\)-expansion coefficients of \(F\) are rational for each \(\alpha \in \widehat{K}^{\ast}\).\par
\par
We now consider forms \(F_{k}\) \((k\geq 0)\) on \(\Gamma _{00}(N)\) of respective weights 0,1,2,..., all but finitely many of which are zero. We assume that each form satisfies the rationality condition of (0.1). For each \(k\) and each \(\alpha \in \widehat{K}^{\ast}\), let \(F_{k,\alpha }\) be the \(q\)-expansion of \(F_{k}\) at the cusp determined by \(\alpha \). We view these \(q\)-expansions as formal power series over \(\mathbb{Q}\), in the variables \(q^{\mu }\) (with \(\mu =0\) or \(\mu \) totally positive).\par
\par
Let \(p\) be a prime. For \(\alpha \in \widehat{K}\), the pth component \(\alpha _{p}\) of \(\alpha \) is the image of \(\alpha \) in \(K\otimes \mathbb{Q}_{p}\); its norm \(\mathcal{N}\alpha _{p}\) is an element of \(\mathbb{Q}_{p}\). When \(\alpha \in \widehat{K}^{\ast}\), the sum\par
\par
\DFormula{S(\alpha )\,=\,\sum _{k\geq 0}\,\mathcal{N}\alpha _{p}^{-k}\,F_{k,\alpha }}
\par
is thus a formal power series with coefficients in \(\mathbb{Q}_{p}\).\par
\par
(0.2) Theorem. If \(S(\alpha )\) has coefficients in \(\mathbb{Z}_{p}\) for one \(\alpha \), then \(S(\alpha )\) has coefficients in \(\mathbb{Z}_{p}\) for each \(\alpha \).\par
\par
(0.3) Corollary. Let \(\alpha \) be such that \(S(\alpha )\) has \(\mathbb{Z}_{p}\) coefficients, except possibly for its constant coefficient. Then for all \(\beta \in \widehat{K}^{\ast}\), \(S(\beta )\) again has non-constant coefficients in \(\mathbb{Z}_{p}\), and the difference between the constant coefficients of \(S(\alpha )\) and \(S(\beta )\) lies in \(\mathbb{Z}_{p}\).\par
\par
These assertions are special cases of (5.13), (5.14) as explained in (5.15).\par
\par
Kummer Congruences\par
\par
When the \(F_{k}\) are Eisenstein series, (0.3) yields congruences among L-values. We now state the principal such congruence (cf. (8.2)). Let \(K\) and \(p\) be as above, and let \(𝔣\) be a non-zero integral ideal of \(K\). Let \(G_{𝔣}\) be the strict ray class group of \(K\) mod \(𝔣\). Let \(G\) be the ray class group of conductor \(𝔣^{\infty }p^{\infty }\), corresponding via class field theory to the largest abelian extension of \(K\) which is unramified at all finite places of \(K\) prime to \(𝔣p\). When \(\varepsilon \) is a complex-valued function on \(G_{𝔣}\), we define as usual\par
\par
\DFormula{L(s,\varepsilon )\,:=\,\sum \,\varepsilon (𝔵)\mathcal{N}𝔵^{-s}\,\,\,\,\,(\operatorname{Re}(s)>1),}
\par
taking the sum over \(prime-to-𝔣\) integral ideals of \(K\). One knows that \(L(s,\varepsilon )\) may be continued to a meromorphic function on \(\mathbb{C}\), holomorphic except for a possible simple pole at \(s=1\). In particular, for each integer \(k\geq 1\), values \(L(1-k,\varepsilon )\) are defined.\par
\par
According to a fundamental theorem of Siegel [38], the association\par
\par
\DFormula{\varepsilon \mapsto L(1-k,\varepsilon )}
\DAnchor{5}{231}
is rational in the sense that \(L(1-k,\varepsilon )\) is a rational number whenever \(\varepsilon \) is \(\mathbb{Q}\)-valued. This being the case, we define by linearity values \(L(1-k,\varepsilon )\in V\) for each function \(\varepsilon \) on \(G_{𝔣}\) with values in a \(\mathbb{Q}\)-vector space \(V\). This construction applies especially when \(V\) is a field of characteristic zero, for example \(\mathbb{Q}_{p}\).\par
\par
We now let \(\mathcal{N}:G\to \mathbb{Z}_{p}^{\ast}\) be that continuous character whose value on the image in \(G\) of a \(prime-to-𝔣p\) ideal is its norm. For \(c\in G\), \(k\geq 1\), and \(\varepsilon \) a \(\mathbb{Q}_{p}\)-valued function on \(G_{𝔣}\), we define\par
\par
\DFormula{\Delta _{c}(1-k,\varepsilon )=L(1-k,\varepsilon )-\mathcal{N}c^{k}\,L(1-k,\varepsilon _{c})\,\in \,\mathbb{Q}_{p},}
\par
where \(\varepsilon _{c}\) is the function \(g\mapsto \varepsilon (cg)\) and the product cg is computed in \(G_{𝔣}\). Let \(\varepsilon _{1},\varepsilon _{2}\),... be a sequence of such functions, only finitely many of which are non-zero. For \(𝔵\) a \(prime-to-𝔣\) ideal, set\par
\par
\DFormula{\varphi (𝔵)=\sum _{k\geq 1}\,\varepsilon _{k}(𝔵)\mathcal{N}𝔵^{k-1}\in \mathbb{Q}_{p}.}
\par
(0.4) Theorem. Suppose that \(\varphi (𝔵)\in \mathbb{Z}_{p}\) for each \(prime-to-𝔣\) ideal \(𝔵\). Then for all \(c\in G\), we have \(\Delta \in \mathbb{Z}_{p}\), where\par
\par
\DFormula{\Delta =\sum _{k\geq 1}\,\Delta _{c}(1-k,\varepsilon _{k}).}
\par
Our main theorem (8.2) is equivalent to (0.4). We prove also some further assertions concerning the case where \(p=2\) and the \(\varepsilon _{k}\) satisfy certain parity conditions. Briefly, if \(p=2\) and if \(\varepsilon _{k}\) is an odd (resp. even) function when \(k\) is odd (resp. even), then \(\Delta \in 2^{r-1}\mathbb{Z}_{2}\). We may furthermore give necessary and sufficient conditions for the divisibility of \(\Delta \) by \(2^{r}\). (See (8.11), (8.12).)\par
\par
\(p\)-adic L-functions\par
\par
It is now well known that the integrality theorem (0.4) may be used to construct \(p\)-adic L-functions [3, 30, 34]. In the above situation, we take \(𝔣\) to be divisible by each prime of \(K\) lying over \(p\). Whenever \(𝔣'\) is an ideal divisible by \(𝔣\) and dividing \((𝔣p)^{n}\) for some \(n\geq 1\), the ideals of \(K\) prime to \(𝔣'\) are just those prime to \(𝔣\). Hence for \(\varepsilon \) a locally constant function on \(G\) with values in a \(\mathbb{Q}\)-vector space \(V\), we may unambiguously define values \(L(1-k,\varepsilon )\in V\) by selecting an \(𝔣'\) such that \(\varepsilon \) factors through \(G\to G_{𝔣'}\) and applying the above procedure, with \(𝔣\) replaced by \(𝔣'\). We may thus view the functions\par
\par
\DFormula{\varepsilon \mapsto L(1-k,\varepsilon )\,\,\,\,\,(k\geq 1)}
\par
as \(\mathbb{Q}\)-valued distributions on \(G\). For \(k\geq 1\) and \(c\in G\), the map\par
\par
\DFormula{\mu _{c,k}:\,\varepsilon \mapsto \Delta _{c}(1-k,\varepsilon )}
\par
is then a \(\mathbb{Q}_{p}\)-valued distribution on \(G\).\par
\par
From (0.4), we immediately obtain (cf. [30, 4.1])\par
\par
(0.5) Theorem. The distribution \(\mu _{c,k}\) is in fact a measure with values in \(\mathbb{Z}_{p}\). We have the formula\textsuperscript{1}\par
\par
\DFormula{\mu _{c,k}=\mathcal{N}^{k-1}\,\mu _{c,1}.}
\par
\(^{1}\) In this formula, the right-hand member represents the product of the function \(\mathcal{N}^{k-1}\) and the measure \(\mu _{c,1}\). Such products are to be distinguished from the (convolution) product of two measures.\par
\DAnchor{6}{232}
Following a normalization of Serre, we put \(\lambda _{c}=\mathcal{N}^{-1}\mu _{c,1}\) for \(c\in G\). For \(k\geq 1\) and for \(\varepsilon \) locally constant, we have\par
\par
\DFormula{\int \,\varepsilon \mathcal{N}^{k}\,d\lambda _{c}\,=\,\Delta _{c}(1-k,\varepsilon ).}
\par
As in [34], we regard the various measures \(\lambda _{c}\) as elements of the ring \(\Lambda =\Lambda _{G}\) of \(\mathbb{Z}_{p}\)-valued measures on \(G\); this ring may alternately be described as a sort of completed group ring \(\mathbb{Z}_{p}[[G]]\). Choose c so that 1-c is not a zero-divisor in \(\Lambda \) and let\par
\par
\DFormula{\lambda \,=\,(1/(1-c))\,\lambda _{c}}
\par
in the total quotient ring of \(\Lambda \). One sees easily that \((1-c')\lambda =\lambda _{c'}\) for \(c'\in G\), so that \(\lambda \) is a pseudo-measure in the sense of [34]. With \(k\) and \(\varepsilon \) as above, the integral against \(\lambda \) of \(\varepsilon \mathcal{N}^{k}\) is \(L(1-k,\varepsilon )\). We thus obtain a result cited by Serre in his discussion of L-values [34, 3.5]. As Serre notes, one now obtains the usual \(p\)-adic L-functions of \(K\) (attached to characters whose conductors divide powers of \(𝔣)\) by integrating suitable characters against \(\lambda \). (See also [30, §4].)\par
\par
The plan of this paper is as follows. §1 begins by recalling the Kummer congruences in the case \(K=\mathbb{Q}\). Although we do not prove these congruences here, they may easily be derived from explicit expressions for values \(L(1-k,\varepsilon )\) in terms of Bernoulli polynomials. Our aim is to motivate what follows by reinterpreting the congruences in terms of measures on \(\widehat{\mathbb{Z}}\) and on \(\widehat{\mathbb{Q}}\). In §2 and §3, we study the analogues \(A\) and \(I\) of these two spaces for the case of an arbitrary totally real field \(K\neq \mathbb{Q}\). Especially, we prove a functional equation for L-functions which involves a somewhat curious Fourier transform for functions on \(I\).\par
\par
In the next two §§, we prove an irreducibility theorem for the Hilbert-Blumenthal moduli problem and give applications to \(p\)-adic modular forms. We obtain (0.2), in particular. Then in §6 and §7 we construct Eisenstein series and theta series for the Hilbert group and compute their \(q\)-expansions. From (0.3) we can then obtain information about L-values. In the final § we make explicit this information and use it to prove (0.4) and its 2-adic refinements. For this, we work with measures on the space \(I\).\par
\par
As suggested above, most of the work presented in this paper was in principle completed in 1976. However, the material concerning theta series grew out of work done in 1977--78. The second author wishes to thank the Sloan Foundation and the N.S.F. for financial support and the I.H.E.S. for its continuing hospitality.\par
\par
\subsection*{Contents}
\par
1. Review of Kummer Congruences over \(\mathbb{Q}\)\dotfill 233\par
2. The Measure Spaces \(A\) and \(I\)\dotfill 238\par
3. \(A\) functional Equation\dotfill 245\par
4. An Irreducibility Theorem\dotfill 254\par
5. Hilbert Modular Forms\dotfill 258\par
6. Eisenstein Series\dotfill 267\par
7. Theta Series\dotfill 273\par
8. Congruences for L-values\dotfill 277\par
\DAnchor{7}{233}
\subsection*{§1. Review of Kummer Congruences over \(\mathbb{Q}\)}
\par
Let \(\varepsilon \) be a complex valued function defined on the set \(\mathbb{Z}/f\mathbb{Z}\) of integers mod f. Viewing \(\varepsilon \) as a function on \(\mathbb{Z}\), periodic mod f, we set\par
\par
\DFormula{(1.1)\,\,L(s,\varepsilon )=\sum _{n\geq 1}\,\varepsilon (n)n^{-s}}
\par
for \(s\in \mathbb{C}\) of real part Re \(s>1\). As is well known, this L-function may be continued to a meromorphic function of \(s\), analytic except for a possible simple pole at \(s=1\). We are interested in its values at the non-negative integral points \(s=1,...,k\) with \(k\geq 1\).\par
\par
Define Bernoulli polynomials \(B_{k}(x)\in \mathbb{Q}[x]\), \(k\geq 0\), by the formal expansion\par
\par
\DFormula{\frac{Ze^{xZ}}{e^Z-1}=\sum_{k\geq0}B_k(x)\frac{Z^k}{k!}.}
\par
(1.2) Theorem. For \(k\geq 1\), we have\par
\par
\DFormula{(1.3)\,\,L(1-k,\varepsilon )=-(f^{k-1}/k)\sum _{t=1}^{f}\,\varepsilon (t)B_{k}(t/f).}
\par
Example. If \(k=1\), the theorem gives\par
\par
\DFormula{L(0,\varepsilon )=-\sum _{t=1}^{f}\,\varepsilon (t)(t/f-1/2).}
\par
In the special case where \(\sum _{t=1}^{f}\) \(\varepsilon (t)=0\), the value \(L(0,\varepsilon )\) is thus the Cesaro sum of the series \(\sum _{t=1}^{\infty }\) \(\varepsilon (t)\): we have\par
\par
\DFormula{\lim _{N\to \infty }\,(1/N)\sum _{n=1}^{N}\sum _{t=1}^{n}\,\varepsilon (t)}
\DFormula{=\,(1/f)\sum _{n=1}^{f}\sum _{t=1}^{n}\,\varepsilon (t)}
\DFormula{=\,(1/f)\sum _{t=1}^{f}\,\varepsilon (t)(f-t+1)}
\DFormula{=\,-(1/f)\sum _{t=1}^{f}\,t\varepsilon (t).}
\par
The general formula for \(L(0,\varepsilon )\) may be recovered from this special one after one knows that \(\zeta (0)=-1/2\), where \(\zeta \) is the Riemann zeta function. This latter fact is just Euler’s formula \(\sum _{n\in \mathbb{Z}}1=0\), as we see by rewriting the sum as\par
\par
\DFormula{1+2\cdot \sum _{n\geq 1}1=1+2\zeta (0).}
\par
Discussions of the analytic continuation of \(L(s,\varepsilon )\) and its evaluation at negative integers may be found for example in [11; \(I\), pp. 72--78], [13, 23, 42].\par
\par
Since the numbers \(B_{k}(t/f)\) are rational, (1.3) may be used as the definition of values \(L(1-k,\varepsilon )\) for \(\varepsilon \) now a function mod f with values in an arbitrary \(\mathbb{Q}\)-vector space \(V\):\par
\par
\DFormula{(1.4)\,\,L(1-k,\varepsilon )=\sum _{t=1}^{f}\,[-f^{k-1}B_{k}(t/f)/k]\cdot \varepsilon (t).}
\DAnchor{8}{234}
Furthermore, we can unambiguously define such values for \(\varepsilon \) a periodic function \(\mathbb{Z}\to V\) (i.e. a function defined mod f for some f) whose modulus of periodicity f is not specified alone with \(\varepsilon \). [To verify this it suffices to examine the case \(V=\mathbb{Q}\) and thus a fortiori it is sufficient to treat the case \(V=\mathbb{C}\); then one simply observes that the number f does not intervene in (1.1).] In other terms, (1.4) defines for a given \(k\geq 1\) a \(\mathbb{Q}\)-linear map\par
\par
(periodic functions \(\mathbb{Z}\to \mathbb{Q})\) \(\to \) \(\mathbb{Q}\),\par
\par
and by tensoring with a \(\mathbb{Q}\)-vector space \(V\) we obtain\par
\par
(periodic function \(\mathbb{Z}\to V)\) \(\to \) \(V\).\par
\par
\DFormula{\varepsilon \mapsto L(1-k,\varepsilon ).}
\par
Since restriction from \(\widehat{\mathbb{Z}}\) to \(\mathbb{Z}\) identifies the locally constant functions on \(\widehat{\mathbb{Z}}\) with the periodic functions on \(\mathbb{Z}\), the map\par
\par
\DFormula{\varepsilon \mapsto L(1-k,\varepsilon )}
\par
may be viewed as a distribution \(T_{k}\) on \(\widehat{\mathbb{Z}}\) with values in \(\mathbb{Q}\). Such a distribution may alternately be regarded as a function on the compact open subsets \(U\) of \(\widehat{\mathbb{Z}}\), here explicitly given by the formula\par
\par
\DFormula{(1.5)\,\,T_{k}(U)=\sum _{n\in U}\,n^{-s}\vert _{s=1-k}.}
\par
(The summation is extended over the positive integers \(n\in U.)\) From (1.5) we may deduce immediately the “invariance”\par
\par
(1.6)  \(T_{k}(nU)=n^{k-1}T_{k}(U)\)\par
\par
for positive integers n. This invariance implies that \(T_{k}\) has a unique extension to a distribution \(\widetilde{T}_{k}\) on \(\widehat{\mathbb{Q}}=\widehat{\mathbb{Z}}\otimes \mathbb{Q}\) which satisfies the invariance\par
\par
(1.7)  \(\widetilde{T}_{k}(nU)=n^{k-1}\widetilde{T}_{k}(U)\)\par
\par
for n a positive rational number and \(U\) a compact-open subset of \(\widehat{\mathbb{Q}}\). This distribution is again defined by (1.5), with the summation now over all positive \(n\in \mathbb{Q}\).\par
\par
Another way to view distributions on \(\widehat{\mathbb{Z}}\) is to regard them as collections of functions \(\alpha _{f}\) on the groups \(\mathbb{Z}/f\mathbb{Z}\) \((f\geq 1)\), linked by a suitable compatibility relation (cf. [26]). The value of \(\alpha _{f}\) on the class of a modulo f is the value of the distribution on the open subset \(a+f\widehat{\mathbb{Z}}\) of \(\widehat{\mathbb{Z}}\). For \(T_{k}\), if we take a in the range \(1\leq a\leq f\), then (1.4) states that \(\alpha _{f}\) maps (a mod f) to the rational number\par
\par
\DFormula{-(f^{k-1}/k)B_{k}(a/f).}
\par
This formula, or else the invariance (1.6), implies that the maps \(\alpha _{f}\) associated to \(T_{1}\) are each obtained by composing the natural maps \(\mathbb{Z}/f\mathbb{Z}\to \mathbb{Q}/\mathbb{Z}\) with a single map \(\alpha :\mathbb{Q}/\mathbb{Z}\to \mathbb{Q}\). In the language of Kubert-Lang ([21]), \(T_{1}\) is thus an ordinary distribution.\par
\DAnchor{9}{235}
We now introduce “twists.” If c is a positive integer and \(\varepsilon \) a periodic function on \(\mathbb{Z}\) with values in a \(\mathbb{Q}\)-vector space, we write \(\varepsilon _{c}\) for the function \(x\mapsto \varepsilon (cx)\). If we are regarding \(\varepsilon \) as a function modulo f, then c will typically be taken prime to f. Set\par
\par
(1.8)  \(\Delta _{c}(1-k,\varepsilon )=L(1-k,\varepsilon )-c^{k}L(1-k,\varepsilon _{c})\).\par
\par
Now let \(\varepsilon _{1},...,\varepsilon _{t}\) be functions mod f with values in a \(p\)-adic field \(\mathbb{Q}_{p}\), and let \(k_{1},...,k_{t}\) be positive integers. The following result summarizes the “generalized Kummer congruences” as considered by Mazur and others (cf. [20, 24, 25]).\par
\par
(1.9) Theorem. Suppose for all \(n\geq 1\) that we have\par
\par
\DFormula{\sum _{i=1}^{t}\,\varepsilon _{i}(n)n^{k_{i}-1}\in \mathbb{Z}_{p}.}
\par
Then for each \(c\geq 1\), prime to pf, we have\par
\par
\DFormula{\sum _{i=1}^{t}\,\Delta _{c}(1-k_{i},\varepsilon _{i})\in \mathbb{Z}_{p}.}
\par
This elementary theorem may be proved directly from (1.4). It is the analogue for \(K=\mathbb{Q}\) of (0.4). As mentioned in the Introduction, we now wish to rephrase (1.9) in terms of measures.\par
\par
First, taking a more “distribution-theoretic” point of view, we eliminate the modulus f from the statement of the theorem. For this, it is desirable to allow twisting by c which are simultaneously prime to all moduli f: these will be elements of the group \(\widehat{\mathbb{Z}}^{\ast}\subset \widehat{\mathbb{Z}}\).\par
\par
So let \(\varepsilon \) be a locally constant function on \(\widehat{\mathbb{Z}}\) with values in a \(\widehat{\mathbb{Q}}\)-module \(V\) \((\widehat{\mathbb{Q}}=\mathbb{Q}\otimes \widehat{\mathbb{Z}}\) is the ring of finite adeles for \(\mathbb{Q})\). Then for \(c\in \widehat{\mathbb{Z}}^{\ast}\) we again define \(\varepsilon _{c}\) and elements \(\Delta _{c}(1-k,\varepsilon )\in V\). We regard \(V=\mathbb{Q}_{p}\) as a \(\widehat{\mathbb{Q}}\)-module via the projection \(\widehat{\mathbb{Q}}\to \mathbb{Q}_{p}\).\par
\par
The hypothesis of (1.9) may be rephrased as follows: the locally constant functions \(\varepsilon _{i}:\widehat{\mathbb{Z}}\to \mathbb{Q}_{p}\) and positive integers \(k_{i}\) \((i=1,...,t)\) are such that for all \(n\in \widehat{\mathbb{Z}}\) we have \(\varphi (n)\in \mathbb{Z}_{p}\) where\par
\par
\DFormula{\varphi (n)=\sum _{i=1}^{t}\,\varepsilon _{i}(n)n^{k_{i}-1}.}
\par
Indeed, \(\varphi \) is a continuous function \(\widehat{\mathbb{Z}}\to \mathbb{Q}_{p}\) and is thus \(\mathbb{Z}_{p}\)-valued if and only if its values on positive integers lie in \(\mathbb{Z}_{p}\). The conclusion may be rephrased as follows: for all \(c\in \widehat{\mathbb{Z}}^{\ast}\), we have\par
\par
\DFormula{\sum _{i=1}^{t}\,\Delta _{c}(1-k_{i},\varepsilon _{i})\in \mathbb{Z}_{p}.}
\par
In fact, one sees immediately that this statement for a given \(c\in \widehat{\mathbb{Z}}^{\ast}\) and the conclusion of (1.9) for a given \(d\geq 1\) are equivalent if \(c\equiv d\) mod \(fp^{n}\), where f is a modulus for all the \(\varepsilon _{i}\) and where \(p^{n}\) is a common denominator for the (finitely many) numbers\par
\par
\DFormula{L(1-k_{i},\varepsilon _{i,c})\,\,\,\,\,(i=1,...,t;\,c\,\operatorname{mod}\,f).}
\DAnchor{10}{236}
Having eliminated the f, we may also eliminate the \(p\). Namely, the following is trivially equivalent to (1.9).\par
\par
(1.9a) Theorem. Let \(\varepsilon _{1},...,\varepsilon _{t}\) be locally constant functions on \(\widehat{\mathbb{Z}}\) with values in \(\widehat{\mathbb{Q}}\). Let \(k_{1},...,k_{t}\) be positive integers. Suppose that\par
\par
\DFormula{\sum _{i=1}^{t}\,\varepsilon _{i}(x)x^{k_{i}-1}\in \widehat{\mathbb{Z}}}
\par
for all \(x\in \widehat{\mathbb{Z}}\). Then for each \(c\in \widehat{\mathbb{Z}}^{\ast}\) we have\par
\par
\DFormula{\sum _{i=1}^{t}\,\Delta _{c}(1-k_{i},\varepsilon _{i})\in \widehat{\mathbb{Z}}.}
\par
We now rephrase this result in terms of measures on \(\widehat{\mathbb{Z}}\). Let \(c\in \widehat{\mathbb{Z}}^{\ast}\). For each locally constant function \(\varepsilon :\widehat{\mathbb{Z}}\to \widehat{\mathbb{Z}}\), the quantity \(\Delta _{c}(0,\varepsilon )\) lies in \(\widehat{\mathbb{Z}}\) by (1.9a). The map \(\varepsilon \mapsto \Delta _{c}(0,\varepsilon )\) is thus a measure \(\mu _{c}\) on \(\widehat{\mathbb{Z}}\) with values in \(\widehat{\mathbb{Z}}\), so that a quantity \(\int \varphi \) \(d\mu _{c}\in \widehat{\mathbb{Z}}\) is defined for all continuous functions \(\varphi :\widehat{\mathbb{Z}}\to \widehat{\mathbb{Z}}\).\par
\par
In particular, let \(x\) be the identity function \((a\mapsto a)\) on \(\widehat{\mathbb{Z}}\), and let \(\varepsilon \) be locally constant. Then (1.9a) implies that we have\par
\par
\DFormula{\int \varepsilon x^{k-1}\,d\mu _{c}=\Delta _{c}(1-k,\varepsilon )}
\par
for all \(k\geq 1\). Indeed, to verify this it suffices to check that the difference between the two sides is congruent to 0 mod \(N\), for each \(N\geq 1\). Let \(N\) be given, and let \(\eta :\widehat{\mathbb{Z}}\to \widehat{\mathbb{Z}}\) be a locally constant function which is congruent mod \(N\) to \(x^{k-1}\). Because \(\mu _{c}\) is a measure we have\par
\par
\DFormula{\int \varepsilon x^{k-1}\,d\mu _{c}\equiv \int \varepsilon \eta \,d\mu _{c}\,\operatorname{mod}\,N\cdot \widehat{\mathbb{Z}}.}
\par
\subsection*{Thus it suffices to check that}
\par
\DFormula{\Delta _{c}(0,\varepsilon \eta )=\int \varepsilon \eta \,d\mu _{c}\equiv \Delta _{c}(1-k,\varepsilon )\,(\operatorname{mod}\,N).}
\par
This follows from (1.9a) and the congruence\par
\par
\DFormula{\varepsilon (x)\eta (x)\equiv \varepsilon (x)x^{k-1}\,\,\,\,\,\text{for}\,x\in \widehat{\mathbb{Z}}.}
\par
Summarizing, we have obtained from the Kummer congruences the following result.\par
\par
(1.9b) Theorem. For each \(c\in \widehat{\mathbb{Z}}^{\ast}\), the map \(\varepsilon \mapsto \Delta _{c}(0,\varepsilon )\) is a measure \(\mu _{c}\) on \(\widehat{\mathbb{Z}}\) with values in \(\widehat{\mathbb{Z}}\). For all \(k\geq 1\), the product \(x^{k-1}\cdot \mu _{c}\) is the distribution \(\varepsilon \mapsto \Delta _{c}(1-k,\varepsilon )\) (which is consequently a measure as well).\par
\par
Conversely, this theorem immediately implies (1.9a). For suppose that we are given \(\varepsilon _{i}\), \(k_{i}\) and c as in the hypothesis of (1.9a). Then we have, assuming (1.9b),\par
\par
\DFormula{\sum _{i=1}^{t}\,\Delta _{c}(1-k_{i},\varepsilon _{i})=\int _{\widehat{\mathbb{Z}}}(\sum _{i=1}^{t}\,\varepsilon _{i}\,x^{k_{i}-1})d\mu _{c}\in \widehat{\mathbb{Z}}.}
\DAnchor{11}{237}
In yet another variant, we consider locally constant functions \(\varepsilon \) with compact support on \(\widehat{\mathbb{Q}}\). For these, we define values \(L(1-k,\varepsilon )\) as the integrals of \(\varepsilon \) against the distributions \(\widetilde{T}_{k}\) introduced above. Formally,\par
\par
\DFormula{(1.10)\,\,L(1-k,\varepsilon )=\sum _{n>0,\,n\in \mathbb{Q}}\,\varepsilon (n)n^{k-1}.}
\par
Introducing twisted values \(\Delta _{c}(1-k,\varepsilon )\) as in (1.8), we obtain from the invariance (1.7) the following integrality properties:\par
\par
(1.11) For \(c\in \widehat{\mathbb{Z}}^{\ast}\), the map \(\mu _{c}\): \(\varepsilon \mapsto \Delta _{c}(0,\varepsilon )\) is a \(\widehat{\mathbb{Z}}\)-valued measure on \(\widehat{\mathbb{Q}}\).\par
\par
(1.12) For \(k\geq 1\), the distribution \(x^{k-1}d\mu _{c}\) on \(\widehat{\mathbb{Q}}\) is the map\par
\par
\DFormula{\varepsilon \mapsto \Delta _{c}(1-k,\varepsilon ).}
\par
Note that the map \(x^{k-1}\) is not bounded when \(k>1\). Hence \(x^{k-1}d\mu _{c}\) will not be a measure on \(\widehat{\mathbb{Q}}\) for \(k>1\).\par
\par
Finally, we discuss an “extra divisibility” at 2 which arises formally from parity considerations. For simplicity, we discuss only functions \(\varepsilon \) supported on \(\widehat{\mathbb{Z}}\).\par
\par
Let \(\delta \) be the Dirac measure at 0 on \(\widehat{\mathbb{Z}}\):\par
\par
\DFormula{\delta :\,\varepsilon \mapsto \varepsilon (0).}
\par
(1.13) Lemma. \(\mu _{(-1)}=-\delta \).\par
\par
Proof. We must verify that \(L(0,\varepsilon )+L(0,\varepsilon _{-1})=-\varepsilon (0)\).\par
\par
\subsection*{This follows immediately from the formula}
\par
\DFormula{L(0,\varepsilon )=-\sum _{t=1}^{f}\,\varepsilon (t)(t/f-1/2),}
\par
which is a special case of (1.3).\par
\par
(1.14) Corollary. For \(k>1\), we have \(x^{k-1}\cdot \mu _{(-1)}=0\), i.e.,\par
\par
\DFormula{L(1-k,\varepsilon )=(-1)^{k}\,L(1-k,\varepsilon _{-1}).}
\par
Proof. We have \(x^{k-1}\cdot \delta =0\).\par
\par
Remark. (1.13) and (1.14) summarize the presence of trivial zeros of L-functions for complex valued characters \(\varepsilon \).\par
\par
Definition. \(A\) function \(\varphi \) on \(\widehat{\mathbb{Z}}\) with values in a \(\mathbb{Q}\)-vector space \(V\) is odd (resp. even) if it satisfies \(\varphi (-x)=-\varphi (x)\) (resp. \(\varphi (-x)=\varphi (x))\) for \(x\in \widehat{\mathbb{Z}}\).\par
\par
From (1.14) we see that if \(k\) is even then \(L(1-k,\varepsilon )=0\) if a locally constant \(\varepsilon \) is odd, whereas \(L(1-k,\varepsilon )=0\) for even functions if \(k>1\) is odd. For \(k=1\), we have \(L(1-k,\varepsilon )=0\) for even \(\varepsilon \) which vanish at 0 (1.13).\par
\par
(1.15) Theorem. Let \(\varphi :\widehat{\mathbb{Z}}\to \widehat{\mathbb{Z}}\) be an odd function. Then we have\par
\par
\DFormula{\int \varphi \,d\mu _{c}\in 2\widehat{\mathbb{Z}}}
\par
for all \(c\in \widehat{\mathbb{Z}}^{\ast}\).\par
\DAnchor{12}{238}
Proof. Let \(\alpha \) be a locally constant function which is congruent to \(\varphi \) mod 4 and which vanishes at 0. Set\par
\par
\DFormula{\alpha ^{+}(x)=[\alpha (x)+\alpha (-x)]/2,\,\,\,\,\,\alpha ^{-}(x)=[\alpha (x)-\alpha (-x)]/2.}
\par
Since \(\alpha ^{+}\) is even and vanishes at 0, we have \(\int \alpha ^{+}\) \(d\mu _{c}=0\), from the above discussion. Moreover, we have \(\varphi \equiv \alpha ^{-}\) mod 2, since\par
\par
\DFormula{2\varphi (x)=\varphi (x)-\varphi (-x)\equiv \alpha (x)-\alpha (-x)=2\alpha ^{-}(x)\,(\operatorname{mod}\,4).}
\par
\subsection*{Hence we have}
\par
\DFormula{\int \varphi \,d\mu _{c}\equiv \int \alpha ^{-}\,d\mu _{c}\,(\operatorname{mod}\,2).}
\par
These computations show that we may assume that \(\varphi \) is in fact locally constant, say mod f. We make this assumption.\par
\par
Since \(\varphi \) is odd, we may find a function \(\varepsilon :\mathbb{Z}/f\mathbb{Z}\to \widehat{\mathbb{Z}}\), vanishing at 0, such that \(\varphi (x)=\varepsilon (x)-\varepsilon (-x)\). We find\par
\par
\DFormula{\int \varphi (x)d\mu _{c}=\Delta _{c}(0,\varepsilon )-\Delta _{c}(0,\varepsilon _{(-1)})=2\Delta _{c}(0,\varepsilon )\in 2\widehat{\mathbb{Z}}.}
\par
(1.16) Corollary. Let \(\varepsilon _{1},...,\varepsilon _{t}\) and \(k_{1},...,k_{t}\) be as in (1.9a). Suppose further for each \(i=1,...,t\) that \(\varepsilon _{i}\) and \(k_{i}\) have the same parity (namely: \(\varepsilon _{i}\) is even if \(k_{i}\) is even and odd if \(k_{i}\) is odd). Then for all \(c\in \widehat{\mathbb{Z}}^{\ast}\) we have\par
\par
\DFormula{\sum _{i=1}^{t}\,\Delta _{c}(1-k_{i},\varepsilon _{i})\in 2\widehat{\mathbb{Z}}.}
\par
\subsection*{§2. The Measure Spaces \(A\) and \(I\)}
\par
Let \(K\) be a totally real field, which will be fixed for the remainder of this paper. In order to avoid certain technical difficulties later on, we will assume right away that \(K\) is different from the rational field \(\mathbb{Q}\). We introduce the following notation and conventions, which will be in force for the remainder of this paper:\par
\par
\(r\) is the degree of \(K\) over \(\mathbb{Q}\).\par
\par
\(\mathcal{N}\) is the norm map \(K\to \mathbb{Q}\) or any norm map derived from it, such as the norm map on ideals of \(K\) or a map of the type \(K\otimes R\to R\), where \(R\) is a \(\mathbb{Q}\)-algebra.\par
\par
An ideal of \(K\) is tacitly understood to be non-zero and fractional (as opposed to integral), unless otherwise described. \(A\) conductor is a (non-zero) integral ideal. The group of ideals is denoted by \(I_{0}\) and the sub-monoid of integral ideals by \(A_{0}\).\par
\par
We denote the ring of integers of \(K\) by \(\mathcal{O}\) and its different by \(𝔇\). The unit ideal (1) is sometimes abbreviated 1.\par
\par
For \(\alpha \in K\), the symbol \(\alpha \gg 0\) means that \(\alpha \) is totally positive. For \(\alpha \in K\otimes \mathbb{R}\), we write \(\alpha \gg 0\) if the image of \(\alpha \) under each \(\mathbb{R}\)-algebra homomorphism\par
\par
\DFormula{K\otimes \mathbb{R}\to \mathbb{R}}
\par
is positive. Thus, if we number the embeddings \(K\to \mathbb{R}\), so that \(K\otimes \mathbb{R}\) may be written \(\mathbb{R}^{r}\), then \(\alpha \gg 0\) means that \(\alpha \) is an r-tuple of positive numbers.\par
\DAnchor{13}{239}
We let \(𝔸\) denote the ring of adeles of \(K\) and write either \(\widehat{K}\) or \(𝔸_{f}\) for the ring of finite adeles of \(K\). If \(x\in \widehat{K}^{\ast}\) is a “finite idele,” we write (x) for the ideal “generated” by \(x\), namely\par
\par
\(\prod _{𝔭}\) \(𝔭^{\mathrm{ord}_{𝔭}(x)}\),\par
\par
the product being extended over the prime ideals \(𝔭\) of \(K\). Also, when \(𝔞\) is an ideal, we write \(\widehat{𝔞}\) for the completion of \(𝔞\), i.e., the closure of \(𝔞\) in \(\widehat{K}\).\par
\par
(2.1) Definition. Let \(𝔞\) and \(𝔟\) be ideals, and let \(𝔣\) be a conductor. The ideals \(𝔞\) and \(𝔟\) are equivalent (or congruent) mod \(𝔣\) if the ideal \(𝔞𝔟^{-1}\) may be written as the principal ideal \((\alpha )\) for some totally positive \(\alpha \) in the set \(1+𝔣𝔟^{-1}\).\par
\par
We write \(\sim _{𝔣}\) for the equivalence relation introduced above. If \(𝔞\) and \(𝔟\) are equivalent mod \(𝔣\), we write \(𝔞\sim _{𝔣}𝔟\) and sometimes say that \(𝔞\) and \(𝔟\) are equivalent mod \(\sim _{𝔣}\). The following results follow easily from the definition of \(\sim _{𝔣}\).\par
\par
(2.2) Suppose that \(𝔞\sim _{𝔣}𝔟\) and that \(𝔡\) is an ideal containing \(𝔣\). Then \(𝔡\) contains \(𝔞\) if and only if \(𝔡\) contains \(𝔟\).\par
\par
(2.3) Corollary. Any ideal equivalent mod \(𝔣\) to an integral ideal is again integral.\par
\par
(2.4) If \(𝔞\) and \(𝔟\) are integral ideals which are equivalent mod \(𝔣\), then the greatest common divisors \((𝔞,𝔣)\) and \((𝔟,𝔣)\) are equal.\par
\par
(2.5) Let \(𝔞\) and \(𝔟\) be ideals and let \(𝔡\) be an integral ideal. We have \(𝔞\sim _{𝔣}𝔟\) if and only if \(𝔞𝔡\sim _{𝔣𝔡}𝔟𝔡\).\par
\par
(2.6) The set \(A_{𝔣}\) of mod \(𝔣\) classes of integral ideals is a monoid under a multiplication \(A_{𝔣}\times A_{𝔣}\to A_{𝔣}\) induced by multiplication of ideals. The invertible elements are those classes represented by ideals which are relatively prime to \(𝔣\). Further, the group \(G_{𝔣}\) of invertible elements of \(A_{𝔣}\) is precisely the group of strict ray classes of \(K\) modulo \(𝔣\): when \(𝔞\) and \(𝔟\) are \(prime-to-𝔣\) integral ideals, we have \(𝔞\sim _{𝔣}𝔟\) if and only if \(𝔞\) and \(𝔟\) represent the same strict ray class mod \(𝔣\).\par
\par
Using the above results, and especially (2.5), (2.6), we obtain the following description of \(A_{𝔣}\). It is a disjoint union of copies of the groups \(G_{𝔟}\), \(𝔟\) running over the divisors of \(𝔣\). The image of \(G_{𝔟}\) in \(A_{𝔣}\) consists of those classes whose greatest common divisor with \(𝔣\) is \(𝔣𝔟^{-1}\). This decomposition generalizes the fact that, for \(f\geq 1\), \(\mathbb{Z}/f\mathbb{Z}\) is a disjoint union of copies of \((\mathbb{Z}/d\mathbb{Z})^{\ast}\), with d running over the divisors of f.\par
\par
(2.7) Suppose that \(𝔞\) and \(𝔟\) are integral ideals divisible by \(𝔣\). Then \(𝔞\) is equivalent to \(𝔟\) mod \(𝔣\) if and only if \(𝔞\) is equivalent to \(𝔟\) mod (1).\par
\par
We have already introduced the notation \(A_{𝔣}\) for the (finite) set of classes of integral ideals mod \(𝔣\). Now let \(I_{𝔣}\) be the set of equivalence classes of all ideals mod \(𝔣\). For each ideal \(𝔟\) containing \(𝔣\), let \(I_{𝔣}^{𝔟}\) be the set of classes mod \(𝔣\) of ideals contained in \(𝔟\), cf. (2.2).\par
\par
(2.8) Each set \(I_{𝔣}^{𝔟}\) is finite. Indeed, if \(𝔟\) is the inverse of an integral ideal \(𝔞\), the map on ideals \(𝔵\mapsto 𝔞𝔵\) induces a bijection\par
\par
\DFormula{I_{𝔣}^{𝔟}\,\xrightarrow{\sim}\,A_{𝔣𝔞}.}
\DAnchor{14}{240}
We now allow \(𝔣\) to vary. Clearly, if \(𝔣'\subseteq 𝔣\) and if \(𝔞\sim _{𝔣'}𝔟\), then \(𝔞\sim _{𝔣}𝔟\). There are thus natural maps \(I_{𝔣'}\to I_{𝔣}\), inducing maps \(A_{𝔣'}\to A_{𝔣}\) and \(G_{𝔣'}\to G_{𝔣}\). We define\par
\par
\DFormula{I=\varprojlim \,I_{𝔣},}
\DFormula{A=\varprojlim \,A_{𝔣}\,\subset \,I,}
\DFormula{G=\varprojlim \,G_{𝔣}\,\subset \,A.}
\par
(We note that \(G\) no longer has the meaning given it in the Introduction.) In the prohibited case \(K=\mathbb{Q}\), these limits become respectively \(\widehat{\mathbb{Q}}\), \(\widehat{\mathbb{Z}}\), and \(\widehat{\mathbb{Z}}^{\ast}\). Also, we recognize \(G\) as the Galois group of the maximal abelian extension \(K^{ab}\) of \(K\). In what follows, we will gradually establish parallels between \(A\) and \(\widehat{\mathbb{Z}}\) and between \(I\) and \(\widehat{\mathbb{Q}}\). \(A\) typical element of \(I\) will be denoted \((𝔞_{𝔣})=𝔞\), indicating that \(𝔞_{𝔣}\) is the \(I_{𝔣}-component\) of \(𝔞\).\par
\par
For \(𝔟\in I_{0}\), we let\par
\par
\(I^{𝔟}=\lbrace (𝔞_{𝔣})\in I\,\vert \,𝔞_{𝔣}\subseteq 𝔟\,\text{for}\,all\,𝔣\subseteq 𝔟\rbrace \)\par
\(=\lbrace (𝔞_{𝔣})\in I\,\vert \,𝔞_{𝔣}\subseteq 𝔟\,\text{for}\,some\,𝔣\subseteq 𝔟\rbrace \);\par
\par
the equivalence of the two definitions resulting from (2.2). It is clear that we have\par
\par
\DFormula{I=\bigcup _{𝔞\in A_{0}}\,I^{𝔞^{-1}}=\bigcup _{𝔟\in I_{0}}\,I^{𝔟}.}
\par
(2.9) Proposition. Let \(𝔟_{1}\) and \(𝔟_{2}\) be ideals. Suppose that \(𝔣_{1}\subseteq 𝔟_{1}\) and \(𝔣_{2}\subseteq 𝔟_{2}\) are conductors such that\par
\par
\DFormula{𝔣_{1}𝔟_{2},\,𝔣_{2}𝔟_{1}\subseteq \mathcal{O}.}
\par
\subsection*{Then multiplication of ideals induces a map}
\par
\DFormula{I_{𝔣_{1}}^{𝔟_{1}}\times I_{𝔣_{2}}^{𝔟_{2}}\to I_{𝔣_{1}𝔟_{2}+𝔣_{2}𝔟_{1}}^{𝔟_{1}𝔟_{2}}.}
\par
Using (2.9), we find a multiplication\par
\par
\DFormula{I\times I\to I.}
\par
Specifically, to calculate a product \(𝔞\cdot 𝔟\) mod \(𝔣\), suppose that \(𝔞\in I^{𝔟_{1}}\), \(𝔟\in I^{𝔟_{2}}\). If \(𝔣_{1}\subseteq 𝔟_{1}\) and \(𝔣_{2}\subseteq 𝔟_{2}\) are such that \(𝔣_{1}𝔟_{2}\), \(𝔣_{2}𝔟_{1}\subseteq 𝔣\), then \(𝔞_{𝔣_{1}}\cdot 𝔟_{𝔣_{2}}\) gives \(𝔞𝔟\) mod \(𝔣\). It is clear that \(𝔟_{1}\cdot I^{𝔟_{2}}=I^{𝔟_{1}𝔟_{2}}\) for ideals \(𝔟_{1}\) and \(𝔟_{2}\). (We write again \(𝔟\) for the image of an ideal \(𝔟\) in \(I.)\) Namely, \(𝔞_{𝔣}\subseteq 𝔟_{2}\) for all \(𝔣\) sufficiently small if and only if \(𝔟_{1}𝔞_{𝔣}\subseteq 𝔟_{1}𝔟_{2}\) for all \(𝔣\) sufficiently small.\par
\par
We observe, incidentally, that the natural map \(I_{0}\to I\) is injective. In fact, if \(𝔞\) and \(𝔟\) are ideals which are equivalent modulo a conductor \(𝔣\subseteq 𝔞\cap 𝔟\), then by (2.2) we have \(𝔞\subseteq 𝔟\) and \(𝔟\subseteq 𝔞\).\par
\par
(2.10) Proposition. Let \(I^{\ast}\) be the set of invertible elements of \(I\) under the multiplication. Then \(I^{\ast}=I_{0}\times G\).\par
\par
Proof. It is immediate from (2.6) that if we have \(𝔞\cdot 𝔟=1\) with \(𝔞,𝔟\in A\), then both \(𝔞\) and \(𝔟\) belong to \(G\). In general, if \(𝔞\cdot 𝔟=1\) with \(𝔞,𝔟\in I\), multiply \(𝔞\) by an ideal\par
\DAnchor{15}{241}
(and \(𝔟\) by its inverse) to have \(𝔞\in A\). We need not have \(𝔟\in A\); let \(𝔡\in A_{0}\) be the denominator of \(𝔟\). Then \(𝔟\in I^{𝔡^{-1}}\). Calculating \(𝔞\cdot 𝔟\) mod 1, we find \(1\sim 𝔞_{𝔡}𝔟_{1}\), giving that \(𝔞_{𝔡}\) is divisible by \(𝔡\). Hence \(𝔞\in I^{𝔡}\), so that in the expression\par
\par
\DFormula{(𝔡^{-1}𝔞)\cdot (𝔡𝔟)=1,}
\par
both terms in the product lie in \(A\). Thus \(𝔡^{-1}𝔞\in G\), so \(𝔞\in I_{0}\times G\).\par
\par
[It is obvious that the intersection \(I_{0}\cap G\) is trivial because an element \((𝔞_{𝔣})\) of \(G\) is such that \(𝔞_{𝔣}\in A_{0}\), \((𝔞_{𝔣},𝔣)=1\) for each \(𝔣.]\)\par
\par
(2.11) Proposition. Let f be a positive integer, and suppose that \(𝔞,𝔟\in I_{0}\) are congruent mod \(\sim _{f}\) (i.e., mod \(\sim _{(f)})\). Let \(𝔡\) be a denominator for \(𝔞\) (and hence for \(𝔟)\). Then\par
\par
\DFormula{\mathcal{N}𝔞\equiv \mathcal{N}𝔟\,\operatorname{mod}\,(\mathcal{N}𝔡^{-1})f,}
\par
i.e.,\par
\par
\(\mathcal{N}𝔞-\mathcal{N}𝔟\in (\mathcal{N}𝔡^{-1})f\mathbb{Z}\).\par
\par
(2.12) Corollary. The norm map \(I_{0}\to \mathbb{N}^{+}\) extends to a continuous map \(\mathcal{N}:I\to \widehat{\mathbb{Q}}\), such that \(\mathcal{N}(𝔞𝔟)=\mathcal{N}(𝔞)\mathcal{N}(𝔟)\). We have \(\mathcal{N}(A)\subseteq \widehat{\mathbb{Z}}\), \(\mathcal{N}(G)\subseteq \widehat{\mathbb{Z}}^{\ast}\).\par
\par
Proof of (2.11). We may assume that \(𝔡=1\). Indeed, under the hypothesis that \(𝔡\) is a denominator, the ideals \(𝔞𝔡\) and \(𝔟𝔡\) are integral and congruent mod f (in fact, mod \(f𝔡)\). The assertion for \(𝔡=1\) gives \(\mathcal{N}(𝔞𝔡)\equiv \mathcal{N}(𝔟𝔡)\) mod f, as desired.\par
\par
Now with \(𝔞,𝔟\in A_{0}\), set \(𝔡=(f,𝔞)=(f,𝔟)\). We have \(𝔞𝔡^{-1}\sim 𝔟𝔡^{-1}\) mod \(f𝔡^{-1}\); let \(f'\) be the largest integer dividing the ideal \(f𝔡^{-1}\). Then f divides \(f'\cdot \mathcal{N}𝔡\), so that the assertion\par
\par
\DFormula{\mathcal{N}(𝔞𝔡^{-1})\equiv \mathcal{N}(𝔟𝔡^{-1})\,\operatorname{mod}\,f'}
\par
implies the proposition. We are thus reduced to the case where \((f,𝔞)=(f,𝔟)=1\).\par
\par
Then, finally, let \(\alpha \) be as in the definition of \(\sim _{f}\) for \(𝔞,𝔟\). We have \(\mathcal{N}𝔞=\mathcal{N}\alpha \cdot \mathcal{N}𝔟\), with \(\alpha \equiv 1\) \(\operatorname{mod}^{\times }\) f. We see immediately that \(\mathcal{N}\alpha \equiv 1\) \(\operatorname{mod}^{\times }\) f, giving\par
\par
\DFormula{\mathcal{N}𝔞\equiv \mathcal{N}𝔟\,(\operatorname{mod}\,f).}
\par
We now wish to compare the space \(I\) with the ring \(\widehat{K}\) of finite adeles of \(K\). For this, we regard \(\widehat{K}\) as the inverse limit \(\varprojlim \) \(K/𝔣\) and consider an element \(\alpha \) of \(\widehat{K}\) as a sequence \((\alpha _{𝔣})\) of totally positive numbers in \(K\), compatible with the transition maps \(K/𝔣'\to K/𝔣\) for \(𝔣'\subseteq 𝔣\). Let \(i(\alpha )\) be the sequence of principal ideals \(((\alpha _{𝔣}))\).\par
\par
(2.13) Proposition. \(i(\alpha )\) is an element of \(I\) which depends only on \(\alpha \), and not on the choice of the \(\alpha _{𝔣}\).\par
\par
Proof. This follows immediately from the following lemma, which in turn is an immediate consequence of the definition of \(\sim _{𝔣}\).\par
\par
(2.14) Lemma. Let \(\alpha ,\beta \in K\) with \(\alpha ,\beta \gg 0\). Then \((\alpha )\sim (\beta )\) mod \(𝔣\) if and only if there is a totally positive unit \(u\in K\) such that \(\alpha \equiv u\beta \) mod \(𝔣\).\par
\par
Now let \(\widehat{U}^{+}\) be the closure in \(\widehat{\mathcal{O}}^{\ast}\) of the group of totally positive units of \(K\).\par
\DAnchor{16}{242}
(2.15) Proposition. Let \(\alpha ,\beta \in \widehat{K}\). Then \(i(\alpha )=i(\beta )\) if and only if \(\alpha =\beta u\) for some \(u\in \widehat{U}^{+}\).\par
\par
Proof. The lemma shows that if \(\alpha =\beta u\) with \(u\in \widehat{U}^{+}\), then \((\alpha _{𝔣})\sim (\beta _{𝔣})\) mod \(𝔣\) for all \(𝔣\), so \(i(\alpha )=i(\beta )\). Conversely, suppose \(i(\alpha )=i(\beta )\), and let us prove that \(\alpha =\beta u\) for some u. We do have \(\alpha _{𝔣}\equiv u_{𝔣}\beta _{𝔣}\) mod \(𝔣\) for each \(𝔣\), where \(u_{𝔣}\) is a totally positive unit. However, the sequence \((u_{𝔣}\) mod \(𝔣)\) may not be compatible because the units \(u_{𝔣}\) may not be unique mod \(𝔣\).\par
\par
To deal with this problem, we multiply \(\alpha \) and \(\beta \) by a positive integer so as to have \(\alpha ,\beta \in \widehat{\mathcal{O}}\). (This changes neither the hypothesis nor the conclusion of the proposition.) For each \(𝔣\), let \(𝔡_{𝔣}=gcd(𝔣,\alpha _{𝔣})=gcd(𝔣,\beta _{𝔣})\). Then the unit \(u_{𝔣}\) is determined modulo \(𝔣𝔡_{𝔣}^{-1}=𝔤_{𝔣}\). For \(𝔣'\subset 𝔣\), we have \(𝔤_{𝔣'}\subset 𝔤_{𝔣}\). Thus we need only to know that \(\widehat{U}^{+}\) maps into \(\varprojlim _{𝔣}\) \(\mathbb{C}_{𝔣}\), where \(\mathbb{C}_{𝔣}\) is the group of totally positive units taken mod \(𝔤_{𝔣}\). This is obvious, for example because \(\widehat{U}^{+}\) is compact and maps onto each \(\mathbb{C}_{𝔣}\).\par
\par
(2.16) Proposition. \(\mathcal{N}(i(\alpha ))=\mathcal{N}\alpha \) for all \(\alpha \in \widehat{K}\).\par
\par
Proof. Indeed, \(\mathcal{N}((\alpha _{𝔣}))=\mathcal{N}\alpha _{𝔣}\) because \(\alpha _{𝔣}\gg 0\).\par
\par
Convention. Let \(𝔞\in I\), \(\alpha \in \widehat{K}\). We write \((𝔞\cdot \alpha )\), or simply \(𝔞\cdot \alpha \), for the product \(𝔞\cdot i(\alpha )\) in \(I\).\par
\par
Note that for \(\alpha \gg 0\) in \(K\) we have \((\alpha )=i(\alpha )\) in \(I\). Thus \((\alpha )\cdot \alpha ^{-1}=1\).\par
\par
(2.17) Lemma. For \(𝔞\in I\), the quantity \((𝔞\cdot 0)\) depends only on the strict ideal class of \(𝔞_{1}\): for \(𝔟\in I\), we have \((𝔞\cdot 0)=(𝔟\cdot 0)\) if and only if \(𝔞_{1}\) and \(𝔟_{1}\) belong to the same strict ideal class of \(K\).\par
\par
Proof. For \(t\gg 0\), we have \((t𝔞\cdot 0)=(𝔞\cdot t0)=(𝔞\cdot 0)\). Hence we can assume that \(𝔞,𝔟\in A\), in which case the assertion to be proved is: \((𝔞\cdot 0)=(𝔟\cdot 0)\) if and only if \(𝔞_{1}\sim 𝔟_{1}\) mod 1. For each \(𝔣\), let \(\alpha _{𝔣}\) be a totally positive number divisible by \(𝔣\). Modulo \(𝔣\), \((𝔞\cdot 0)\) is \(𝔞_{𝔣}\alpha _{𝔣}\) and \((𝔟\cdot 0)\) is \(𝔟_{𝔣}\alpha _{𝔣}\). By (2.7) we see that \(𝔞_{𝔣}\alpha _{𝔣}\sim 𝔟_{𝔣}\alpha _{𝔣}\) mod \(𝔣\) if and only if \(𝔞_{𝔣}\alpha _{𝔣}\sim 𝔟_{𝔣}\alpha _{𝔣}\) mod 1, which is true if and only if \(𝔞_{1}\) and \(𝔟_{1}\) are in the same strict ideal class.\par
\par
(2.18) Proposition. Let \(𝔞\in I\). Suppose that \((𝔞\cdot 0)=(1\cdot 0)\). Then \(𝔞=(1\cdot \alpha )\) for some \(\alpha \in \widehat{K}\).\par
\par
Proof. As usual, we can multiply \(𝔞\) by some \(t\gg 0\) to assume that \(𝔞\in A\). Then by hypothesis we have, for each \(𝔣\), \(𝔞_{𝔣}=(\alpha _{𝔣})\), where \(\alpha _{𝔣}\in K\) is a totally positive integer, which by (2.14) is well defined, mod \(𝔣\), modulo the action of the group of totally positive units. Let \(\mathbb{C}_{𝔣}\) be the set of integers mod \(𝔣\), taken modulo this action. Then the compact \(\widehat{\mathcal{O}}\) maps onto each \(\mathbb{C}_{𝔣}\) and hence on to the limit \(\varprojlim \) \(\mathbb{C}_{𝔣}\). This gives what is needed.\par
\par
(2.19) Corollary. Let \(𝔞=(𝔞_{𝔣})\in I\), and let \(𝔟\in I_{0}\) be an ideal in the same ideal class as \(𝔞_{1}\). Then \(𝔞=(𝔟\cdot \alpha )\) for some \(\alpha \in \widehat{K}\).\par
\par
Proof. By (2.17), \((𝔞𝔟^{-1}\cdot 0)=(1\cdot 0)\). Hence \(𝔞𝔟^{-1}\) may be written \((1\cdot \alpha )\) for some \(\alpha \in \widehat{K}\).\par
\DAnchor{17}{243}
We recall that for \(x\) an invertible element of \(\widehat{K}\), we denote by (x) the ideal “generated” by \(x\).\par
\par
(2.20) Proposition. For each \(x\in \widehat{K}^{\ast}\), we have \((x)\cdot x^{-1}\in G\).\par
\par
Proof. It suffices to prove that \((x)\cdot x^{-1}\in A\), using the symmetry and the fact that an element of \(A\) with an inverse in \(A\) must in fact belong to \(G\). This follows from the following result, whose proof we omit.\par
\par
(2.21) Lemma. Let \(𝔞,𝔟\in I_{0}\), and \(\alpha \in \widehat{K}\). Then\par
\par
\DFormula{(𝔞\cdot \alpha )\in I^{𝔟}}
\par
if and only if \(\alpha \) belongs to the completion \(\widehat{𝔞^{-1}𝔟}\) of the ideal \(𝔞^{-1}𝔟\).\par
\par
As a complement to (2.21), we mention the following fact:\par
\par
(2.22) Let \(𝔞\) be an ideal and let \(\alpha ,\beta \in \widehat{K}\). We have \((𝔞\cdot \alpha )\sim (𝔞\cdot \beta )\) mod \(𝔣\) if and only if there is a totally positive unit u such that we have\par
\par
\DFormula{\alpha -u\beta \in \widehat{𝔞^{-1}𝔣}.}
\par
Referring now to (2.20), we let \(j(x)=(x)\cdot x^{-1}\) for \(x\in \widehat{K}^{\ast}\). Then j may be viewed as a homomorphism\par
\par
\DFormula{j:\widehat{K}^{\ast}\to G}
\par
which is trivial on the set \(K^{\gg 0}\) of totally positive elements of \(K\). On the other hand, for each conductor \(𝔣\) there is a natural map\par
\par
\DFormula{\psi _{𝔣}:\widehat{K}^{\ast}\to G_{𝔣},}
\par
trivial on \(K^{\gg 0}\), which maps each \(x\in \widehat{K}^{\ast}\) such that \(x\equiv 1\) \(\operatorname{mod}^{\times }\) \(𝔣\) to the class in \(G_{𝔣}\) of the ideal (x) generated by \(x\), cf. [22, pp. 146--147]. (Here we should recall that \(G_{𝔣}\) is the ray class group of \(K\) mod \(𝔣\), so that each \(prime-to-𝔣\) ideal of \(K\), integral or not, has a well defined image in \(G_{𝔣}\). From the point of view of the equivalence \(\sim _{𝔣}\), only the integral \(prime-to-𝔣\) ideals map to \(G_{𝔣}.)\) The resulting map\par
\par
\DFormula{\psi :\widehat{K}^{\ast}\to G=\varprojlim \,G_{𝔣}}
\par
is surjective. By abuse of language, we may refer to it as an Artin symbol.\par
\par
(2.33) Proposition. We have \(\psi =j\).\par
\par
Proof. Since both \(\psi \) and j vanish on \(K^{\gg 0}\), it suffices to check that \(\psi _{𝔣}\) and the composition of j and the quotient \(G\to G_{𝔣}\) agree on elements of the set \(1+\widehat{𝔣}\). Now, mod \(𝔣\), j(x) is represented by any ideal of the form \((\alpha )(x)\), where \(\alpha \) is a totally positive element of \(K\) with \(\alpha -x^{-1}\in \widehat{𝔣}\). We must verify that\par
\par
\DFormula{(\alpha )(x)\sim _{𝔣}(x),}
\par
and it suffices to see that\par
\par
\DFormula{\alpha \in 1+\widehat{(x)^{-1}𝔣}.}
\DAnchor{18}{244}
This, finally, follows from the defining property of \(\alpha \), together with the fact that\par
\par
\DFormula{1-x^{-1}=x^{-1}(x-1)\in x^{-1}\widehat{𝔣}.}
\par
(2.24) Frobenius Elements\par
\par
Let \(𝔣\) be a conductor, and let v be a real place of \(K\). Let \(\alpha \in 1+𝔣\) be a number which is negative at v and positive at each real place of \(K\) different from v. The class \(\sigma _{v}\) of \((\alpha )\) in \(G_{𝔣}\) is independent of v and has order 1 or 2. It is trivial if and only if \(\alpha \) may be chosen to be a unit. If we interpret \(G_{𝔣}\) as the Galois group over \(K\) of the ray class field \(K_{𝔣}\) of \(K\) of conductor \(𝔣\), then \(\sigma _{v}\) becomes complex conjugation after we choose an embedding \(K_{𝔣}\to \mathbb{C}\) which induces v on \(K\).\par
\par
For fixed v and varying \(𝔣\), the elements \(\sigma _{v}\) piece together to give a Frobenius element in \(G\), which we again call \(\sigma _{v}\). Each \(\sigma _{v}\) now has order 2, and in fact the subgroup \(\Sigma ^{\pm }\) of \(G\) generated by the \(\sigma _{v}\) is an elementary 2-group whose order is \(2^{r}\). Each \(\sigma _{v}\) has norm \(-1\). If we let \(\Sigma \) be the kernel of the norm map\par
\par
\DFormula{\mathcal{N}:\Sigma ^{\pm }\to \lbrace \pm 1\rbrace ,}
\par
then \(\Sigma \) has order \(2^{r-1}\) and is generated by products \(\sigma _{v}\sigma _{w}\), with v and w running over the real places of \(K\). For \(𝔣\) a conductor, we let \(\Sigma _{𝔣}^{\pm }\) (resp. \(\Sigma _{𝔣})\) be the image of \(\Sigma ^{\pm }\) (resp. \(\Sigma )\) in \(G_{𝔣}\). We will be especially interested in the case \(𝔣=(1)\).\par
\par
In a natural way, the \(\sigma _{v}\) lead to the notion of functions with parity. Let \(\varphi \) be a function on \(G\) with values in an abelian group \(\mathcal{V}\). For each real place v, let \(a_{v}\) be one of the two integers 0, 1. We say that \(\varphi \) has parity \((a_{v})\) if we have\par
\par
\(\varphi (\sigma _{v}g)=(-1)^{a_{v}}\varphi (g)\)\par
\par
for all g,v. When the \(a_{v}\) are all equal to 0 (resp. 1) we say that \(\varphi \) is even (resp. odd). The phrase “\(\varphi \) has parity \((-1)^{k}\)” means that \(\varphi \) is odd if \(k\) is odd and even if \(k\) is even.\par
\par
(2.25) Invertible Modules\par
\par
As is well known, the group of ideal classes of \(K\) may be interpreted as the group of isomorphism classes of invertible \(\mathcal{O}\)-modules. We shall now similarly interpret the constructions of this §. For this, we must introduce the notion of a signed \(\mathcal{O}\)-module.\par
\par
Let \(\mathcal{L}\) be an invertible \(\mathcal{O}\)-module. For each real place v of \(K\), let \(\mathcal{L}_{v}\) be the tensor product \(\mathcal{L}\otimes _{\mathcal{O}}\) \(\mathbb{R}\), with \(\mathbb{R}\) viewed as an \(\mathcal{O}\)-module via v. Then \(\mathcal{L}_{v}\) is a free \(\mathbb{R}\)-module of rank 1. \(A\) positivity for \(\mathcal{L}\) is the choice, for each v, of an element of\par
\par
\DFormula{\operatorname{Isom}(\mathcal{L}_{v},\mathbb{R})\cong \mathbb{R}^{\ast},}
\par
taken modulo the action of the group of positive real numbers. There are precisely \(2^{r}\) possible positivities for \(\mathcal{L}\).\par
\par
\(A\) pair consisting of an invertible \(\mathcal{O}\)-module \(\mathcal{L}\), together with a positivity \(+\) for \(\mathcal{L}\), is a signed \(\mathcal{O}\)-module. Two such modules \((\mathcal{L},+)\), \((\mathcal{L}',+')\) are isomorphic if there is an \(\mathcal{O}-isomorphism\) \(\mathcal{L}\xrightarrow{\sim}\mathcal{L}'\) compatible with \(+\), \(+'\). It is clear that the isomorphism classes of signed \(\mathcal{O}\)-modules are the strict ideal classes of \(K\). More precisely, this works as follows. Each ideal \(𝔞\) of \(K\) has a canonical positivity\par
\DAnchor{19}{245}
\(+_{\mathrm{can}}\) arising from the canonical isomorphism \(𝔞\otimes _{\mathcal{O}}\) \(K=K\). With \(𝔞\) and \(𝔟\) ideals, we have \((𝔞,+_{\mathrm{can}})\cong (𝔟,+_{\mathrm{can}})\) if and only if \(𝔞\) and \(𝔟\) are in the same strict class. Finally, any pair \((\mathcal{L},+)\) is isomorphic to \((𝔞,+_{\mathrm{can}})\) for some ideal \(𝔞\).\par
\par
Now let \(𝔣\) be a conductor. We consider triples \((\mathcal{L},+,\varphi )\), where \(\varphi \) is an \(\mathcal{O}\)-linear map \(\mathcal{L}\to K/𝔣\). Two such triples \((\mathcal{L},+,\varphi )\), \((\mathcal{L}',+',\varphi ')\) are said to be isomorphic if there is an isomorphism \((\mathcal{L},+)\xrightarrow{\sim}(\mathcal{L}',+')\) so that the diagram\par
\par
\begin{center}\begin{tikzcd}[column sep=large,row sep=small] \mathcal L \arrow[dd] \arrow[dr,"\varphi"] & \\ & K/\mathfrak f \\ \mathcal L' \arrow[ur,"\varphi'"'] & \end{tikzcd}\end{center}
\par
is commutative. For \(𝔞\in I_{0}\), we let \(\varphi _{\mathrm{can}}\) be the map \(𝔞\hookrightarrow K\to K/𝔣\). \(A\) computation shows that \((𝔞,+_{\mathrm{can}},\varphi _{\mathrm{can}})\) and \((𝔟,+_{\mathrm{can}},\varphi _{\mathrm{can}})\) are isomorphic if and only if \(𝔞\sim _{𝔣}𝔟\).\par
\par
Furthermore, let \((\mathcal{L},+,\varphi )\) be a given triple, and let \(𝔞\) be an ideal such that \((\mathcal{L},+)\) and \((𝔞,+_{\mathrm{can}})\) are isomorphic. Choosing an isomorphism we find a map\par
\par
\DFormula{𝔞\xrightarrow{\sim}\mathcal{L}\,\xrightarrow{\varphi }\,K/𝔣,}
\par
necessarily given by a multiplication \(\gamma :𝔞\to \gamma 𝔞\) (with \(\gamma \in K)\), followed by the reduction \(K\to K/𝔣\). We may choose \(\gamma \) to be totally positive. Then the isomorphism\par
\par
\DFormula{𝔞\gamma \,\xrightarrow{\gamma ^{-1}}\,𝔞\cong \mathcal{L}}
\par
shows that we have \((\mathcal{L},+,\varphi )\cong (𝔞\gamma ,+_{\mathrm{can}},\varphi _{\mathrm{can}})\). Hence \(I_{𝔣}\) is just the set of isomorphism classes of triples \((\mathcal{L},+,\varphi )\).\par
\par
To describe \(I\), we consider triples \((\mathcal{L},+,\varphi )\), where \(\varphi \) is now an \(\mathcal{O}\)-linear map \(\mathcal{L}\to \widehat{K}\). Such triples give rise to elements of \(I\) as follows. Choosing an isomorphism \((\mathcal{L},+)\cong (𝔞,+_{\mathrm{can}})\), we find an \(\mathcal{O}\)-linear \(𝔞\to \widehat{K}\), which is necessarily of the form “multiplication by \(\gamma \)” for some \(\gamma \in \widehat{K}\). The product \((𝔞\cdot \gamma )\) in \(I\) is easily seen to depend only on \((\mathcal{L},+,\varphi )\). Note that the image mod \(𝔣\) of this product is represented by \((\mathcal{L},+)\) together with the map “\(\varphi \) mod \(𝔣\)”:\par
\par
\DFormula{\mathcal{L}\,\xrightarrow{\varphi }\,\widehat{K}\to \widehat{K}/\widehat{𝔣}=K/𝔣.}
\par
In either situation, the action of a real Frobenius \(\sigma _{v}\) is as follows: we have\par
\par
\DFormula{\sigma _{v}\cdot (\mathcal{L},+,\varphi )=(\mathcal{L},+',\varphi ),}
\par
where \(+'\) is deduced from \(+\) by changing \(+\) at v.\par
\par
\subsection*{§3. \(A\) Functional Equation}
\par
In this § we define the L-series attached to a (complex valued) Schwartz function on the measure space \(I\). We then prove a functional equation for such L-series, involving a Fourier transform on the space \(I\). Defining this Fourier transform is the first order of business.\par
\par
We shall view the ring \(𝔸\) of adeles of \(K\) as the product \(𝔸_{f}\times K_{\infty }\), where \(𝔸_{f}=\widehat{K}\) is the ring of finite adeles of \(K\) and \(K_{\infty }=K\otimes \mathbb{R}\) is the product of the real completions of \(K\). For computing, we number these real places of \(K\), so that \(K_{\infty }=\mathbb{R}^{r}\). Also, for \(x\in 𝔸\), we write \(x_{f}\) (resp. \(x_{\infty })\) for the image of \(x\) in \(𝔸_{f}\) (resp. \(K_{\infty })\).\par
\par
Let \(\vert \) \(\vert :K_{\infty }\to \mathbb{R}\) be the standard absolute value, namely\par
\par
\DFormula{(x_{1},...,x_{r})\mapsto \prod \vert x_{i}\vert .}
\DAnchor{20}{246}
We write instead \(\Vert \) \(\Vert \) for the standard absolute value on \(𝔸_{f}\), defined for example by the equation\par
\par
\(d(\gamma x)=\Vert \gamma \Vert \cdot dx\),\par
\par
where dx is a Haar measure on \(𝔸_{f}\). For \(\gamma \in \widehat{K}^{\ast}\), we have \(\Vert \gamma \Vert =\mathcal{N}(\gamma )^{-1}\), where \((\gamma )\) is the ideal of \(K\) generated by \(\gamma \).\par
\par
Let \(\psi :𝔸\to \mathbb{C}^{\ast}\) be the standard additive character, trivial on \(K\), whose restriction to a completion \(K_{v}\) of \(K\) is given as follows:\par
\par
For v real, \(\psi :x\mapsto e^{2\pi ix}\);\par
\par
For v \(p\)-adic, \(\psi :x\mapsto e^{-2\pi i\cdot \operatorname{tr}\,x}\),\par
\par
where \(\operatorname{tr}:K_{v}\to \mathbb{Q}_{p}\) is the trace, and the exponential \(e^{2\pi it}\) for \(t\in \mathbb{Q}_{p}\) is defined in the usual way. The restriction of \(\psi \) to \(𝔸_{f}\) (resp. \(K_{\infty })\) is denoted \(\psi _{f}\) (resp. \(\psi _{\infty })\), or simply \(\psi \), according to the context. For \(\varepsilon \) a Schwartz function on \(𝔸_{f}\) (resp. \(K_{\infty }\), resp. \(𝔸)\), we define the Fourier transform \(\widehat{\varepsilon }\) of \(\varepsilon \) by the formula\par
\par
(3.1)    \(\widehat{\varepsilon }(y)=\int \varepsilon (x)\cdot \psi (xy)dx\).\par
\par
Here dx denotes Haar measure, normalized so that the formula\par
\par
\DFormula{\widehat{\widehat{\varepsilon }}(x)=\varepsilon (-x)}
\par
holds, cf. §3.3 and §4.1 of [39].\par
\par
We now turn our attention to Schwartz functions on \(\widehat{K}\) and on \(I\). We recall the following “variance” formula for \(\widehat{K}\):\par
\par
(3.2) If \(\eta (x)=\varepsilon (\gamma x)\) with \(\gamma \in \widehat{K}^{\ast}\), then\par
\par
\DFormula{\widehat{\eta }(y)=\widehat{\varepsilon }(\gamma ^{-1}y)\Vert \gamma \Vert ^{-1}.}
\par
Now let \(\varepsilon \) be a locally constant function with compact support on \(I\). Specifically, assume that \(\varepsilon \) is supported on \(I^{𝔟}\) and defined modulo \(\sim _{𝔣}\). If \(𝔲\in I_{0}\), let \(\varepsilon _{𝔲}:\widehat{K}\to \mathbb{C}\) be the function \(x\mapsto \varepsilon (𝔲\cdot x)\). By (2.21) and (2.22), \(\varepsilon _{𝔲}\) is supported on \(\widehat{𝔲^{-1}𝔟}\) and defined modulo \(\widehat{𝔲^{-1}𝔣}\). Let \(\widehat{\varepsilon }_{𝔲}\) denote the Fourier transform of \(\varepsilon _{𝔲}\). For \(\gamma \in K\) totally positive, one has the formula\par
\par
\DFormula{(3.3)\,\,\,\,\widehat{\varepsilon }_{\gamma 𝔲}(x)=\widehat{\varepsilon }_{𝔲}(\gamma ^{-1}x)\cdot \Vert \gamma \Vert ^{-1},}
\par
as follows easily from (3.2).\par
\par
We define the Fourier transform \(T\varepsilon :I\to \mathbb{C}\) of \(\varepsilon \) by the formula\par
\par
\DFormula{(3.4)\,\,\,\,(T\varepsilon )(𝔲^{-1}\cdot x)=(\mathcal{N}𝔲^{-1})\widehat{\varepsilon }_{𝔲}(x),}
\par
for \(𝔲\in I_{0}\), \(x\in \widehat{K}\). [We observe that each element of \(I\) is a product \(𝔲^{-1}\cdot x\), by (2.19).] To check that \(T\varepsilon \) is well defined, suppose that \(𝔲^{-1}\cdot \alpha =𝔳^{-1}\cdot \beta \); we then wish to verify that\par
\par
\DFormula{\mathcal{N}𝔲^{-1}\cdot \widehat{\varepsilon }_{𝔲}(\alpha )=\mathcal{N}𝔳^{-1}\cdot \widehat{\varepsilon }_{𝔳}(\beta ).}
\par
The equation \(𝔲^{-1}\cdot \alpha =𝔳^{-1}\cdot \beta \) gives \(𝔲^{-1}\cdot 0=𝔳^{-1}\cdot 0\), so that \(𝔳=\gamma 𝔲\) for some \(\gamma \gg 0\) by (2.17).\par
\DAnchor{21}{247}
We find that \((𝔲^{-1}\cdot \alpha )=(𝔲^{-1}\cdot \gamma ^{-1}\beta )\), so that\par
\par
\DFormula{\alpha =\gamma ^{-1}\beta t}
\par
for some \(t\in \widehat{U}^{+}\) (2.15). By (3.3), we have\par
\par
\DFormula{\widehat{\varepsilon }_{𝔳}(\beta )/\mathcal{N}𝔳\,=\,\widehat{\varepsilon }_{𝔲}(\gamma ^{-1}\beta )/\mathcal{N}𝔲\,=\,\widehat{\varepsilon }_{𝔲}(\alpha t^{-1})/\mathcal{N}𝔲.}
\par
The required formula then follows from (3.2), since \(\varepsilon _{𝔲}\) is invariant under \(x\mapsto xt\), and \(\Vert t\Vert =1\).\par
\par
We now give some formal properties of \(T\) that follow directly from the definition.\par
\par
\DFormula{(3.5)\,\,\,\,T(T\varepsilon )(x)=\varepsilon (x\cdot -1).}
\par
Proof. Let \(\eta =T\varepsilon \). For \((𝔲^{-1}\cdot \alpha )\in I\), we have\par
\par
\((T\eta )(𝔲^{-1}\cdot \alpha )=1/(\mathcal{N}𝔲)\) \(\int \eta _{𝔲}(x)\psi (x\alpha )dx\)\par
\(=\int \widehat{\varepsilon }_{𝔲^{-1}}(x)\psi (x\alpha )dx\)\par
\DFormula{=\widehat{\widehat{\varepsilon }}_{𝔲^{-1}}(\alpha )=\varepsilon _{𝔲^{-1}}(-\alpha )}
\DFormula{=\varepsilon (𝔲^{-1}\cdot -\alpha )=\varepsilon ((𝔲^{-1}\cdot \alpha )\cdot -1).}
\par
(3.6) If \(\eta (x)=\varepsilon (𝔟x)\) for some \(𝔟\in I_{0}\), then\par
\par
\DFormula{(T\eta )(y)=(T\varepsilon )(𝔟^{-1}y)\cdot \mathcal{N}𝔟.}
\par
Proof. We have\par
\par
\(\mathcal{N}𝔲\cdot T\eta (𝔲^{-1}\cdot \alpha )=\int \eta (𝔲\cdot x)\psi (x\alpha )dx\)\par
\(=\int \varepsilon (𝔲𝔟\cdot x)\psi (x\alpha )dx\)\par
\DFormula{=(T\varepsilon )((𝔲𝔟)^{-1}\cdot \alpha )\cdot \mathcal{N}(𝔲𝔟).}
\par
(3.7) Corollary. Suppose that \(\eta (x)=\varepsilon (cx)\) with \(c\in G\). Then \((T\eta )(y)=(T\varepsilon )(c^{-1}y)\).\par
\par
Proof. As noted in §2, each \(c\in G\) may be written \((\gamma )\cdot \gamma ^{-1}\) for some \(\gamma \in \widehat{K}^{\ast}\). Since \(\Vert \gamma \Vert \cdot \mathcal{N}(\gamma )=1\), the result follows from (3.6), (3.3).\par
\par
Using the above properties, we will now verify that \(T\varepsilon \) is locally constant and compactly supported. For the application that we ultimately have in mind, the ideal \(𝔟\) governing the support of \(\varepsilon \) is (1): that is, \(\varepsilon \) is supported on \(A\). Let us assume that this is the case, for simplicity. [If \(𝔟\neq (1)\), one may use (3.6) to reduce to the case \(𝔟=(1).]\) We suppose (as above) that \(\varepsilon \) is defined mod \(\sim _{𝔣}\).\par
\par
Let \(\eta \) be defined by the formula\par
\par
\DFormula{\eta (x)=(T\varepsilon )(𝔇^{-1}x).}
\par
(3.8) Theorem. The function \(\eta \) is supported on \(I^{𝔣^{-1}}\) and defined mod \(\sim 1\).\par
\par
Proof. We first verify the periodicity. If \(x,y\in I\) are congruent mod 1, we may write \(x=(𝔲^{-1}\cdot \alpha )\), \(y=(𝔲^{-1}\cdot \beta )\) for some \(𝔲\in I_{0}\) and elements \(\alpha ,\beta \) of \(\widehat{K}\) with \(\alpha -\beta \in \widehat{𝔲}\).\par
\DAnchor{22}{248}
Thus the periodicity property asserted for \(T\eta \) reduces to the statement that \(\widehat{\varepsilon }_{𝔲𝔇}\) is periodic mod \(\widehat{𝔲}\). As already noted, \(\varepsilon _{𝔲𝔇}\) is supported on \(𝔲^{-1}𝔇^{-1}\). Looking at the formula defining \(\widehat{\varepsilon }_{𝔲𝔇}\), we see that the required periodicity follows from the fact that the additive character \(\psi \) is trivial on \(𝔇^{-1}\).\par
\par
For the question of support, it is now enough to check that \(\eta \) vanishes on each ideal \(𝔲\) with \(𝔲\not\subset𝔣^{-1}\). (This remark follows from the periodicity and the fact that for \(x\sim y\) mod 1 we have\par
\par
\DFormula{x\in I^{𝔣^{-1}}\,\Leftrightarrow \,y\in I^{𝔣^{-1}}.)}
\par
So we must check that \((T\varepsilon )(𝔲)=0\) provided that \(𝔲\not\subset𝔣^{-1}𝔇^{-1}\). The vanishing of \((T\varepsilon )(𝔲)\) amounts to the vanishing of the integral\par
\par
\(\int \varepsilon (𝔲^{-1}\cdot x)\psi (x)dx\).\par
\par
But the hypothesis on \(𝔲\) insures the existence of an \(\alpha \in \widehat{𝔲𝔣}\) with \(\psi (\alpha )\neq 1\). The integral thus vanishes because of the formula\par
\par
\DFormula{\varepsilon (𝔲^{-1}\cdot x)=\varepsilon (𝔲^{-1}\cdot (x+\alpha )),}
\par
cf. (2.22).\par
\par
(3.9) Corollary. The function \(T\varepsilon \) is supported on \(I^{𝔣^{-1}𝔇^{-1}}\). It is defined modulo \(𝔇^{-1}\) in the sense that \((T\varepsilon )(x)=(T\varepsilon )(y)\) whenever \(𝔇x\sim 𝔇y\) mod 1.\par
\par
(3.10) Remark. The introduction of \(\eta \) is a consequence of our failure to have defined \(\sim _{𝔣}\) for \(𝔣\) a fractional ideal. For later use we now define two elements x,y of \(I\) to be congruent mod \(𝔇^{-1}\) if \(𝔇x\) and \(𝔇y\) are congruent mod 1.\par
\par
We now begin our study of L-functions. Let \(\varepsilon :I\to \mathbb{C}\) be locally constant and compactly supported. We set\par
\par
\DFormula{L(s,\varepsilon )=\sum _{𝔞}\,\varepsilon (𝔞)\mathcal{N}𝔞^{-s}}
\par
for \(s\) with large real part, where the sum runs over all (fractional) ideals of \(K\). It is easy to rewrite \(L(s,\varepsilon )\) as a linear combination of functions of the form\par
\par
\DFormula{\mathcal{N}𝔟^{s}\,\zeta _{𝔣}(s,c),}
\par
where \(𝔟\) is some ideal of \(K\) and \(\zeta _{𝔣}(s,c)\) is the partial zeta function of a class c mod \(𝔣\) in the sense of Siegel. Using results of Hecke, we may thus admit a priori that \(L(s,\varepsilon )\) may be continued to a meromorphic function on \(\mathbb{C}\), with at worst a simple pole at \(s=1\) and with no other poles. We are interested in obtaining a functional equation linking \(L(s,\varepsilon )\) and \(L(s,T\varepsilon )\).\par
\par
For this, we shall now assume that \(\varepsilon \) has parity \((a_{v})\) for some collection of integers \(a_{v}=0\) or 1, cf. (2.24). By (3.7), \(T\varepsilon \) again has parity \((a_{v})\). Let \(\Gamma _{\mathbb{R}}(s)=\pi ^{-s/2}\Gamma (s/2)\), and set\par
\par
\DFormula{(3.11a)\,\,\,\,\gamma (s)=\prod _{v}\,\Gamma _{\mathbb{R}}(s+a_{v}),}
\par
with the product taken over the real places of \(K\). We put\par
\par
\DFormula{(3.11b)\,\,\,\,\Lambda (s,\varepsilon )=\gamma (s)L(s,\varepsilon ).}
\DAnchor{23}{249}
(3.12) Theorem. We have\par
\par
\DFormula{\Lambda (s,\varepsilon )=i^{\sum a_{v}}\cdot \Lambda (1-s,T\varepsilon ).}
\par
In fact, the case of primary interest for us is that where all \(a_{v}\) are either 0 \((\varepsilon \) even) or 1 \((\varepsilon \) odd) and \(s\) is an integer \(k\geq 1\) with the same parity as \(\varepsilon \). In this case we will obtain the following\par
\par
(3.13) Corollary. If \(\varepsilon \) has parity \((-1)^{k}\) and \(k\geq 1\), then\par
\par
\DFormula{L(k,\varepsilon )=2^{-r}\alpha _{k}\,L(1-k,T\varepsilon ),}
\par
where\par
\par
\DFormula{\alpha _{k}=\lbrace ((2\pi i)^{k})/(k-1)!\rbrace ^{r}.}
\par
Before beginning the proof of (3.12), (3.13), we remark that these equations are essentially those discussed by Siegel [38].\par
\par
Proof of (3.12). We will in fact establish an identity between the partial L-function obtained by summing \(\varepsilon (𝔞)\mathcal{N}𝔞^{-s}\) over a given wide ideal class of \(K\) and that made by summing \((T\varepsilon )(𝔞)\mathcal{N}𝔞^{s-1}\) over the inverse of the class. Let \(𝔲\) be an ideal. We will prove that\par
\par
\DFormula{(3.14)\,\,\,\,\sum _{\beta \in K^{\ast}/U}\,\varepsilon (\beta 𝔲)/\vert \beta \vert ^{s}}
\DFormula{=h(s)\,\sum _{\beta \in K^{\ast}/U}\,\mathcal{N}𝔲\cdot T\varepsilon (\beta 𝔲^{-1})\vert \beta \vert ^{s-1},}
\par
where the notation is as follows:\par
\par
\(U=unit\) group of \(K\)\par
\par
\DFormula{h(s)=i^{\sum a_{v}}\cdot \gamma _{\varepsilon }(1-s)\gamma _{\varepsilon }(s)^{-1}}
\par
\(\vert \beta \vert \) is the archimedean absolute value of \(\beta \), i.e., \(\vert \beta \vert =\vert \mathcal{N}\beta \vert \).\par
\par
Let \(\rho \) be the function \(x\mapsto \prod _{v}(\operatorname{sgn}\) \(v(\beta ))^{a_{v}}\). We have, by the parity hypothesis,\par
\par
\DFormula{(3.15)\,\,\,\,\varepsilon (\beta 𝔲)=\varepsilon (𝔲\cdot \beta )\rho (\beta )=\varepsilon _{𝔲}(\beta )\rho (\beta );}
\par
\DFormula{(T\varepsilon )(\beta 𝔲^{-1})=(T\varepsilon )(𝔲^{-1}\cdot \beta )\rho (\beta )}
\DFormula{=(\mathcal{N}𝔲)^{-1}\widehat{\varepsilon }_{𝔲}(\beta )\rho (\beta ).}
\par
Thus the formula to be verified reads:\par
\par
\DFormula{(3.16)\,\,\,\,\sum _{\beta \in K^{\ast}/U^{+}}\,[\varepsilon _{𝔲}(\beta )\rho (\beta )]/\vert \beta \vert ^{s}}
\DFormula{?=\,h(s)\,\sum _{\beta \in K^{\ast}/U^{+}}\,[\widehat{\varepsilon }_{𝔲}(\beta )/\rho (\beta )]\cdot \vert \beta \vert ^{s-1},}
\par
where \(U^{+}\) is the group of totally positive units. (We prefer to sum over this latter unit group because \(\varepsilon _{𝔲}\), \(\widehat{\varepsilon }_{𝔲}\) and \(\rho \) (as well as \(\vert \) \(\vert )\) are invariant by it, whereas only the products \(\varepsilon _{𝔲}\cdot \rho \), \(\widehat{\varepsilon }_{𝔲}\rho \) need be invariant by \(U\). In (3.14), we have the right to replace \(U\) by any subgroup of finite index in it.) Let \(\chi \) be a quasicharacter\par
\par
\DFormula{K_{\infty }^{\ast}\to \mathbb{C}^{\ast},}
\par
and let \(\varphi _{\infty }\) be a function such that both it and its Fourier transform are rapidly decreasing at infinity (cf. [39], §2.4). For Re \(s\gg 0\), the integral\par
\par
\DFormula{(3.17)\,\,\,\,\int _{K_{\infty }^{\ast}}\,\chi (x)\vert x\vert ^{s}\,\varphi _{\infty }(x)d^{\times }x}
\DAnchor{24}{250}
converges to an analytic function of \(s\) which has a meromorphic continuation to all of \(\mathbb{C}\). Furthermore, there is a non-zero meromorphic function \(\alpha (\chi ,s)\) of \(s\), independent of \(\varphi _{\infty }\), such that\par
\par
\DFormula{(3.18)\,\,\,\,\alpha (\chi ,s)\int _{K_{\infty }^{\ast}}\,\chi ^{-1}(x)\vert x\vert ^{s}\,\varphi _{\infty }(x)d^{\times }x}
\DFormula{=\int _{K_{\infty }^{\ast}}\,\chi (x)\vert x\vert ^{1-s}\,\widehat{\varphi }_{\infty }(x)d^{\times }x.}
\par
This formula, which is just a semi-local variant of the functional equation of [39, loc. cit.], expresses the fact that the Fourier transform of the distribution \(\chi \vert \) \(\vert ^{-s}=\chi _{s}\) on \(K_{\infty }\) is proportional to \(\chi _{s}^{-1}\vert \) \(\vert ^{-1}\).\par
\par
Let us take \(\chi \) equal to the “sign” character \(\rho \), viewed as a function on \(K_{\infty }^{\ast}\). If we write the collection of numbers \((a_{v})\) as a tuple \((a_{1},...,a_{r})\), using the chosen numbering of the real places, then\par
\par
\DFormula{\rho (x)=\prod _{j=1}^{r}\,\operatorname{sgn}(x_{j})^{a_{j}}.}
\par
Following Tate, we compute \(\alpha (\rho ,s)\) by choosing\par
\par
\DFormula{\varphi _{\infty }(x)=\prod _{j}(x_{j}^{a_{j}}\cdot e^{-x_{j}^{2}}).}
\par
With this choice for \(\varphi _{\infty }\), each integral in (3.18) may be rewritten as a product of \(r\) different elementary integrals of the type computed on \(p\). 317 of [39]. We find that\par
\par
\DFormula{\alpha (\rho ,s)=h(s),}
\par
with h as in (3.14). This (3.14) is a special case of\par
\par
(3.19) Theorem. Let \(\chi :K_{\infty }^{\ast}\to \mathbb{C}^{\ast}\) be a quasicharacter invariant under a subgroup \(U_{0}\) of finite index in \(U\). Let \(\varepsilon \) be a locally constant compactly supported function on \(\widehat{K}\) which is also invariant under \(U_{0}\). Then we have the equality of meromorphic functions\par
\par
\DFormula{\sum _{\beta \in K^{\ast}/U_{0}}\,\varepsilon (\beta )\chi (\beta )\vert \beta \vert ^{-s}}
\DFormula{=\alpha (\chi ,s)\sum _{\beta \in K^{\ast}/U_{0}}\,\widehat{\varepsilon }(\beta )\chi (\beta )^{-1}\vert \beta \vert ^{s-1}.}
\par
Remarks. 1. More generally, let \(K\) be a global field, and let \(S\) be a finite set of primes of \(K\) which contains the archimedean primes. Let \(\chi \) be a quasicharacter of \(\prod _{v\in S}K_{v}^{\ast}\), invariant by a subgroup \(U_{0}\) of finite index in the group of S-units of \(K\). The Fourier transform \(\widehat{\chi }\) of \(\chi \) is a multiple \(\alpha (\chi )\cdot \vert \) \(\vert ^{-1}\cdot \chi ^{-1}\) of \(\vert \) \(\vert ^{-1}\chi ^{-1}\). Let \(\varepsilon \) be a \(U_{0}\)-invariant locally constant function with compact support on \(\prod _{v\notin S}K_{v}\), and let \(\widehat{\varepsilon }\) be its Fourier transform. The proof given below yields the identity\par
\par
\DFormula{(3.20)\,\,\,\,\sum _{x\in K^{\ast}/U_{0}}\,\chi (x)\varepsilon (x)=\sum _{x\in K^{\ast}/U_{0}}\,\widehat{\chi }(x)\widehat{\varepsilon }(x),}
\par
where both sides are defined by analytic continuation in the family \(\chi \vert \) \(\vert ^{-s}\). Suppose that we (formally) allow \(S\) to be the set of all places of \(K\), and take \(U_{0}=K^{\ast}\). Then (3.20) yields the identity \(\widehat{\chi }=\chi ^{-1}\vert \) \(\vert ^{-1}\) (i.e., \(\alpha (\chi )=1)\) for \(\chi \) a grossencharacter. This identity is Tate’s global functional equation, as rewritten by Weil [40].\par
\par
2. Our method of proof is similar to that given in Sato-Shintani [31]. \(A\) common generalization should be possible.\par
\DAnchor{25}{251}
Proof of (3.19). For simplicity, we will write \(U\) instead of \(U_{0}\) for the unit group appearing in (3.19). Furthermore, let us continue to write \(\chi _{s}=\chi \cdot \vert \) \(\vert ^{-s}\).\par
\par
Step \(I\). Let \(\varphi _{\infty }\) be a function on \(K_{\infty }\) as in (3.17), chosen so that\par
\par
\DFormula{\int _{K_{\infty }^{\ast}}\,\chi _{s}(x)\varphi _{\infty }(x)d^{\times }x}
\par
is not identically zero. Let \(\varphi \) be the function \(\varepsilon \otimes \varphi _{\infty }\) on \(𝔸\):\par
\par
\DFormula{\varphi (x)=\varphi _{\infty }(x_{\infty })\varepsilon (x_{f}).}
\par
For each \(x\in K_{\infty }^{\ast}\hookrightarrow 𝔸\) we consider the function\par
\par
\(a\mapsto \varphi (xa)\)    \((a\in 𝔸)\),\par
\par
whose Fourier transform is\par
\par
\DFormula{a\mapsto \vert x\vert ^{-1}\,\widehat{\varphi }(x^{-1}a).}
\par
We apply the Poisson summation formula to this function and obtain\par
\par
\DFormula{(3.21)\,\,\,\,\sum _{\beta \in K}\,\varphi (x\beta )=\sum _{\beta \in K}\vert x\vert ^{-1}\,\widehat{\varphi }(x^{-1}\beta ),}
\par
cf. [39, p.333].\par
\par
Step II. We have an equality of meromorphic functions\par
\par
\DFormula{(3.22)\,\,\,\,\int _{K_{\infty }^{\ast}/U}\,\chi _{s}(x)(\sum _{\beta \in K^{\ast}}\varphi (x\beta ))d^{\times }x}
\DFormula{=\int _{K_{\infty }^{\ast}/U}\,\chi _{s}(x^{-1})\vert x\vert (\sum _{\beta \in K^{\ast}}\widehat{\varphi }(x\beta ))d^{\times }x.}
\par
[Note that the integrands are \(U\)-invariant. E.g., if \(\gamma \in U\) and \(x'=\gamma x\), then \(\varphi (x'\beta )\) is by definition \(\varphi (x\gamma _{\infty }\beta )\). Since \(\varepsilon \) is invariant under \(U\), we may rewrite this as \(\varphi (x\cdot \beta \gamma ).]\)\par
\par
Let \(K_{\infty }^{>}\) (resp. \(K_{\infty }^{<})\) be the subset of \(K_{\infty }^{\ast}\) consisting of elements with \(\vert x\vert >1\) (resp. \(\vert x\vert <1)\). Over \(K_{\infty }^{>}\), both integrals converge to analytic functions of \(s\) because of the rapid decrease of \(\varphi ,\widehat{\varphi }\) at \(\infty \). Over \(K_{\infty }^{<}\), the left-hand integral is well behaved for Re \(s\ll 0\) and the right hand integral for Re \(s\gg 0\). Let \(L\) and \(R\), respectively, denote the integrands of these integrals.\par
\par
Suppose that we multiply (3.21) by \(\chi _{s}(x)\), integrate over \(K_{\infty }^{>}/U\), change variables \(x\mapsto x^{-1}\) on the right-hand side, and isolate the terms corresponding to \(\beta =0\). We obtain\par
\par
\DFormula{\int _{K_{\infty }^{>}/U}\,L}
\DFormula{=\int _{K_{\infty }^{<}/U}\,R+\widehat{\varphi }(0)\int _{K_{\infty }^{<}/U}\vert x\vert \chi _{s}(x^{-1})d^{\times }x-\varphi (0)\int _{K_{\infty }^{>}/U}\chi _{s}(s)d^{\times }x,}
\par
for Re \(s\gg 0\). Similarly, for Re \(s\ll 0\) we have\par
\par
\DFormula{\int _{K_{\infty }^{>}/U}\,R}
\DFormula{=\int _{K_{\infty }^{<}/U}\,L+\varphi (0)\int _{K_{\infty }^{<}/U}\chi _{s}(x)d^{\times }x-\widehat{\varphi }(0)\int _{K_{\infty }^{>}/U}\chi _{s}(x^{-1})\vert x\vert d^{\times }x.}
\par
Now the point is that each integral multiplying \(\varphi (0)\) and \(\widehat{\varphi }(0)\) is meromorphic on \(\mathbb{C}\), so that \(\int _{K_{\infty }^{<}/U}L\), \(\int _{K_{\infty }^{<}/U}R\) are continuable to meromorphic functions on \(\mathbb{C}\).\par
\DAnchor{26}{252}
Furthermore, when we take the difference of the two equations just written, the \(\varphi (0)\) and the \(\widehat{\varphi }(0)\) terms cancel, giving the desired equality (3.22).\par
\par
For example, let us show that \(\int _{K_{\infty }^{<}/U}\chi _{s}(x)d^{\times }x\) and \(\int _{K_{\infty }^{>}/U}\chi _{s}(x)d^{\times }x\) are meromorphic in \(s\), and that each is the negative of the other.\par
\par
We let \(\pi :K_{\infty }^{\ast}/U\to (0,\infty )\) be the surjection \(x\mapsto \vert x\vert \). By Dirichlet’s theorem, the fibres of \(\pi \) are compact. Hence for each \(t\in (0,\infty )\), the integral\par
\par
\DFormula{\int _{\pi ^{-1}(t)}\,\chi _{s}(x)d^{\times }x}
\par
is defined for all \(s\). Because \(\pi \) and \(\chi _{s}\) are homomorphisms, the integral vanishes if \(\chi \) (and hence \(\chi _{s})\) is non-trivial on the kernel of \(\pi \). In that case, by Fubini’s theorem, the two integrals to be calculated are identically zero. So we can assume that \(\chi \) is trivial on Ker \(\pi \), giving that \(\chi (x)=\vert x\vert ^{s_{0}}\) for some \(s_{0}\in \mathbb{C}\). We then find:\par
\par
\DFormula{\int _{K_{\infty }^{>}/U}\,\chi _{s}(x)d^{\times }x=V/(s-s_{0})}
\par
and\par
\par
\DFormula{\int _{K_{\infty }^{<}/U}\,\chi _{s}(x)d^{\times }x=V/(s_{0}-s),}
\par
where \(V\) is the volume of \(\pi ^{-1}(1)\). Hence our claim, and thus (3.22), is verified.\par
\par
Step III. We apply Fubini’s theorem to the left hand side of (3.22), obtaining\par
\par
\DFormula{\sum _{\beta \in K^{\ast}/U}\,\int _{K_{\infty }^{\ast}}\,\chi _{s}(x)\varphi (x\beta )d^{\times }x.}
\par
We replace \(x\) by \(\beta ^{-1}x\), i.e., \(\beta _{\infty }^{-1}x\), in the integral. It becomes\par
\par
\DFormula{\varepsilon (\beta )\chi _{s}(\beta )^{-1}\,\int _{K_{\infty }^{\ast}}\,\chi _{s}(x)\varphi _{\infty }(x)d^{\times }x,}
\par
so that the left-hand side of (3.22) may be rewritten\par
\par
\DFormula{[\sum _{\beta \in K^{\ast}/U}\,\varepsilon (\beta )\chi _{s}(\beta )^{-1}]\cdot \int _{K_{\infty }^{\ast}}\chi _{s}(x)\varphi _{\infty }(x)d^{\times }x.}
\par
This decomposition gives a meromorphic continuation to the first factor, i.e., the sum: the product of the two factors is meromorphic by the above, whereas the integral is meromorphic (and non-zero by the choice of \(\varphi _{\infty })\). We similarly write the right-hand side of (3.22) as\par
\par
\DFormula{[\sum _{\beta \in K^{\ast}/U}\,\widehat{\varepsilon }(\beta )\chi _{s}(\beta )\vert \beta \vert ^{-1}]\cdot \int _{K_{\infty }^{\ast}}\chi (x^{-1})\vert x\vert ^{1+s}\widehat{\varphi }_{\infty }(x)d^{\times }x.}
\par
This gives, because of (3.18), the equality\par
\par
\DFormula{\sum _{\beta \in K^{\ast}/U}\,\varepsilon (\beta )\chi ^{-1}(\beta )\vert \beta \vert ^{s}}
\DFormula{=\alpha (\chi ^{-1},-s)\sum _{\beta \in K^{\ast}/U}\,\widehat{\varepsilon }(\beta )\chi (\beta )\vert \beta \vert ^{-1-s},}
\par
which is just (3.19) with \(\chi ^{-1}\) and \(-s\) instead of \(\chi \) and \(s\).\par
\DAnchor{27}{253}
Proof of (3.13). The factor \(\alpha (\chi ,s)\) is in this case\par
\par
\DFormula{h(s)=\begin{cases}\left[\dfrac{\Gamma_{\mathbb R}(1-s)}{\Gamma_{\mathbb R}(s)}\right]^r&\text{if }k\text{ is even},\\\left[\dfrac{i\Gamma_{\mathbb R}(2-s)}{\Gamma_{\mathbb R}(1+s)}\right]^r&\text{if }k\text{ is odd}.\end{cases}}
\par
However, we have\par
\par
\DFormula{\Gamma _{\mathbb{R}}(s)/\Gamma _{\mathbb{R}}(1-s)=2^{1-s}\pi ^{-s}\operatorname{cos}(\pi s/2)\Gamma (s),}
\par
\DFormula{\Gamma _{\mathbb{R}}(1+s)/(i\Gamma _{\mathbb{R}}(2-s))=-i2^{1-s}\pi ^{-s}\operatorname{sin}(\pi s/2)\Gamma (s),}
\par
as remarked in [39, p.317]. The assertion now follows on setting \(s=k\) in each case.\par
\par
Trivial Zeros\par
\par
(3.23) Proposition. The function \(\Lambda (s,\varepsilon )\) of (3.11b) is holomorphic on \(\mathbb{C}\) except in the case where \(\varepsilon \) is even, i.e., where all \(a_{v}\) vanish. In this case, \(\Lambda (s,\varepsilon )\) is holomorphic for all \(s\neq 0,1\) and has at worst simple poles at these points.\par
\par
As is well known, by considering the poles of \(\Gamma (s)\) one deduces the existence of “trivial zeros” for \(L(s,\varepsilon )\):\par
\par
(3.24) Corollary. We have \(L(1-k,\varepsilon )=0\) for all even integers \(k\geq 1\) except when \(\varepsilon \) is an even function. Similarly, if \(\varepsilon \) is not an odd function, we have \(L(1-k,\varepsilon )=0\) for all odd integers \(k\geq 1\).\par
\par
[The assertion concerning \(k=1\) would be false in the excluded case \(K=\mathbb{Q}\). Note that in (3.24) the function \(\varepsilon \) is still assumed to have some parity, as in (3.12).]\par
\par
Proof of (3.23). We review the proof of (3.19), making the normalization \(\chi =\rho \). The proof of the equality (3.22) shows that the functions in (3.22) are everywhere holomorphic, except for the case where \(\chi \) is trivial, in which case there may be simple poles only at the points \(s=0,1\). We now make in Step III of the proof of (3.19) the choice\par
\par
\DFormula{\varphi _{\infty }(x)=\prod _{j}(x_{j}^{a_{j}}\cdot e^{-x_{j}^{2}}).}
\par
Then we may explicitly calculate the integrals appearing in Step III, obtaining in particular the equation\par
\par
\DFormula{\int _{K_{\infty }^{\ast}}\,\chi _{s}(x)\varphi _{\infty }(x)d^{\times }x=\prod _{v}\,\Gamma _{\mathbb{R}}(a_{v}-s).}
\par
Except for the change of sign in \(s\), the product is the \(\Gamma -factor\) \(\gamma (s)\) used to define \(\Lambda (s,\varepsilon )\). We may conclude that the product of \(\gamma (s)\) by the L-series in the statement\par
\DAnchor{28}{254}
of (3.19) has the property asserted of \(\Lambda (s,\varepsilon )\): it is holomorphic if the \(a_{v}\) are not all zero, and otherwise has possible simple poles at 0 and 1 and no other poles. Expressing \(\Lambda (s,\varepsilon )\) in terms of such series, we find that \(\Lambda (s,\varepsilon )\) has this property as well.\par
\par
\subsection*{§4. An Irreducibility Theorem}
\par
This § concerns the “Hilbert-Blumenthal” moduli problem of classifying abelian varieties with real multiplication. We wish to complement the work of Rapoport [28] with an irreducibility theorem in characteristic \(p\). We recall that a Hilbert-Blumenthal abelian variety (or: HBAV) relative to the integer ring \(\mathcal{O}\) of \(K\) over a base \(S\) is an abelian scheme X/S, furnished with a homomorphism \(m:\mathcal{O}\hookrightarrow \operatorname{End}\) \(X\) making \(\operatorname{Lie}(X/S)\) into a locally free \(\mathcal{O}\otimes \mathcal{O}_{S}\)-module of rank 1. It follows in particular from this definition that the relative dimension of \(X\) over \(S\) is the degree \(r\) of \(K\) over \(\mathbb{Q}\). Conversely, if \(S\) is of characteristic 0 and X/S is an abelian scheme of relative dimension \(r\) furnished with an \(m:\mathcal{O}\hookrightarrow \operatorname{End}\) \(X\), then the condition on \(\operatorname{Lie}(X/S)\) is automatically satisfied. Also, if \(S\) is of characteristic \(p>0\), then the condition is satisfied at least whenever \(X\) is ordinary.\par
\par
Let \(X\) be a HBAV over an algebraically closed field \(k\). If \(N\) is prime to the characteristic of \(k\), the group scheme \(X_{N}\) of \(N\)-division points of \(X\) is etale, and so may be identified with the group of its points: a free \(\mathcal{O}/N\mathcal{O}\)-module of rank 2. Also, we attach to \(X\) its polarization module \(\mathcal{P}(X)\); this is the invertible \(\mathcal{O}\)-module consisting of the symmetric \(\mathcal{O}\)-linear homomorphisms \(X\to X^{\ast}\), where \(X^{\ast}\) is the dual of \(X\). This module has a natural positivity, in which we declare positive those homomorphisms defined by a polarization [28, §1].\par
\par
In the special case \(k=\mathbb{C}\), we may interpret \(X_{N}\) and \(\mathcal{P}(X)\) in terms of the homology group \(T(X)=H_{1}(X(\mathbb{C}),\mathbb{Z})\). For \(X_{N}\), we have\par
\par
\DFormula{(4.1)\quad X_N=(1/N)T(X)/T(X)\xrightarrow[\sim]{N}T(X)/NT(X).}
\par
For \(\mathcal{P}(X)\), we regard T(X) and \(T(X^{\ast})\) as being paired over \(\mathbb{Z}\) into \(\mathbb{Z}(1)=2\pi i\mathbb{Z}\). We have\par
\par
\DFormula{T(X^{\ast})=\operatorname{Hom}_{\mathcal{O}}(T(X),𝔇^{-1}(1)),}
\par
which leads to\par
\par
\DFormula{(4.2)\,\,\,\,\mathcal{P}(X)=\operatorname{Hom}_{\mathcal{O}}(\wedge ^{2}T(X),𝔇^{-1}(1)).}
\par
(The exterior power is taken over \(\mathcal{O}.)\) The notion of positivity on \(\mathcal{P}(X)\) results from the fact that the real vector space \(T(X)\otimes _{\mathbb{Z}}\) \(\mathbb{R}\) has a natural complex structure [28, 1.26].\par
\par
If we combine (4.1) and (4.2), we obtain an isomorphism\par
\par
(4.3)    \(\wedge ^{2}X_{N}\) \(\xrightarrow{\sim}\) \(\operatorname{Hom}_{\mathcal{O}}(\mathcal{P}(X),𝔇^{-1}(1))\otimes \mathbb{Z}/N\mathbb{Z}\).\par
\DAnchor{29}{255}
Over an arbitrary base \(S\), the groups \((X\overline{_{s}})_{N}\) (with \(N\) invertible on \(S)\) and \(\mathcal{P}(X\overline{_{s}})\) fit together to give local systems \(X_{N}\), \(\mathcal{P}(X)\) (for the etale topology). From the \(e_{m}-pairings\) of Weil, we may construct an analogue of (4.3):\par
\par
\DFormula{(4.4)\,\,\,\,\wedge ^{2}X_{N}\,\xrightarrow{\sim}\,\operatorname{Hom}_{\mathcal{O}}(\mathcal{P}(X),𝔇^{-1})\otimes \mu _{N}.}
\par
Let \(N\geq 3\) be an integer, and let \(\mathcal{P}\) be an invertible \(\mathcal{O}\)-module of rank 1, given with a positivity. For each base \(S\), we consider triples \((X,\lambda ,\alpha )\) over \(S\) consisting of: a HBAV X/S, a positive isomorphism \(\lambda :\mathcal{P}(X)\xrightarrow{\sim}\mathcal{P}\), and a level-N structure \(\alpha :(\mathcal{O}/N\mathcal{O})^{2}\xrightarrow{\sim}X_{N}\). Let \(F[\mathcal{P},N]\) be the functor\par
\par
\(S\mapsto \lbrace isomorphism\,classes\,of\,triples\,(X,\lambda ,\alpha )\,over\,S\rbrace \).\par
\par
According to [28], \(F[\mathcal{P},N]\) is represented by an algebraic space \(\mathcal{M}=\mathcal{M}_{N}[\mathcal{P}]\), smooth over Z[1/N]. One can find a compactification \(\overline{\mathcal{M}}\) of \(\mathcal{M}\), proper and smooth over Z[1/N], so that the complement of \(\mathcal{M}\) in \(\overline{\mathcal{M}}\) is a relative divisor with normal crossings. (We note in passing that the level-N structure is imposed in order to eliminate automorphisms.) It follows that \(\mathcal{M}\) has the “same” geometric connected components in characteristic \(p\nmid N\) as over \(\mathbb{C}\). These components are parameterized by the set of “invertible” elements of the free \(\mathcal{O}/N\mathcal{O}\)-module\par
\par
\(H=(\operatorname{Hom}_{\mathcal{O}}(\mathcal{P},𝔇^{-1})\otimes \mathbb{Z}/N\mathbb{Z})(1)\).\par
\par
Explicitly, given a triple \((X,\lambda ,\alpha )\), we obtain an element of \(H\) according to the following recipe. From \(\alpha \) and (4.4) we obtain an isomorphism\par
\par
\DFormula{\mathcal{O}/N\mathcal{O}=\wedge ^{2}(\mathcal{O}/N\mathcal{O})^{2}\,\xrightarrow{\sim}\,\wedge ^{2}X_{N}\approx \operatorname{Hom}_{\mathcal{O}}(\mathcal{P}(X),𝔇^{-1})\otimes \mu _{N}=H',}
\par
and using \(\lambda \) we have also an isomorphism \(H\approx H'\). Composing the two, we obtain an isomorphism \(\mathcal{O}/N\mathcal{O}\approx H\), and hence a canonically given “basis vector” of \(H\).\par
\par
Let \(p\nmid N\) be a prime, and let \(p^{n}>1\) be a power of \(p\). Let \(\mathcal{M}^{o}=\mathcal{M}_{N}^{o}[\mathcal{P},p]\) be the open subset of the reduction of \(\mathcal{M}\) modulo \(p\) which corresponds to ordinary abelian varieties. (We recall that \(X\) is ordinary if \(X_{p}\) is an extension of an etale group scheme by a group scheme of multiplicative type.) We shall see below that \(\mathcal{M}^{o}\) meets each connected component of \(\mathcal{M}\) mod \(p\), so that it is dense in \(\mathcal{M}\) mod \(p\).\par
\par
If X/S is an ordinary Hilbert abelian variety in characteristic \(p\), then \(X_{p^{n}}\) is locally (over \(S\) for the etale topology) an extension of an etale group scheme, \(\mathcal{O}-isomorphic\) to \(\mathcal{O}/p^{n}\mathcal{O}\), by a group scheme of multiplicative type, isomorphic to \(\mathcal{O}\otimes \mu _{p^{n}}\):\par
\par
\DFormula{(X_{p^{n}})^{o}\sim \mathcal{O}\otimes \mu _{p^{n}}.}
\par
Let \(\mathcal{J}=\mathcal{J}[\mathcal{P},N,p^{n}]\) be the covering of \(\mathcal{M}^{o}\) defined by:\par
\par
\DFormula{\mathcal{J}=\operatorname{Isom}_{\mathcal{M}^{o}}(\mu _{p^{n}}\otimes \mathcal{O},(X_{p^{n}})^{o}),}
\DAnchor{30}{256}
\(X/\mathcal{M}^{o}\) being the universal Hilbert-Blumenthal abelian variety. It is a principal homogeneous space (“torsor”) over \(\mathcal{M}^{o}\), with structural group \((\mathcal{O}/p^{n}\mathcal{O})^{\ast}\). It may alternately be obtained as the reduction (mod p) of the space \(\mathcal{J}'\) associated to the moduli problem of classifying quadruples \((X,\lambda ,\alpha ,\beta )\), with \(X\), \(\lambda \), \(\alpha \) as above, and \(\beta \) an \((\mathcal{O}\)-linear) immersion\par
\par
\DFormula{\mathcal{O}\otimes \mu _{p^{n}}\hookrightarrow X.}
\par
The space \(\mathcal{J}'\) is etale and quasi-finite over \(\mathcal{M}\), but not finite.\par
\par
(4.5) Theorem. The covering \(\mathcal{J}[\mathcal{P},N,p^{n}]\) of \(\mathcal{M}_{N}^{o}[\mathcal{P},p]\) is geometrically irreducible. I.e., if \(\overline{\mathbb{F}}\) is an algebraic closure of \(\mathbb{F}_{p}\), then \(\mathcal{J}\) defines an irreducible covering of each connected component of \(\mathcal{M}^{o}\otimes _{\mathbb{F}_{p}}\overline{\mathbb{F}}\).\par
\par
Proof. Let \(q\) be large enough so that all the geometric connected components of \(\mathcal{M}^{o}\) are defined over \(\mathbb{F}_{q}\). (By the above discussion, it suffices in fact to have \(\mathbb{F}_{q}\) containing the Nth roots of 1.) Let \(\mathcal{M}_{1}^{o}\subset \mathcal{M}^{o}\otimes _{\mathbb{F}_{p}}\mathbb{F}_{q}\) be one of the components, and let \(m\in \mathcal{M}_{1}^{o}(\overline{\mathbb{F}})\) be a base point of \(\mathcal{M}_{1}\). The covering of \(\mathcal{M}_{1}^{o}\) induced by \(\mathcal{J}\) is determined (up to unique isomorphism) by its fibre \(\mathcal{J}_{m}\) and the monodromy action of \(\pi _{1}(\mathcal{M}_{1}^{o},m)\) on \(\mathcal{J}_{m}\). This action is given by a character\par
\par
\DFormula{\rho :\pi _{1}(\mathcal{M}_{1}^{o},m)\to (\mathcal{O}/p^{n}\mathcal{O})^{\ast}.}
\par
Irreducibility of the covering, after extension from \(\mathbb{F}_{q}\) to \(\mathbb{F}_{q^{a}}\), amounts to the surjectivity of \(\rho \vert \pi _{1}(\mathcal{M}_{1}^{o}\otimes _{\mathbb{F}_{q}}\mathbb{F}_{q^{a}})\). Geometric irreducibility amounts to surjectivity for all \(a\geq 1\): any geometric connected component of \(\mathcal{J}\) is defined over some \(\mathbb{F}_{q^{a}}\).\par
\par
Each closed point \(x\) of \(\mathcal{M}_{1}^{o}\otimes _{\mathbb{F}_{q}}\mathbb{F}_{q^{a}}\) defines a Frobenius element \(F_{x}\) in \(\pi _{1}(\mathcal{M}_{1}^{o}\otimes \mathbb{F}_{q^{a}})\). (Although \(F_{x}\) is in fact not well defined, its image under \(\rho \) is well defined.) We will prove the required surjectivity by using these Frobenius elements and their powers. (The Čebotarev density theorem tells us that we will be able to do this if the theorem is true.) The powers come in because of the following “base change” formula. Let b be a multiple of a, and let \(x\in \mathcal{M}_{1}^{o}(\mathbb{F}_{q^{b}})\). Then \(x\) is a closed point of \(\mathcal{M}_{1}^{o}\otimes _{\mathbb{F}_{q}}\mathbb{F}_{q^{b}}\) but defines as well as closed point \(x'\) of \(\mathcal{M}_{1}^{o}\otimes _{\mathbb{F}_{q}}\mathbb{F}_{q^{a}}\). Let \(\mathbb{F}_{q^{c}}\) be the residue field of \(x'\), so that \(a\vert c\vert b\). Then we have \(F_{x}=F_{x'}^{b/c}\) in \(\pi _{1}(\mathcal{M}_{1}^{o}\otimes _{\mathbb{F}_{q}}\mathbb{F}_{q^{a}})\).\par
\par
Suppose that \(x\in \mathcal{M}_{1}^{o}(\mathbb{F}_{q^{b}})\) corresponds to a triple \((X,\lambda ,\alpha )\) over \(\mathbb{F}_{q^{b}}\). Then \(\rho (F_{x})\in (\mathcal{O}/p^{n}\mathcal{O})^{\ast}\) is the number giving the action of Frobenius on the etale group scheme \(\operatorname{Hom}(\mathcal{O}\otimes \mu _{p^{n}},(X_{p^{n}})^{o})\). Equivalently, we write the Frobenius endomorphism of the formal group \(\widehat{X}\) of \(X\) as a product \(q^{b}\cdot u\), where \(u\in (\mathcal{O}\otimes \mathbb{Z}_{p})^{\ast}\). Then \(\rho (F_{x})\) is the image of u in \((\mathcal{O}/p^{n}\mathcal{O})^{\ast}\) cf. [15, (4.2.1)]. (Depending on conventions chosen, one may prefer to take \(u^{-1}\) instead of u at this point; this problem is irrelevant for our purposes.)\par
\par
Notice that the recipe just given for \(\rho (F_{x})\) does not involve either the level \(N\)-structure \(\alpha \) or the isomorphism \(\lambda \). On the other hand, the action of \(\operatorname{GL}(2,\mathcal{O}/N\mathcal{O})\) on \(\mathcal{M}^{o}\) deduced from\par
\par
\DFormula{g:(X,\lambda ,\alpha )\mapsto (X,\lambda ,\alpha \circ g^{-1})}
\par
permutes the geometric connected components of \(\mathcal{M}^{o}\). This enables us to ignore the distinction between \(\mathcal{M}_{1}^{o}(\mathbb{F}_{q^{b}})\) and \(\mathcal{M}^{o}(\mathbb{F}_{q^{b}})\) in finding a triple \((X,\lambda ,\alpha )\) such that the associated u maps mod \(p^{n}\) to a desired element \(\overline{u}\) of \((\mathcal{O}/p^{n}\mathcal{O})^{\ast}\); we are\par
\DAnchor{31}{257}
required only to find an ordinary \(X/\mathbb{F}_{q^{b}}\), for some multiple b of a, with the following properties:\par
\par
(i) \(X\) admits a level-N structure \(\alpha \) over \(\mathbb{F}_{q^{b}}\);\par
(ii) The number \(u\in (\mathcal{O}\otimes \mathbb{Z}_{p})^{\ast}\) computed for \(X\) maps to the given \(\overline{u}\) in \((\mathcal{O}/p^{n}\mathcal{O})^{\ast}\);\par
(iii) The polarization module \(\mathcal{P}(X)\) is isomorphic to the given \(\mathcal{P}\) (as an invertible \(\mathcal{O}\)-module with positivity).\par
\par
Carrying out this construction will prove the surjectivity, and thus the theorem.\par
\par
We eliminate consideration of requirement (iii) by the following trick. Let \(M\) be a positive integer, prime to pN, such that the ideal \(M\mathcal{O}\) is divisible by some ideal of \(\mathcal{O}\) in each strict ideal class of \(K\). Assume that \(X/\mathbb{F}_{q^{b}}\) satisfies (ii) and admits a level NM structure over \(\mathbb{F}_{q^{b}}\), but does not necessarily satisfy (iii). Choose a level NM structure for \(X\), and view it as a level \(M\) structure plus a level \(N\) structure. If \(𝔠\) is an ideal with \(M\mathcal{O}\subseteq 𝔠\subseteq \mathcal{O}\), we may divide \(X\) by the image in \(X\) of the group \(((\mathcal{O})\times 𝔠\) mod \(M)\) under the immersion \((\mathcal{O}/M\mathcal{O})^{2}\hookrightarrow X\). We obtain a variety \(X'/\mathbb{F}_{q^{b}}\) which satisfies (i) and (ii). It is easily seen (and at least guessed from (4.2)) that the polarization module \(\mathcal{P}(X')\) is isomorphic to \(\mathcal{P}(X)\otimes 𝔠^{-1}\). Since we may choose \(𝔠\) in any strict ideal class of \(K\), we may arrange \(\mathcal{P}(X')\) to be isomorphic to the given \(\mathcal{P}\). We may thus forget about condition (iii), after replacing \(N\) by a suitable multiple.\par
\par
To construct varieties \(X\) satisfying (i) and (ii), we use the description given in [4] of the category of ordinary abelian varieties over a finite field. The equivalence of categories furnished by [4] identifies ordinary Hilbert-Blumenthal abelian varieties with pairs \((\mathcal{L},F)\), where \(\mathcal{L}\) is a locally free \(\mathcal{O}\)-module of rank 2, and \(F\) is a “Frobenius” endomorphism. One requires that \(F\) satisfy:\par
\par
\DFormula{(a)\,\operatorname{det}\,F=q^{b},\,\text{and}\,(\operatorname{tr}\,F)^{2}-4\,\operatorname{det}\,F\ll 0,}
(b) the completion \(\mathcal{L}\otimes \mathbb{Z}_{p}\) of \(\mathcal{L}\) decomposes as the sum \(\mathcal{L}'\oplus \mathcal{L}''\) of two \(F\)-stable free \((\mathcal{O}\otimes \mathbb{Z}_{p})\)-modules of rank 1, with \(F\) acting as \(q^{b}\) u on \(\mathcal{L}'\) and as \(u^{-1}\) on \(\mathcal{L}''\), for some \(u\in (\mathcal{O}\otimes \mathbb{Z}_{p})^{\ast}\).\par
\par
[The number u appearing in (b) is the number u associated above to the abelian variety corresponding to \((\mathcal{L},F).]\) Varieties which admit level-N structures correspond to pairs \((\mathcal{L},F)\) which satisfy in addition:\par
\par
(c) We have \(F\equiv 1\) mod \(N\).\par
\par
We will construct pairs \((\mathcal{L},F)\) by considering irreducible polynomials \(x^{2}+cx+d\) \((c,d\in \mathcal{O})\) and taking \(\mathcal{L}\) to be the integer ring of \(K[x]/(x^{2}+cx+d)\) and \(F\) to be multiplication by \(x\). To satisfy conditions (a) and (b), we need:\par
\par
\DFormula{(1)\,d=q^{b};}
\DFormula{(2)\,c^{2}-4q^{b}\ll 0;}
\DFormula{(3)\,c\in (\mathcal{O}\otimes \mathbb{Z}_{p})^{\ast}.}
\par
For (c), we require that \((x-1)/N\) be integral:\par
\par
\DFormula{(4)\,N\mid (2+c),\,N^{2}\mid (1+q^{b}+c).}
\par
Then \(u^{-1}\), being the unit root of \(x^{2}+cx+d\), will be congruent to \(-c\), mod \(q^{b}\).\par
\DAnchor{32}{258}
To complete the proof of the theorem, we pick \(c\in \mathcal{O}\) so that \(c\equiv -2\) mod \(N^{2}\) and \(c\equiv -\overline{u}^{-1}\) mod \(p^{n}\). Next we choose b so that \(c-4q^{b}\ll 0\), \(q^{b}\equiv 1\) mod \(N\), and \(p^{n}\mid q^{b}\). Conditions (1) through (4) are satisfied, and we have \(u^{-1}\equiv -c\equiv \overline{u}^{-1}\) mod \(p^{n}\) as required.\par
\par
We now define a \(\Gamma _{00}(N)\) structure on a Hilbert-Blumenthal abelian variety to be an \(\mathcal{O}\)-linear immersion \(\mathcal{O}\otimes \mu _{N}\hookrightarrow X\). (Here \(N\) is a positive integer, not necessarily prime to p.) Let \(\mathcal{P}\) be a locally free rank-one \(\mathcal{O}\)-module with positivity, and let \(F[\mathcal{P},N]\) be the functor\par
\par
\DFormula{S\mapsto\left\{\begin{array}{l}\text{isomorphism classes of triples }(X,\lambda,\alpha)\mid X\text{ is a Hilbert-Blumenthal}\\\text{abelian variety over }S;\quad\lambda\text{ is a positive isomorphism }\mathcal P(X)\approx\mathcal P;\quad\alpha\text{ is}\\\text{a }\Gamma_{00}(N)\text{-structure}\end{array}\right\}.}
\par
For \(N\) sufficiently large, such triples have no automorphisms, and \(F[\mathcal{P},N]\) is represented by an algebraic space \(\mathcal{M}[\mathcal{P},N]\), smooth over \(\mathbb{Z}\).\par
\par
(4.6) Corollary. For each prime \(p\), the reduction mod \(p\) of \(\mathcal{M}[\mathcal{P},N]\) is geometrically irreducible.\par
\par
Proof. Let us write \(N=p^{n}N'\), where \(p\nmid N'\). Since \(\mathcal{M}[\mathcal{P},N]\) is surjective mod \(p\) when \(\widetilde{N}\) is a multiple of \(N\), we may prove the corollary with \(N\) replaced by any multiple of \(N\). Hence we may assume that \(N'\) is so large that the moduli space \(\mathcal{M}[\mathcal{P},N']\) exists.\par
\par
Over the algebraic closure \(\overline{\mathbb{F}}\) of \(\mathbb{F}_{p}\), let us fix an embedding \(\mu _{N'}\hookrightarrow (\mathbb{Z}/N'\mathbb{Z})^{2}\). After tensoring with \(\mathcal{O}\), we get a map \(g:\mathcal{O}\otimes \mu _{N'}\hookrightarrow (\mathcal{O}/N'\mathcal{O})^{2}\). Via g, a level \(N'\) structure on a Hilbert-Blumenthal abelian variety gives rise to a \(\Gamma _{00}(N')\)-structure on the variety, so that we have a map \(\mathcal{M}_{N'}[\mathcal{P}]\to \mathcal{M}[\mathcal{P},N']\) over \(\overline{\mathbb{F}}\). If \(\Gamma \) is the subgroup of \(\operatorname{GL}(2,\mathcal{O}/N'\mathcal{O})\) fixing g, we find an isomorphism \(\mathcal{M}_{N'}[\mathcal{P}]/\Gamma \) \(\xrightarrow{\sim}\) \(\mathcal{M}[\mathcal{P},N']\) over \(\overline{\mathbb{F}}\). Now \(\Gamma \) acts transitively on the geometric connected components of \(\mathcal{M}_{N'}[\mathcal{P}]\) because det \(\Gamma =(\mathcal{O}/N'\mathcal{O})^{\ast}\); this gives the irreducibility of \(\mathcal{M}[\mathcal{P},N']\) over \(\overline{\mathbb{F}}\), proving the corollary in case \(N=N'\). If \(p^{n}>1\) we remark that similarly there is an isomorphism\par
\par
\(\mathcal{J}_{N'}[\mathcal{P},p^{n}]/\Gamma \) \(\xrightarrow{\sim}\) \(\mathcal{M}[\mathcal{P},N]\)    (over \(\overline{\mathbb{F}})\).\par
\par
The theorem implies that \(\Gamma \) permutes the geometric connected components of \(\mathcal{J}_{N'}[\mathcal{P},p^{n}]\) (they are the same as those of \(\mathcal{M}_{N'}[\mathcal{P}])\), so we again get the statement of the corollary.\par
\par
\subsection*{§5. Hilbert Modular Forms}
\par
In this §, we review the connections among complex, algebraic, and \(p\)-adic Hilbert modular forms (for \(K)\), emphasizing the \(q\)-expansions of these forms. Our irreducibility theorem (4.6) implies a \(q\)-expansion principle for such forms, which states (roughly speaking) that a form is determined by any one of its \(q\)-expansions.\par
\DAnchor{33}{259}
In our exposition, we will follow very closely Chap. \(I\) of Katz’s paper [19]. In order to make our discussion consistent with that of [19], we shall slightly change the definition of a \(\Gamma _{00}(N)\)-structure on a HBAV X/S. Namely, we shall now understand such a structure to mean an embedding\par
\par
\DFormula{i:(N^{-1}𝔇^{-1}/𝔇^{-1})(1)\hookrightarrow X}
\par
rather than an embedding\par
\par
\DFormula{\alpha :(\mathcal{O}/N\mathcal{O})(1)\hookrightarrow X}
\par
as above. This change has little importance for what follows, in light of the correspondence \(X\mapsto X\otimes 𝔇^{-1}\) between HBAV’s provided with a \(\Gamma _{00}(N)\) structure in the previous sense and those provided with a \(\Gamma _{00}(N)\)-structure in the new sense.\par
\par
In order to define modular forms, we choose (and fix) an ideal \(𝔠\) of \(K\), which will serve as polarization module. We view \(𝔠\) as a signed \(\mathcal{O}\)-module (2.25), endowing it with its canonical positivity \(+_{\mathrm{can}}\). This choice being made, modular forms are functions of quadruples \((X,\lambda ,\omega ,i)\), where:\par
\par
\(X\) is a HBAV over a ring \(R\),\par
\(\lambda :\mathcal{P}(X)\xrightarrow{\sim}𝔠\) is a positive isomorphism,\par
\(\omega \) is a basis of the \(\mathcal{O}\otimes R\)-module \(\underline{\omega }_{X/R}\), a priori a locally free \(\mathcal{O}\otimes R\)-module of rank 1.\par
\(i:N^{-1}𝔇^{-1}/𝔇^{-1}\hookrightarrow X\) is a \(\Gamma _{00}(N)\)-structure (over \(R)\).\par
\par
Note that \(\omega \) amounts to an \(\mathcal{O}\)-linear isomorphism\par
\par
\DFormula{(5.1)\,\,\,\,\operatorname{Lie}(X/R)\xrightarrow{\sim}𝔇^{-1}\otimes R.}
\par
For \(k\) a positive integer, we let \(M_{k}(\Gamma _{00}(N),R)\) be the \(R\)-module of (Hilbert) modular forms of weight \(k\) on \(\Gamma _{00}(N)\) over \(R\), for the polarization module \(𝔠\). Recall that an element of \(M_{k}(\Gamma _{00}(N),R)\) is a function \(F\) defined on quadruples \((X,\lambda ,\omega ,i)\) over \(R\)-algebras \(R'\) such that:\par
\par
For \((X,\lambda ,\omega ,i)\) over \(R'\), the value \(F((X,\lambda ,\omega ,i))\) is an element of \(R'\) which depends only on the \(R'-isomorphism\) class of \((X,\lambda ,\omega ,i)\).\par
\par
For \((X,\lambda ,\omega ,i)\) over \(R'\) and \(\alpha \in (\mathcal{O}\otimes R')^{\ast}\), we have\par
\par
\DFormula{F((X,\lambda ,\alpha \omega ,i))=\mathcal{N}\alpha ^{-k}\cdot F((X,\lambda ,\omega ,i)).}
\par
For \((X,\lambda ,\omega ,i)\) over \(R'\) and \(f:R'\to R''\) a homomorphism of \(R\)-algebras, we have\par
\par
\DFormula{F((X,\lambda ,\omega ,i)/_{R''})=f(F((X,\lambda ,\omega ,i))).}
\par
That is, the function \(F\) commutes with any extension of scalars \(R'\to R''\).\par
\par
An (unramified) cusp on \(\Gamma _{00}(N)\) over \(R\) is determined by the following data:\par
\par
Two ideals \(𝔞,𝔟\in I_{0}\) such that \(𝔞𝔟^{-1}=𝔠\).\par
An \(\mathcal{O}\)-linear isomorphism \(\varepsilon :N^{-1}\mathcal{O}/\mathcal{O}\xrightarrow{\sim}N^{-1}𝔞^{-1}/𝔞^{-1}\).\par
An \(\mathcal{O}\otimes R\)-linear isomorphism\par
\par
\DFormula{j:𝔞^{-1}\otimes R\xrightarrow{\sim}\mathcal{O}\otimes R.}
\DAnchor{34}{260}
Roughly speaking, these three data give rise, respectively to:\par
\par
\(A\) HBAV analogue \(X_{𝔞,𝔟}\) of the Tate curve [15].\par
\(A\) \(\Gamma _{00}(N)\)-structure i on \(X_{𝔞,𝔟}\).\par
\(A\) basis vector \(\omega =\omega (j)\) for \(\underline{\omega }_{X_{𝔞,𝔟}}\),\par
\par
all over a suitable power series ring. Further, \(X_{𝔞,𝔟}\) is provided with a canonical \(\lambda \). By evaluating a modular form \(F\in M_{k}(\Gamma _{00}(N),R)\) on the resulting quadruple, we obtain a formal power series\par
\par
\DFormula{F(𝔞,𝔟,\varepsilon ,j)=\sum _{\mu \in 𝔞𝔟,\,\mu \gg 0\,\text{or}\,\mu =0}\,c_{\mu }q^{\mu },}
\par
whose coefficients \(c_{\mu }=c_{\mu }(𝔞,𝔟,\varepsilon ,j)\) are elements of \(R\). This is the \(q\)-expansion of \(F\) at the cusp (determined by) \(𝔞,𝔟,\varepsilon ,j\).\par
\par
Now suppose that \((𝔞,𝔟,\varepsilon ,j)\) is a cusp, and suppose that \(\alpha \) is an invertible element of \(\mathcal{O}\otimes R\). Let \(\alpha \cdot j\) be the composite\par
\par
\DFormula{\mathfrak a^{-1}\otimes R\overset{j}{\xrightarrow{\sim}}\mathcal O\otimes R\xrightarrow{\sim}\mathcal O\otimes R,}
\par
with the latter map being “multiplication by \(\alpha \).” Then, since \(F\) is of weight \(k\), we have the formula\par
\par
\DFormula{(5.2)\,\,\,\,F(𝔞,𝔟,\varepsilon ,\alpha \cdot j)=\mathcal{N}\alpha ^{-k}F(𝔞,𝔟,\varepsilon ,j).}
\par
(5.3) Example. If \(R\) is a \(\mathbb{Q}\)-algebra, the equality\par
\par
\DFormula{𝔞^{-1}\otimes \mathbb{Q}=\mathcal{O}\otimes \mathbb{Q}}
\par
provides us with a canonical isomorphism \(j_{\mathrm{can}}\) between \(𝔞^{-1}\otimes R\) and \(\mathcal{O}\otimes R\). Suppose, further, that we are given a finite idele \(\alpha \in \widehat{K}^{\ast}\). We may then take \(𝔞=(\alpha )\) to be the ideal generated by \(\alpha \), \(𝔟=𝔞𝔠^{-1}\), and \(\varepsilon \) to be the isomorphism\par
\par
\DFormula{N^{-1}\mathcal{O}/\mathcal{O}\approx N^{-1}\widehat{\mathcal{O}}/\widehat{\mathcal{O}}\,\xrightarrow{\alpha ^{-1}}\,N^{-1}\widehat{𝔞}^{-1}/\widehat{𝔞}^{-1}\approx N^{-1}𝔞^{-1}/𝔞^{-1}.}
\par
We shall refer to the cusp \((𝔞,𝔟,\varepsilon ,j_{\mathrm{can}})\) simply as the cusp determined by \(\alpha \).\par
\par
Our irreducibility theorem (4.6) implies the following “\(q\)-expansion principles” cf. [6, VII, §3], [28, §6]):\par
\par
(5.4) Let \((𝔞,𝔟,\varepsilon ,j)\) be a cusp on \(\Gamma _{00}(N)\) over \(R\). Suppose that \(F(𝔞,𝔟,\varepsilon ,j)=0\), where \(F\in M_{k}(\Gamma _{00}(N),R)\). Then \(F=0\).\par
\par
(5.5) Further, let \(R_{0}\) be a subring of \(R\), and suppose that \((𝔞,𝔟,\varepsilon ,j)\) is in fact a cusp on \(\Gamma _{00}(N)\) over \(R_{0}\). Suppose that the \(q\)-expansion coefficients \(c_{\mu }\) of \(F(𝔞,𝔟,\varepsilon ,j)\) all lie in \(R_{0}\). Then \(F\) is an element of \(M_{k}(\Gamma _{00}(N),R_{0})\).\par
\par
[In connection with the conclusion of (5.5), we remark that there is a natural map \(M_{k}(\Gamma _{00}(N),R_{0})\to M_{k}(\Gamma _{00}(N),R)\) because every \(R\)-algebra \(R'\) is in particular an \(R_{0}\)-algebra. The injectivity of this map may be deduced from (5.4).]\par
\DAnchor{35}{261}
Over the complex field \(\mathbb{C}\), we may express quadruples \((X,\lambda ,\omega ,i)\) in terms of lattices \(\mathcal{L}\) in the \(\mathbb{C}\)-vector space \(K\otimes _{\mathbb{Q}}\) \(\mathbb{C}\). (We tacitly suppose all such lattices to be stable under \(\mathcal{O}.)\) Specifically, the data \((X,\lambda ,\omega ,i)\) amount to the giving of:\par
\par
\(A\) lattice \(\mathcal{L}\subset K\otimes \mathbb{C}\).\par
An \((\mathcal{O}\)-linear) isomorphism \(\lambda :\wedge ^{2}\mathcal{L}\xrightarrow{\sim}𝔇^{-1}𝔠^{-1}\).\par
An \((\mathcal{O}\)-linear) embedding \(i:N^{-1}𝔇^{-1}/𝔇^{-1}\hookrightarrow N^{-1}\mathcal{L}/\mathcal{L}\).\par
\par
We may view \(\lambda \) as an \(\mathcal{O}\)-linear alternating form \(\lambda (x,y)\) on \(\mathcal{L}\times \mathcal{L}\), giving rise by linearity to a \(K\otimes \mathbb{R}\)-linear alternating form on \(K\otimes \mathbb{C}\). For \(\alpha \in K\otimes \mathbb{C}\), we have\par
\par
\DFormula{\lambda (\alpha x,\alpha y)=(\alpha \overline{\alpha })\cdot \lambda (x,y),}
\par
where \(\overline{\alpha }\) is the conjugate of \(\alpha \) in \(K\otimes \mathbb{C}\).\par
\par
From this point of view, an element of \(M_{k}(\Gamma _{00}(N),\mathbb{C})\) is simply a holomorphic function \(F\) on the space of triples \((\mathcal{L},\lambda ,i)\) which satisfies:\par
\par
\DFormula{(5.6)\,\,\,\,F(\alpha \mathcal{L},(\alpha \overline{\alpha })^{-1}\lambda ,\alpha i)=\mathcal{N}\alpha ^{-k}F(\mathcal{L},\lambda ,i)}
\par
for \(\alpha \in (K\otimes \mathbb{C})^{\ast}\). Given such a function, we may obtain the \(q\)-expansion of the corresponding modular form at a cusp of the form \((𝔞,𝔟,\varepsilon ,j_{\mathrm{can}})\) (cf. (5.3)) by evaluating the function on a suitable triple \((\mathcal{L},\lambda ,i)\).\par
\par
This triple is constructed as follows. First, we introduce the lattice \(\mathcal{L}=2\pi i(𝔇^{-1}𝔞^{-1}+𝔟\tau )\), where \(\tau \) is a “variable” element of the “upper half plane”\par
\par
\DFormula{\mathfrak{H}=\lbrace \tau \in K\otimes \mathbb{C}\,\vert \,\operatorname{Im}\,\tau \gg 0\rbrace .}
\par
Second, we define an alternating form \(\lambda =\lambda _{\mathrm{can}}\) by the formula\par
\par
\DFormula{((2\pi i)(\alpha +\beta \tau ),(2\pi i)(\gamma +\delta \tau ))\mapsto \alpha \delta -\beta \gamma .}
\par
Finally, we let i be the embedding\par
\par
\DFormula{N^{-1}𝔇^{-1}/𝔇^{-1}\,\xrightarrow{\varepsilon }\,N^{-1}𝔞^{-1}𝔇^{-1}/𝔞^{-1}𝔇^{-1}\,\hookrightarrow \,N^{-1}\mathcal{L}/\mathcal{L},}
\par
the first map being the isomorphism deduced from \(\varepsilon \) and the second being induced by multiplication by \(2\pi i\).\par
\par
If we evaluate an \(F\) as above on \((\mathcal{L},\lambda ,i)\), we obtain a holomorphic function \(F(\tau )\) which is invariant under translations by elements of \(𝔞^{-1}𝔟^{-1}𝔇^{-1}\). It may consequently be expanded as a Fourier series\par
\par
\DFormula{F(\tau )=F_{𝔞,𝔟,\varepsilon }(\tau )=\sum _{\mu \in 𝔞𝔟}\,c_{\mu }e^{2\pi i\cdot \operatorname{tr}(\tau \cdot \mu )},}
\par
the summation being restricted to those elements \(\mu \) of \(𝔞𝔟\) which are either 0 or totally positive. (Here tr is the trace map \(K\otimes \mathbb{C}\to \mathbb{C}.)\) Then the numbers \(c_{\mu }\) are precisely the coefficients \(c_{\mu }(𝔞,𝔟,\varepsilon ,j_{\mathrm{can}})\) of the \(q\)-expansion of the modular form determined by \(F\), at the cusp (over \(\mathbb{C})\) determined by \(𝔞,𝔟\), and \(\varepsilon \). It is usual to make the formal substitution\par
\par
\DFormula{q^{\mu }=e^{2\pi i\cdot \operatorname{tr}(\mu \cdot \tau )};}
\DAnchor{36}{262}
this underscores the equality between the Fourier expansion of \(F(\tau )\) and the \(q\)-expansion of the corresponding modular form.\par
\par
We now wish to discuss “slashing” by matrices. Let us fix an integer \(k\geq 0\). Suppose that \(F\) is a complex-valued function on \(\mathfrak{H}\), and that \(M=\begin{pmatrix}\alpha&\beta\\\gamma&\delta\end{pmatrix}\) is an element of \(\operatorname{SL}(2,K\otimes \mathbb{R})\). We define \(F\vert M\) to be the function on \(\mathfrak{H}\) whose value at \(\tau \) is\par
\par
\DFormula{\mathcal{N}(\gamma \tau +\delta )^{-k}F((\alpha \tau +\beta )/(\gamma \tau +\delta )),}
\par
as in §0.\par
\par
In particular, this formula defines a right action of \(\operatorname{SL}(2,K)\) on the space of functions on \(\mathfrak{H}\). Because of the strong approximation theorem for \(\operatorname{SL}(2)\), we may define \(F\vert M\) whenever \(F\) is invariant by a congruence subgroup of \(\operatorname{SL}(2,K)\) and \(M\in \operatorname{SL}(2,\widehat{K})\) is a matrix with entries in the ring \(\widehat{K}\) of finite adeles of \(K\). (Explicitly, suppose that we have \(F\vert M=F\) for all \(M\) in a congruence subgroup \(\Gamma \) of \(\operatorname{SL}(2,K)\). Let \(\widehat{\Gamma }\) be the closure of \(\Gamma \) in \(\operatorname{SL}(2,\widehat{K})\). Given \(M\in \operatorname{SL}(2,\widehat{K})\), we may write \(M=M_{1}M_{2}\) with \(M_{1}\in \widehat{\Gamma }\) and \(M_{2}\in \operatorname{SL}(2,K)\). We define \(F\vert M\) to be \(F\vert M_{2}.)\)\par
\par
As is well known, modular forms of weight \(k\) on \(\Gamma _{00}(N)\) over \(\mathbb{C}\) correspond to holomorphic functions of \(\tau \) which are invariant (under slashing) by suitable subgroups of \(\operatorname{SL}(2,K)\). To establish this correspondence, we fix ideals \(𝔞\) and \(𝔟\) of \(K\) such that \(𝔞𝔟^{-1}=𝔠\), together with an isomorphism\par
\par
\DFormula{\varepsilon :N^{-1}\mathcal{O}/\mathcal{O}\xrightarrow{\sim}N^{-1}𝔞^{-1}/𝔞^{-1}.}
\par
Each lattice \(\mathcal{L}\subset K\otimes \mathbb{C}\), equipped with a polarization \(\lambda :\wedge ^{2}\mathcal{L}\xrightarrow{\sim}𝔇^{-1}𝔠^{-1}\) and a level-N structure\par
\par
\DFormula{i:𝔇^{-1}N^{-1}/𝔇^{-1}\hookrightarrow N^{-1}\mathcal{L}/\mathcal{L}}
\par
is isomorphic for some \(\tau \in \mathfrak{H}\) to the lattice \(2\pi i(𝔞^{-1}𝔇^{-1}+𝔟\tau )\), equipped with the canonical polarization \(\lambda _{\mathrm{can}}\) and the inclusion \(i=i(\varepsilon )\) discussed above. Furthermore, two elements \(\tau \) and \(\tau '\) of \(\mathfrak{H}\) give isomorphic lattices precisely when we have\par
\par
\DFormula{\tau '=(\alpha \tau +\beta )/(\gamma \tau +\delta )}
\par
for some \(\begin{pmatrix}\alpha&\beta\\\gamma&\delta\end{pmatrix}\) in \(\operatorname{SL}(2,K)\) such that:\par
\par
\DFormula{\alpha ,\delta \in \mathcal{O},\,\,\,\,\,\,\,\,\alpha \equiv \delta \equiv 1\,\operatorname{mod}\,N\mathcal{O},}
\DFormula{\gamma \in N𝔞𝔟𝔇,\,\,\,\,\,\,\beta \in 𝔞^{-1}𝔟^{-1}𝔇^{-1}.}
\par
We denote by \(\Gamma _{00}(N;𝔞,𝔟)\), or simply \(\Gamma _{00}(N)\), the subgroup of \(\operatorname{SL}(2,K)\) consisting of matrices which satisfy these conditions. Then we have\par
\par
(5.7) Proposition. If \(F\in M_{k}(\Gamma _{00}(N),\mathbb{C})\) is a modular form on \(\Gamma _{00}(N)\) of weight \(k\), then the function \(F_{𝔞,𝔟,\varepsilon }(\tau )\) on \(\mathfrak{H}\) is invariant under slashing by the group \(\Gamma _{00}(N)\). Conversely, any holomorphic function on \(\mathfrak{H}\) which is invariant under \(\Gamma _{00}(N)\) is of the form \(F_{𝔞,𝔟,\varepsilon }(\tau )\) for some \(F\in M_{k}(\Gamma _{00}(N),\mathbb{C})\).\par
\DAnchor{37}{263}
Remark. \(A\) \(\Gamma _{00}(N)\)-invariant function on \(\mathfrak{H}\) is holomorphic at infinity because of our assumption \(K\neq \mathbb{Q}\).\par
\par
We now wish to consider the various \(q\)-expansions of a form \(F\in M_{k}(\Gamma _{00}(N),\mathbb{C})\) at the cusps defined by elements of \(\widehat{K}^{\ast}\), as in (5.3). For each element \(x\) of \(\widehat{K}^{\ast}\) (we have another use now for the letter \(\alpha )\), the \(q\)-expansion of \(F\) at the cusp determined by \(x\) is the formal power series corresponding to the Fourier series \(F_{𝔞,𝔟,\varepsilon }(\tau )\), where \(𝔞\), \(𝔟\), and \(\varepsilon \) are derived from \(x\) as in (5.3). We write simply \(F_{x}(\tau )\) for this function.\par
\par
(5.8) Proposition. For each \(x\in \widehat{K}^{\ast}\), we have\par
\par
\DFormula{F_{x}=F_{1}\,\vert \,\begin{pmatrix}x&0\\0&x^{-1}\end{pmatrix}.}
\par
Proof. By (5.7), \(F_{1}(\tau )\) is invariant under the group \(\Gamma _{00}(N;\mathcal{O},𝔠^{-1})\). According to the definition of \(F_{1}\) \(\vert \) \(\begin{pmatrix}x&0\\0&x^{-1}\end{pmatrix}\), we may compute this function by slashing \(F_{1}\) by any matrix \(\begin{pmatrix}\alpha&\beta\\\gamma&\delta\end{pmatrix}\in \operatorname{SL}(2,K)\) subject to the conditions:\par
\par
\DFormula{\alpha \in 𝔞,\,\,\,\,\,\,\delta \in 𝔞^{-1},\,\,\,\,\,\,\beta \in 𝔇^{-1}𝔟^{-1},\,\,\,\,\,\,\gamma \in N𝔇𝔟.}
\DFormula{\alpha \equiv x\,\operatorname{mod}\,N\widehat{𝔞},\,\,\,\,\,\,\delta \equiv x^{-1}\,\operatorname{mod}\,N\widehat{𝔞}^{-1}.}
\par
(Here we have put \(𝔞=(x)\), \(𝔟=𝔠^{-1}𝔞\). Cf. [10, pp. 234--235].)\par
\par
We view the modular form \(F\) as a function of triples \((\mathcal{L},\lambda ,i)\). Since \(F\) has weight \(k\), we find for each \(\tau \in \mathfrak{H}\) that \([F_{1}\) \(\vert \) \(\begin{pmatrix}\alpha&\beta\\\gamma&\delta\end{pmatrix}](\tau )\) is the value of \(F\) on the lattice\par
\par
\DFormula{(5.9)\,\,\,\,2\pi i[𝔇^{-1}(\gamma \tau +\delta )+𝔠^{-1}(\alpha \tau +\beta )],}
\par
with an appropriate polarization and level structure. The conditions satisfied by \(\alpha ,\beta ,\gamma \), and \(\delta \) insure that the lattice (5.9) is \(2\pi i(𝔇^{-1}𝔞^{-1}+𝔟\tau )\) with the polarization \(\lambda _{\mathrm{can}}\) and the \(\Gamma _{00}(N)\)-structure \(i(\varepsilon )\), where \(\varepsilon \) is the isomorphism defined by \(x\) as in (5.3). We omit details of this calculation.\par
\par
The Modular Forms of §0\par
\par
According to (5.7), the weight-k modular forms defined in the Introduction are precisely the standard \(q\)-expansions \(F_{1}(\tau )\) of forms \(F\in M_{k}(\Gamma _{00}(N),\mathbb{C})\), with the polarization module \(𝔠\) chosen to be \(\mathcal{O}\). The “other” \(q\)-expansions of these forms were defined by a formula which is now justified by (5.8). Therefore (0.1) is a special case of the \(q\)-expansion principle (5.4), (5.5).\par
\par
\(p\)-adic Modular Forms\par
\par
Let \(p\) be a prime. \(A\) \(\Gamma _{00}(Np^{\infty })\)-structure on a HBAV \(X\) is a compatible system of \(\Gamma _{00}(Np^{n})\)-structures, for \(n\geq 0\). Let \(R\) be a \(p\)-adic ring: \(R\xrightarrow{\sim}\varprojlim \) \(R/p^{n}R\). \(A\) \(p\)-adic\par
\DAnchor{38}{264}
Hilbert modular form on \(\Gamma _{00}(N)\) over \(R\) is a function \(F\) of triples \((X,\lambda ,i)\), where:\par
\par
\(X\) is a HBAV over a \(p\)-adic \(R\)-algebra \(R'\),\par
\(\lambda :\mathcal{P}(X)\xrightarrow{\sim}𝔠^{-1}\) is a positive isomorphism,\par
i is a \(\Gamma _{00}(Np^{\infty })\)-structure on \(X\).\par
\par
It is required that \(F(X,\lambda ,i)\) depend only on the \(R'-isomorphism\) class of \((X,\lambda ,i)/R'\) and that \(F\) commute with extensions of scalars \(R'\to R''\).\par
\par
We let \(V(\Gamma _{00}(N),R)\) be the \(R\)-module of \(p\)-adic Hilbert modular forms over \(R\).\par
\par
The key point concerning \(p\)-adic modular forms is that there is a canonical way to associate a \(p\)-adic modular form to each (Hilbert) modular form, of any weight, over \(R\). We recall that whenever \(X/R\) has a \(\Gamma _{00}(Np^{\infty })\)-structure i, we may deduce from i an isomorphism of formal groups\par
\par
\DFormula{𝔇^{-1}\otimes _{\mathbb{Z}}\,\widehat{G}_{m}\,\xrightarrow{\sim}\,\widehat{X}}
\par
and hence an isomorphism of their Lie algebras\par
\par
\DFormula{𝔇^{-1}\otimes R\xrightarrow{\sim}\operatorname{Lie}(X/R).}
\par
There is thus a canonically chosen basis \(\omega (i)\) of \(\underline{\omega }_{X/R}\).\par
\par
\subsection*{The map on test objects}
\par
\DFormula{(X,\lambda ,i)\mapsto (X,\lambda ,\omega (i),i)}
\par
induces a map\par
\par
\DFormula{M_{k}(\Gamma _{00}(N),R)\to V(\Gamma _{00}(N),R)}
\par
for each \(k\). We denote this map by \(\pi _{k}\), or simply \(\pi \).\par
\par
To make \(q\)-expansions for \(p\)-adic modular forms, let \(𝔞\) and \(𝔟\) be ideals of \(K\) with \(𝔞𝔟^{-1}=𝔠\) and let \(\varepsilon =(\varepsilon _{n})\) be a compatible system of isomorphisms\par
\par
\DFormula{\varepsilon _{n}:p^{-n}N^{-1}\mathcal{O}/\mathcal{O}\xrightarrow{\sim}p^{-n}N^{-1}𝔞^{-1}/𝔞^{-1}.}
\par
These data provide us with a “Tate variety” \(X_{𝔞,𝔟}\), as before, together with a \(\Gamma _{00}(Np^{\infty })\)-structure on \(X_{𝔞,𝔟}\), which we denote by \(i(\varepsilon )\). Given a \(p\)-adic modular form \(F\), its \(q\)-expansion \(F(𝔞,𝔟,\varepsilon )\) is obtained by evaluating \(F\) on the triple consisting of \(X_{𝔞,𝔟}\), the canonical isomorphism\par
\par
\DFormula{\lambda :\mathcal{P}(X_{𝔞,𝔟})\xrightarrow{\sim}𝔠,}
\par
and the level structure \(i(\varepsilon )\). (We refer to \((𝔞,𝔟,\varepsilon )\) as a “\(p\)-adic cusp.”)\par
\par
It is easy to describe the differential \(\omega (i(\varepsilon ))\) on \(X_{𝔞,𝔟}\) in terms of an isomorphism\par
\par
\DFormula{j=j(\varepsilon ):𝔞^{-1}\otimes R\xrightarrow{\sim}\mathcal{O}\otimes R.}
\par
Namely, the \(\varepsilon _{n}\) give isomorphisms\par
\par
\DFormula{p^{-n}\mathcal{O}/\mathcal{O}\xrightarrow{\sim}p^{-n}𝔞^{-1}/𝔞^{-1}}
\par
for each n, and hence an isomorphism\par
\par
\DFormula{\mathcal{O}\otimes \mathbb{Z}_{p}\xrightarrow{\sim}𝔞^{-1}\otimes \mathbb{Z}_{p}.}
\DAnchor{39}{265}
Tensoring with \(R\) gives an isomorphism\par
\par
\DFormula{\mathcal{O}\otimes R\xrightarrow{\sim}𝔞^{-1}\otimes R,}
\par
which is just the inverse of the desired map j. This gives the compatibility:\par
\par
(5.10) If \(F\in M_{k}(\Gamma _{00}(N),R)\) is a modular form of weight \(k\), and \((𝔞,𝔟,\varepsilon )\) is a \(p\)-adic cusp, then\par
\par
\((\pi _{k}F)(𝔞,𝔟,\varepsilon )=F(𝔞,𝔟,\varepsilon _{0},j(\varepsilon ))\).\par
\par
(5.11) Example. Let \(\alpha \in \widehat{K}^{\ast}\). \(A\) \(p\)-adic cusp \((𝔞,𝔟,\varepsilon )\) is defined by \(\alpha \) as follows. We take \(𝔞=(\alpha )\) and \(𝔟=𝔞𝔠^{-1}\), as in (5.3), and for each \(n\geq 0\) we define \(\varepsilon _{n}\) to be the composite\par
\par
\DFormula{p^{-n}N^{-1}\mathcal{O}/\mathcal{O}\,\xrightarrow{\alpha ^{-1}}\,N^{-1}p^{-1}\widehat{𝔞}^{-1}/\widehat{𝔞}^{-1}\approx p^{-n}N^{-1}𝔞^{-1}/𝔞^{-1},}
\par
cf. (5.3). Let \(\alpha _{p}\) be the image of \(\alpha \) in \((K\otimes \mathbb{Q}_{p})^{\ast}\). Multiplication by \(\alpha _{p}\) induces an isomorphism\par
\par
\DFormula{(5.12)\,\,\,\,𝔞^{-1}\otimes \mathbb{Z}_{p}\xrightarrow{\sim}\mathcal{O}\otimes \mathbb{Z}_{p}.}
\par
Tensoring with \(R\), we obtain an isomorphism between \(𝔞^{-1}\otimes R\) and \(\mathcal{O}\otimes R\), which is just the map \(j(\varepsilon )\) discussed above.\par
\par
We now consider the situation where \(R\) is a flat \(\mathbb{Z}_{p}\)-module. We define \(V(\Gamma _{00}(N),E)\), where \(E=R\otimes \mathbb{Q}_{p}\), to be the tensor product\par
\par
\DFormula{V(\Gamma _{00}(N),R)\otimes _{\mathbb{Z}_{p}}\mathbb{Q}_{p}.}
\par
\subsection*{Then for each \(k\) we have a map}
\par
\DFormula{\pi _{k}:M_{k}(\Gamma _{00}(N),E)=M_{k}(\Gamma _{00}(N),R)\otimes \mathbb{Q}_{p}\to V(\Gamma _{00}(N),E),}
\par
obtained by tensoring the previous \(\pi _{k}\) with \(\mathbb{Q}_{p}\). Given data \((𝔞,𝔟,\varepsilon )\) as above, we may define for each \(F\in V(\Gamma _{00}(N),E)\) its \(q\)-expansion \(F(𝔞,𝔟,\varepsilon )\) at \((𝔞,𝔟,\varepsilon )\). Then we have the following \(q\)-expansion principle.\par
\par
(5.13) Let \(F\in V(\Gamma _{00}(N),E)\), and let \((𝔞,𝔟,\varepsilon )\) be a \(p\)-adic cusp. Then the \(q\)-expansion \(F(𝔞,𝔟,\varepsilon )\), a priori a power series with coefficients in \(E\), in fact has coefficients in \(R\) if and only if \(F\) is an element of \(V(\Gamma _{00}(N),R)\).\par
\par
(5.14) Corollary. Suppose that \(F\in V(\Gamma _{00}(N),E)\) has at one cusp \((𝔞,𝔟,\varepsilon )\) a \(q\)-expansion whose non-constant coefficients are all elements of \(R\). Then the difference between the constant terms of the \(q\)-expansions of \(F\) at any two cusps is an element of \(R\).\par
\par
Proof. Let \(t\in E\) be the constant term of \(F(𝔞,𝔟,\varepsilon )\). We may view \(t\) as an element of \(M_{0}(\Gamma _{00}(N),E)\) and hence in particular as an element of \(V(\Gamma _{00}(N),E)\). The \(q\)-expansion of \(t\) at each cusp is the constant power series \(t\). By the hypothesis on \(F\), and because of (5.13), \(F-t\) belongs to \(V(\Gamma _{00}(N),R)\). The \(q\)-expansion of this form at each cusp has coefficients in \(R\), and so in particular its constant term belongs to \(R\). This proves the required “integrality.”\par
\DAnchor{40}{266}
(5.15) Let \(F\in M_{k}(\Gamma _{00}(N),E)\), and for \(\alpha \in \widehat{K}^{\ast}\), let \(F_{\alpha }\) denote the \(q\)-expansion of \(F\) at the cusp determined by \(\alpha \) (5.3). Then the \(q\)-expansion of \(\pi _{k}F\) at the \(p\)-adic cusp determined by \(\alpha \) (5.11) is given by\par
\par
\DFormula{\mathcal{N}\alpha _{p}^{-k}F_{\alpha },}
\par
where \(\alpha _{p}\in (K\otimes \mathbb{Q}_{p})^{\ast}\) is the component of \(\alpha \) at \(p\). To see this, we write the \(p\)-adic \(q\)-expansion as\par
\par
\DFormula{F(𝔞,𝔟,\varepsilon _{0},j),}
\par
where j is the isomorphism \(𝔞^{-1}\otimes E\xrightarrow{\sim}\mathcal{O}\otimes E\) induced by multiplication by \(\alpha _{p}\), viewed as a map \(𝔞^{-1}\otimes \mathbb{Z}_{p}\xrightarrow{\sim}\mathcal{O}\otimes \mathbb{Z}_{p}\). (The ideals \(𝔞\) and \(𝔟\), and the map \(\varepsilon _{0}\), are defined in (5.11).) The \(q\)-expansion \(F_{\alpha }\) is\par
\par
\DFormula{F(𝔞,𝔟,\varepsilon _{0},j_{\mathrm{can}}),}
\par
where \(j_{\mathrm{can}}\) is the isomorphism \(𝔞^{-1}\otimes E\xrightarrow{\sim}\mathcal{O}\otimes E\) derived from the equality \(𝔞^{-1}\otimes \mathbb{Q}=\mathcal{O}\otimes \mathbb{Q}\). We thus have \(j=\alpha _{p}\cdot j_{\mathrm{can}}\), so that desired formula follows from (5.2).\par
\par
Now for each \(k\geq 1\), let \(F_{k}\) be an element of \(M_{k}(\Gamma _{00}(N),E)\). Assume that \(F_{k}=0\) for \(k\) sufficiently large. Then, by the above discussion, the \(p\)-adic modular form \(\sum _{k}\pi _{k}F_{k}\) has \(q\)-expansion\par
\par
\DFormula{\sum _{k}\mathcal{N}\alpha _{p}^{-k}F_{k,\alpha }}
\par
at the \(p\)-adic cusp determined by \(\alpha \). (Here we have written \(F_{k,\alpha }\) for the \(q\)-expansion \((F_{k})_{\alpha }\) of \(F_{k}\) at the cusp determined by \(\alpha .)\) By (5.13), we see that this power series has coefficients in \(R\) for one \(\alpha \) if and only if it has coefficients in \(R\) for each \(\alpha \). Similarly, by (5.14), if the non-constant coefficients of one of these series all lie in \(R\), then the difference between the constant coefficients of any two of these series lies in \(R\).\par
\par
We thus see, in particular, how (0.2) and (0.3) are special cases of (5.13) and (5.14), respectively.\par
\par
Variant: Forms on \(\Gamma _{00}(𝔣)\)\par
\par
We suppose that \(𝔣\) is a conductor, i.e., a non-zero integral ideal of \(K\). \(A\) \(\Gamma _{00}(𝔣)\)-structure on a HBAV \(X\) is an embedding\par
\par
\DFormula{(𝔣^{-1}𝔇^{-1}/𝔇^{-1})(1)\hookrightarrow X.}
\par
\(A\) modular form of weight \(k\) on \(\Gamma _{00}(𝔣)\) over \(R\) is a function of tuples \((X,\lambda ,\omega ,i)\) as before, with i now a \(\Gamma _{00}(𝔣)\)-structure on \(X\) over \(R\). We denote the space of such forms by \(M_{k}(\Gamma _{00}(𝔣),R)\).\par
\par
Let \(N\) be a positive integer divisible by \(𝔣\). Then an element of \(M_{k}(\Gamma _{00}(𝔣),R)\) is simply an element of \(M_{k}(\Gamma _{00}(N),R)\) which enjoys the following extra invariance: its value on a tuple \((X,\lambda ,\omega ,i)\), where i is a \(\Gamma _{00}(N)\)-structure on \(X\), depends only on \(X,\lambda ,\omega \), and the composite\par
\par
\DFormula{(𝔣^{-1}𝔇^{-1}/𝔇^{-1})(1)\hookrightarrow (N^{-1}𝔇^{-1}/𝔇^{-1})(1)\,\xrightarrow{i}\,X.}
\DAnchor{41}{267}
In particular, a \(q\)-expansion \(F(𝔞,𝔟,\varepsilon ,j)\) is defined for each form \(F\) on \(\Gamma _{00}(𝔣)\) and each cusp \((𝔞,𝔟,\varepsilon ,j)\) on \(\Gamma _{00}(N)\). When \(R\) is a \(\mathbb{Q}\)-algebra, we may thus speak of the \(q\)-expansion of \(F\) at the cusp determined by an \(\alpha \in \widehat{K}^{\ast}\).\par
\par
We may similarly define \(p\)-adic modular forms on \(\Gamma _{00}(𝔣)\).\par
\par
\subsection*{§6. Eisenstein Series}
\par
Let \(𝔅\) be an ideal of \(K\). In this § we shall construct Eisenstein series for the Hilbert group, following Hecke and Siegel. The polarization module \(𝔠\) of §5 will be the ideal \(𝔅^{-1}\).\par
\par
For \(\varepsilon \) a Schwartz function on \(I\), we define the modification of \(\varepsilon \) (with respect to \(𝔅)\) to be the unique Schwartz function \(\widetilde{\varepsilon }\) on \(I\) which is supported on \(I^{𝔅}\) and whose value on an ideal \(𝔵\subseteq 𝔅\) is given by \(\varepsilon (𝔵^{-1}\cdot 0)\). Since this quantity depends only on the strict ideal class of \(𝔵\), \(\widetilde{\varepsilon }\) is defined mod \(𝔣\) for every conductor \(𝔣\subseteq 𝔅\). For \(c\in G\), we define \(\widetilde{\varepsilon }_{c}\) to be \((\widetilde{\varepsilon })_{c}\), i.e., the twist by c of the modification of \(\varepsilon \).\par
\par
(6.1) Theorem. Let \(k\geq 1\) be an integer, and let \(\varepsilon :I\to \mathbb{C}\) be a function with parity \((-1)^{k}\) which is supported on \(A\) and defined modulo the conductor \(𝔣\). Then there exists a modular form\par
\par
\DFormula{G_{k,\varepsilon }\in M_{k}(\Gamma _{00}(𝔣),\mathbb{C})}
\par
whose \(q\)-expansion at the cusp determined by each \(\alpha \in \widehat{K}^{\ast}\) is given by the formulas\par
\par
(6.2) If \(k>1\)\par
\par
\DFormula{\mathcal{N}\alpha ^{k}\lbrace 2^{-r}L(1-k,\varepsilon _{c})+\sum _{\mu \gg 0}(\sum _{𝔵\subset 𝔅𝔞^{2}}\,\varepsilon _{c}(\mu 𝔵^{-1})\mathcal{N}(\mu 𝔵^{-1})^{k-1})q^{\mu }\rbrace .}
\par
If \(k=1\)\par
\par
\DFormula{\mathcal{N}\alpha \lbrace 2^{-r}L(0,\varepsilon _{c}+\widetilde{\varepsilon }_{c})+\sum _{\mu \gg 0}(\sum _{𝔵\subset 𝔅𝔞^{2}}\,\varepsilon _{c}(\mu 𝔵^{-1}))q^{\mu }\rbrace .}
\par
Here:\par
\par
\(𝔞=(\alpha )\) is the ideal of \(K\) “generated” by \(\alpha \),\par
\(c=(𝔞\cdot \alpha ^{-1})\) is the element of \(G\) defined by \(\alpha \), cf. (2.23).\par
\par
Proof. We will construct \(G_{k,\varepsilon }\) as a function of \((\mathcal{O}-)lattices\) \(\mathcal{L}\) in \(K\otimes \mathbb{C}\). We are required from this point of view to produce a complex number \(G_{k,\varepsilon }(\mathcal{L})\) each time that we are given a lattice \(\mathcal{L}\) together with the supplementary data of a \(𝔠-polarization\) and a \(\Gamma _{00}(𝔣)\) structure on \(\mathcal{L}\). Given such data, we first ignore the polarization completely (!). We then view the \(\Gamma _{00}(𝔣)\) structure as the giving of an overlattice \(\mathcal{L}'\supseteq \mathcal{L}\) together with an isomorphism\par
\par
\DFormula{\mathcal{L}'/\mathcal{L}\approx 𝔇^{-1}𝔣^{-1}/𝔇^{-1}.}
\par
We consider the various \(\mathcal{O}-submodules\) \(\mathcal{M}\) of \(\mathcal{L}'\) which are invertible, i.e., of rank 1. Given an \(\mathcal{M}\), we choose arbitrarily a positivity \(+\) for \(\mathcal{M}\) (2.25). We will then evaluate the Fourier transform \(T\varepsilon \) of \(\varepsilon \) on an element of \(I\) which is defined\par
\DAnchor{42}{268}
mod \(𝔇^{-1}\) in the sense of (3.9). This element is obtained as in (2.25) from the triple consisting of \(\mathcal{M}\), the positivity \(+\), and the map given as the composite\par
\par
\DFormula{\mathcal{M}\subset \mathcal{L}'\to \mathcal{L}'/\mathcal{L}\approx 𝔇^{-1}𝔣^{-1}/𝔇^{-1}\hookrightarrow K/𝔇^{-1}.}
\par
We write \(T\varepsilon (\mathcal{M}\to K/𝔇^{-1})\) for the value of \(T\varepsilon \) on this element.\par
\par
We further associate a volume to \((\mathcal{M},+)\). Namely, let \(𝔞\) be an ideal which defines the same strict ideal class as \((\mathcal{M},+)\). Choose a positive isomorphism \((𝔞,+_{\mathrm{can}})\approx (\mathcal{M},+)\). This isomorphism induces an inclusion \(K\hookrightarrow K\otimes \mathbb{C}\), and if \(m\in K\otimes \mathbb{C}\) is the image of \(1\in K\), then \(\mathcal{M}=𝔞\cdot m\). We define\par
\par
\DFormula{\operatorname{vol}(\mathcal{M},+)=\mathcal{N}𝔞\cdot \mathcal{N}m\in \mathbb{C},}
\par
the product of the norms of the ideal \(𝔞\) and the number \(m\in K\otimes \mathbb{C}\). The volume is well defined because replacing \((\mathcal{M},+)\approx (𝔞,+_{\mathrm{can}})\) by another isomorphism would only replace \(𝔞\) by \(𝔞\gamma ^{-1}\) and m by \(\gamma m\), where \(\gamma \) is totally positive, and in particular of positive norm.\par
\par
Similarly, if we modify \(+\) at one real place v, then \(\operatorname{vol}(\mathcal{M},+)\) is replaced by its negative, so that \(\operatorname{vol}(\mathcal{M},+)^{k}\) is multiplied by \((-1)^{k}\). However, because \(\varepsilon \) has parity \((-1)^{k}\), so that \(T\varepsilon \) again has this parity, changing \(+\) at v multiplies \(T\varepsilon (\mathcal{M},+,\mathcal{M}\to K/𝔇^{-1})\) by this same factor. Hence the quotient\par
\par
\DFormula{T\varepsilon (\mathcal{M},+,\alpha :\mathcal{M}\to K/𝔇^{-1})\,/\,\operatorname{vol}(\mathcal{M},+)^{k}}
\par
is independent of the choice of \(+\), so that a quantity\par
\par
\DFormula{T\varepsilon (\mathcal{M},\mathcal{M}\to K/𝔇^{-1})\,/\,(\operatorname{vol}\mathcal{M})^{k}}
\par
is well defined. We define the value of \(G_{k,\varepsilon }\) on \(\mathcal{L}\) to be \([(-1)^{k}(k-1)!]^{r}\cdot S\), where \(S\) is the sum\par
\par
\DFormula{\sum _{\mathcal{M}\subset \mathcal{L}'}\,T\varepsilon (\mathcal{M},\mathcal{M}\to K/𝔇^{-1})/(\operatorname{vol}\mathcal{M})^{k},}
\par
computed in the sense of analytic continuation:\par
\par
\DFormula{(6.3)\,\,\,\,S=[\sum _{\mathcal{M}\subset \mathcal{L}'}\,T\varepsilon (\mathcal{M},\mathcal{M}\to K/𝔇^{-1})/((\operatorname{vol}\mathcal{M})^{k}\vert \operatorname{vol}\mathcal{M}\vert ^{s})]_{s=0}.}
\par
[The sum is absolutely convergent for \(k>2\); the continuation is necessary only for \(k=1,2.]\)\par
\par
To verify that this sum defines a modular form whose \(q\)-expansion is that given in (6.2), we must compute the value of \(G_{k,\varepsilon }\) at the cusp defined by \(𝔅,\alpha \). That is, we take\par
\par
\DFormula{\mathcal{L}=2\pi i(𝔞^{-1}𝔇^{-1}+𝔞𝔅\tau )}
\DFormula{\mathcal{L}'=2\pi i(𝔞^{-1}𝔣^{-1}𝔇^{-1}+𝔞𝔅\tau )}
\DAnchor{43}{269}
together with the isomorphism\par
\par
\DFormula{\mathcal{L}'/\mathcal{L}=𝔞^{-1}𝔣^{-1}𝔇^{-1}/𝔞^{-1}𝔇^{-1}\xrightarrow{\sim}𝔣^{-1}𝔇^{-1}/𝔇^{-1}}
\par
induced by multiplication by \(\alpha \): \(𝔞^{-1}𝔣^{-1}𝔇^{-1}/𝔞^{-1}𝔇^{-1}\to \widehat{𝔣^{-1}𝔇^{-1}}/\widehat{𝔇}^{-1}\). The map \(\mathcal{L}'\to \widehat{𝔇^{-1}𝔣^{-1}}/\widehat{𝔇}^{-1}\) is then just the composition of the three maps\par
\par
division by \(2\pi i\) and projection: \(\mathcal{L}'\subset 2\pi i(𝔞^{-1}𝔣^{-1}𝔇^{-1}+𝔞𝔅\tau )\to 𝔞^{-1}𝔣^{-1}𝔇^{-1}\);\par
\par
multiplication: \(𝔇^{-1}𝔞^{-1}𝔣^{-1}\) \(\xrightarrow{\alpha }\) \(\widehat{𝔣^{-1}𝔇^{-1}}\);\par
\par
the canonical map: \(\widehat{𝔣^{-1}𝔇^{-1}}\to \widehat{𝔣^{-1}𝔇^{-1}}/\widehat{𝔇}^{-1}\).\par
\par
Hence, for \(\mathcal{M}\subset \mathcal{L}'\), \(T\varepsilon (\mathcal{M},+,\mathcal{M}\to \widehat{𝔣^{-1}𝔇^{-1}}/\widehat{𝔇}^{-1})\) is the value of \(T\varepsilon \) on the element of \(I\) given by the triple\par
\par
\DFormula{(\mathcal{M},+,\mathcal{M}\to 𝔞^{-1}𝔣^{-1}𝔇^{-1}\,\xrightarrow{\alpha }\,\widehat{K}).}
\par
The sum (6.3) may now be written\par
\par
\DFormula{S=(2\pi i)^{-kr}\left[\sum_{\substack{\mathcal M\subset K+\mathfrak a\mathfrak B\tau\\\mathcal M\text{ invertible}}}\frac{(T\varepsilon)(\mathcal M,+,\mathcal M\xrightarrow{\mathrm{pr}}K\xrightarrow{\alpha}\widehat K)}{\operatorname{vol}(\mathcal M,+)^k\lvert\operatorname{vol}(\mathcal M,+)\rvert^s}\right]_{s=0},}
\par
where the “pr” is simply projection \(\mathcal{M}\to K\) onto the first factor. Indeed, let \(\mathcal{M}\subset K+𝔞𝔅\tau \) be an invertible \(\mathcal{O}\)-module. If \(\operatorname{pr}(\mathcal{M})\) is the ideal \(𝔟\), then \((\mathcal{M},+,\mathcal{M}\) \(\xrightarrow{\operatorname{pr}}\) \(K\) \(\xrightarrow{\alpha }\) \(\widehat{K})\) is simply \((𝔟\cdot \alpha )\) for a suitable choice of \(+\). Since \((𝔟\cdot \alpha )\in I^{𝔣^{-1}𝔇^{-1}}\) if and only if \(𝔟𝔞\subset 𝔣^{-1}𝔇^{-1}\), \(\mathcal{M}\) gives a non-zero contribution to the sum only if it is contained in \(𝔞^{-1}𝔣^{-1}𝔇^{-1}+𝔞𝔅\tau =(2\pi i)^{-1}\mathcal{L}'\).\par
\par
Let \(\eta (x)=(T\varepsilon )(x\cdot \alpha )\), \(𝔟=𝔅𝔞\). Then we have\par
\par
\DFormula{S=(2\pi i)^{-kr}\sum_{\mathcal M\subset K+\mathfrak b\tau}\frac{\eta(\mathcal M,+,\mathcal M\xrightarrow{\mathrm{pr}}K)}{\operatorname{vol}(\mathcal M,+)^k},}
\par
calculated (if necessary) by analytic continuation.\par
\par
Let \(\mathcal{M}\subset K+𝔟\tau \), and suppose that \(𝔲\) is an ideal isomorphic to \(\mathcal{M}\) as an invertible \(\mathcal{O}\)-module. Then \(\mathcal{M}=𝔲\cdot (\beta +\lambda \tau )\) for some \(\beta ,\lambda \in K\), and the pair \((\beta ,\lambda )\) is well defined up to multiplication by a unit of \(K\). Fix one choice \((\beta ,\lambda )\) and endow \(\mathcal{M}\) with the positivity induced by \(+_{\mathrm{can}}\) on \(𝔲\) via \((\beta +\lambda \tau ):𝔲\xrightarrow{\sim}\mathcal{M}\). For this positivity we have \(\operatorname{vol}\mathcal{M}=(\mathcal{N}𝔲)(\mathcal{N}(\beta +\lambda \tau ))\), and the element of \(I\) given by \((\mathcal{M},\mathcal{M}\to K)\) is \((𝔲\cdot \beta )\). Hence \(S\) may be written\par
\par
\((2\pi i)^{-kr}\sum _{𝔲}\) \(1/\mathcal{N}𝔲^{k}\) \(\lbrace \sum _{\beta \in K^{\ast}/U}\,\eta _{𝔲}(\beta )/\mathcal{N}\beta ^{k}\,+\,\sum _{0\neq \lambda \in 𝔟𝔲^{-1}\,\operatorname{mod}\,U}\sum _{\beta \in K}\,\eta _{𝔲}(\beta )/\mathcal{N}(\lambda \tau +\beta )^{k}\rbrace \),\par
\par
where the outer sum is over a set \(\lbrace 𝔲\rbrace \) of representations for the set of wide ideal classes of \(K\), and where \(\eta _{𝔲}(\beta )=\eta (𝔲\cdot \beta )\) as in §3.\par
\par
Thus \(S\) is naturally the sum of two terms, the first being\par
\par
\((2\pi i)^{-kr}L(k,\eta )=(2\pi i)^{-kr}L(k,(T\varepsilon )(x\cdot \alpha ))\).\par
\DAnchor{44}{270}
Changing variables, we rewrite this as\par
\par
\((2\pi i)^{-kr}L(k,T\varepsilon (c^{-1}x))\cdot \mathcal{N}𝔞^{k}\)                    \((c=𝔞\cdot \alpha ^{-1})\)\par
\par
\(=\) \([\mathcal{N}𝔞^{k}\cdot 2^{-r}/[(k-1)!]^{r}]\) \(L(1-k,T(T\varepsilon (c^{-1}x)))\)          (by (3.13))\par
\par
\(=\) \([\mathcal{N}𝔞^{k}\cdot 2^{-r}/[(k-1)!]^{r}]\) \(L(1-k,(T^{2}\varepsilon )(cx))\)              (by the invariance (3.7))\par
\par
\(=\) \([\mathcal{N}𝔞^{k}\cdot 2^{-r}/[(k-1)!]^{r}]\) \(L(1-k,\varepsilon (cx\cdot -1))\)               (by (3.5))\par
\par
\(=\) \([(-1)^{rk}\mathcal{N}𝔞^{k}2^{-r}/[(k-1)!]^{r}]\) \(L(1-k,\varepsilon _{c})\)            (by the parity of \(\varepsilon )\).\par
\par
Since \(G_{k,\varepsilon }=(-1)^{rk}[(k-1)!]^{r}\cdot S\), this first contribution to \(S\) gives precisely the term \(2^{-r}\mathcal{N}𝔞^{k}L(1-k,\varepsilon _{c})\) in (6.2).\par
\par
For the second term, we must find the value at \(s=0\) of the series\par
\par
\DFormula{\sum _{0\neq \lambda \in 𝔟𝔲^{-1}\,\operatorname{mod}\,U}\sum _{\beta \in K}\,\eta _{𝔲}(\beta )/\mathcal{N}(\lambda \tau +\beta )^{k+\vert s\vert }.}
\par
(We write \(x^{k+\vert s\vert }\) for \(x^{k}\vert x\vert ^{s}.)\) As Katz has recently recalled [19], Hecke’s technique for treating such expressions is to apply the Poisson summation formula to \(\sum _{\beta }\) with Re \(s\gg 0\), and then to show that one obtains the value of the resulting double series by setting \(s=0\) in each term.\par
\par
More precisely, for Re \(s\gg 0\), our series may be written\par
\par
\(\sum _{0\neq \lambda \in 𝔟𝔲^{-1}\,\operatorname{mod}\,U}\sum _{\gamma \in K}\) \(\widehat{\eta }_{𝔲}(\gamma )\int _{K\otimes \mathbb{R}}\) \(e^{2\pi i\cdot \operatorname{tr}(t\gamma )}dt/\mathcal{N}(\lambda \tau +t)^{k+\vert s\vert }\).\par
\par
[The summation modulo \(U\) continues to be well defined, as one can verify explicitly from the formula\par
\par
\DFormula{(6.4)\,\,\,\,\widehat{\eta }_{𝔲}(\gamma )=(\operatorname{sgn}\,\zeta )^{k}\eta _{𝔲}(\zeta \gamma )}
\par
for all units \(\zeta \) of \(K\). (Here we write sgn \(\zeta \) for the sign of \(\mathcal{N}\zeta .)\) The formula (6.4) is an immediate consequence of the parity of \(\varepsilon .]\)\par
\par
We consider the term in the inner sum with \(\gamma =0\):\par
\par
(6.5)    \(\widehat{\eta }_{𝔲}(0)\sum _{0\neq \lambda \in 𝔟𝔲^{-1}\,\operatorname{mod}\,U}\int _{K\otimes \mathbb{R}}\) \(dt/\mathcal{N}(\lambda \tau +t)^{k+\vert s\vert }\).\par
\par
By (6.4), we have \(\widehat{\eta }_{𝔲}(0)=0\) if there is a unit \(\zeta \) with (sgn \(\zeta )^{k}=-1\). We assume from now on, therefore, that this is not the case.\par
\par
In the integral, we write \(y\) for the totally positive number Im \(\tau \in (K\otimes \mathbb{R})^{\ast}\).\par
\DAnchor{45}{271}
\subsection*{Then the integral becomes}
\par
\DFormula{[1/(\mathcal{N}\lambda ^{k}\cdot \mathcal{N}y^{k}\cdot \vert \lambda \vert ^{s-1}\vert y\vert ^{s-1})]\cdot \varphi _{k}(s)^{r},}
\par
where\par
\par
\(\varphi _{k}(s)=\int _{-\infty }^{+\infty }\) \(dt/((i+t)^{k}\vert i+t\vert ^{s})\).\par
\par
We have, in fact \(\vert y\vert =\mathcal{N}y\) because \(y\) is totally positive). Then (6.5) may be rewritten\par
\par
\DFormula{(6.5\,bis)\,\,\,\,[\widehat{\eta }_{𝔲}(0)/(\mathcal{N}y^{k-1}\vert y\vert ^{s})]\varphi _{k}(s)^{r}\cdot \sum _{0\neq \lambda \in 𝔟𝔲^{-1}\,\operatorname{mod}\,U}\,(\operatorname{sgn}\,\lambda )^{k}/\vert \lambda \vert ^{k+s-1}.}
\par
By writing the sum in terms of partial zeta functions (for example), we see that it represents a meromorphic function of \(s\) with a possible simple pole at \(s=2-k\). Thus if \(k>2\), (6.5) represents a function of \(s\) which is holomorphic at least in the half-plane \(\operatorname{Re}(s)>-1\). Further, its value at \(s=0\) is then zero, because we easily compute \(\varphi _{k}(0)=0\).\par
\par
Now Hecke [9, \(p\). 393] yields the expressions\par
\par
\DFormula{\varphi _{2}(s)=-s\sqrt{\pi }\,\Gamma ((s+1)/2)/[\Gamma (s/2+1)(s+2)],}
\par
\DFormula{\varphi _{1}(s)=\sqrt{\pi }\,\Gamma ((s+1)/2)/[i\Gamma (s/2+1)].}
\par
For \(k=2\), we thus find again that (6.5) is holomorphic for \(\operatorname{Re}(s)>-1/2\) and has the value 0 at \(s=0\). (Because \(r>1\), \(\varphi _{2}(s)^{r}\) has at least a double zero at \(s=2-k=0.)\)\par
\par
For \(k=1\), the possible pole at \(s=1\) of the series\par
\par
\DFormula{(6.6)\,\,\,\,\sum _{0\neq \lambda \in 𝔟𝔲^{-1}\,\operatorname{mod}\,U}\,\operatorname{sgn}\,\lambda /\vert \lambda \vert ^{s}}
\par
does not appear. Indeed, (6.6) is, up to the factor \(\mathcal{N}(𝔟𝔲^{-1})^{-s}\), the L-series attached to an odd function on the group of strict ideal classes of \(K\). Hence (6.5) is again holomorphic for \(\operatorname{Re}(s)>-1/2\), and its value at \(s=0\) is\par
\par
\DFormula{\widehat{\eta }_{𝔲}(0)(-\pi i)^{r}[\sum _{\lambda \in 𝔟𝔲^{-1},\,\lambda \neq 0\,\operatorname{mod}\,U}\,\operatorname{sgn}\,\lambda /\vert \lambda \vert ^{s}]_{s=0}.}
\DAnchor{46}{272}
\subsection*{Thus the term we are considering becomes}
\par
\DFormula{(6.7)\,\,\,\,(-2)^{-r}\sum _{𝔲}\,\mathcal{N}𝔲^{-1}\widehat{\eta }_{𝔲}(0)[\sum _{\lambda \in 𝔟𝔲^{-1},\,\lambda \neq 0\,\operatorname{mod}\,U}\,\operatorname{sgn}\,\lambda /\mathcal{N}(\lambda 𝔲)^{s}]_{s=0}.}
\par
For \(\lambda \in K\) we have\par
\par
\DFormula{T\eta (\lambda ^{-1}𝔲^{-1}\cdot 0)=(\operatorname{sgn}\,\lambda )\mathcal{N}𝔲^{-1}\widehat{\eta }_{𝔲}(0).}
\par
Hence (6.7) may be rewritten as the L-value which is formally\par
\par
\DFormula{(-2)^{-r}\sum _{𝔵\subset 𝔟}\,T\eta (𝔵^{-1}\cdot 0).}
\par
Now \(\eta (x)=T\varepsilon (x\cdot \alpha )\). Hence\par
\par
\DFormula{T\eta (𝔵)=\varepsilon (𝔵\cdot \alpha ^{-1})(-1)^{r}\mathcal{N}𝔞.}
\par
Thus (6.7) becomes\par
\par
\DFormula{\mathcal{N}𝔞2^{-r}\sum _{𝔵\subset 𝔟}\,\varepsilon (𝔵^{-1}\cdot 0)=\mathcal{N}𝔞2^{-r}\sum _{𝔵\subset 𝔅𝔞}\,\varepsilon (𝔵^{-1}\cdot 0).}
\par
Since \(c=(𝔞\cdot \alpha ^{-1})\), this becomes\par
\par
\DFormula{2^{-r}\mathcal{N}𝔞\sum _{𝔵\subset 𝔅}\,\varepsilon ((c𝔵)^{-1}\cdot 0)=2^{-r}\mathcal{N}𝔞L(0,\widetilde{\varepsilon }_{c}).}
\par
If we now take into account the factor \([(-1)^{k}(k-1)!]^{r}\) in \(G_{k,\varepsilon }\), we finally arrive at the conclusion that the term we are calculating contributes to the q-exapansion (6.2) a term\par
\par
\DFormula{(-1)^{r}2^{-r}\mathcal{N}𝔞L(0,\widetilde{\varepsilon }_{c})}
\par
when \(k=1\). Under the assumption made above that all units of \(K\) have norm 1, we can ignore the \((-1)^{r}\), and we obtain precisely the perturbing term in (6.2) for \(k=1\). If there is a unit of \(K\) with norm \(-1\), then as we have already noted, (6.5) vanishes. On the other hand, \(\widetilde{\varepsilon }=0\) for odd \(\varepsilon \) if there is a unit of \(K\) of negative norm. Hence the perturbing term of (6.4) vanishes as well.\par
\par
We must finally consider the behavior for \(s\to 0\) of the sum\par
\par
(6.8)    \(\sum _{0\neq \lambda \in 𝔟𝔲^{-1}\,\operatorname{mod}\,U}\sum _{\gamma \in K^{\ast}}\) \(\widehat{\eta }_{𝔲}(\gamma )\int _{K\otimes \mathbb{R}}\) \(e^{2\pi i\cdot \operatorname{tr}(t\gamma )}dt/\mathcal{N}(\lambda \tau +t)^{k+\vert s\vert }\).\par
\par
As Katz explains, it suffices to continue each integral as a holomorphic function of \(s\), set \(s=0\), and evaluate the resulting double sum. Referring to Katz’s calculation ([19], 3.2.32), we find that\par
\par
\([\int _{K\otimes \mathbb{R}}\) \(e^{2\pi i\cdot \operatorname{tr}(t\gamma )}dt/\mathcal{N}(\lambda \tau +t)^{k+\vert s\vert }]_{s=0}\)\par
\par
\DFormula{=\begin{cases}\left[\dfrac{(2\pi i)^k}{(k-1)!}\right]^r e^{-2\pi i\cdot\operatorname{tr}(\gamma\lambda\tau)}\mathcal N\gamma^k\lvert\gamma\rvert^{-1},&\text{if }\gamma\lambda\ll0;\\0,&\text{otherwise}.\end{cases}}
\DAnchor{47}{273}
The inner sum of (6.8) may thus be written\par
\par
\DFormula{[((2\pi i)^{k})/(k-1)!]^{r}\,\sum _{\beta \in K,\,\beta \lambda \gg 0}\,\widehat{\varepsilon }_{𝔲}(\beta )\mathcal{N}\beta ^{k}\vert \beta \vert ^{-1}q^{\beta \lambda },}
\par
because on the substitution \(\beta \mapsto -\beta \) the functions \(\widehat{\varepsilon }_{𝔲}(\beta )\) and \(\mathcal{N}\beta ^{k}\) both take on factors \((-1)^{rk}\). Thus (6.8) gives to \(S\) the contribution\par
\par
\DFormula{(6.9)\,\,\,\,1/[(k-1)!]^{r}\,\sum _{𝔲}\,(\sum _{0\neq \lambda \in 𝔟𝔲^{-1}\,\operatorname{mod}\,U}\sum _{\beta \in K,\,\beta \lambda \gg 0}[\widehat{\eta }_{𝔲}(\beta )(\operatorname{sgn}\,\beta )^{k}\mathcal{N}𝔲^{-1}]\mathcal{N}(\beta 𝔲^{-1})^{k-1}q^{\beta \lambda }).}
\par
We recognise the bracketed term as \((T\eta )(\beta 𝔲^{-1})\). Setting \(\mu =\beta \lambda \), we change the order of summation in (6.9) and rewrite it as\par
\par
\DFormula{1/[(k-1)!]^{r}\,\sum _{\mu \gg 0,\,\mu \in K}\sum _{𝔲}\sum _{\lambda \in 𝔟𝔲^{-1}\,\operatorname{mod}\,U}(T\eta )(𝔲^{-1}\lambda ^{-1}\mu )\mathcal{N}(𝔲^{-1}\lambda ^{-1}\mu )^{k-1}q^{\mu }}
\par
\DFormula{=1/[(k-1)!]^{r}\,\sum _{\mu \gg 0}\sum _{𝔵\subset 𝔟}(T\eta )(𝔵^{-1}\mu )\mathcal{N}(𝔵^{-1}\mu )^{k-1}q^{\mu }.}
\par
We have \((T\eta )(𝔵)=\varepsilon (𝔵\cdot \alpha ^{-1})\mathcal{N}𝔞\cdot (-1)^{rk}\). Multiplying by the factor \((-1)^{rk}[(k-1)!]^{r}\) which appears in \(G_{k,\varepsilon }\), we have then\par
\par
\DFormula{\mathcal{N}𝔞\sum _{\mu \gg 0}\sum _{𝔵\subset 𝔅𝔞}\,\varepsilon (𝔵^{-1}\mu \cdot \alpha ^{-1})\mathcal{N}(𝔵^{-1}\mu )^{k-1}q^{\mu }}
\par
\DFormula{=\mathcal{N}𝔞^{k}\sum _{\mu \gg 0}\sum _{𝔵\subset 𝔅𝔞^{2}}\,\varepsilon (𝔵^{-1}\mu 𝔞\cdot \alpha ^{-1})\mathcal{N}(𝔵^{-1}\mu )^{k-1}q^{\mu }}
\par
\DFormula{=\mathcal{N}𝔞^{k}\sum _{\mu \gg 0}(\sum _{𝔵\subset 𝔅𝔞^{2}}\,\varepsilon _{c}(𝔵^{-1}\mu )\mathcal{N}(𝔵^{-1}\mu )^{k-1})q^{\mu }.}
\par
The proof of (6.1) is complete.\par
\par
\subsection*{§7. Theta Series}
\par
In this §, we study certain series of weight one, with the polarization module \(𝔠\) of §5 taken to be the ring of integers of \(K\).\par
\par
Let u be a totally positive element of \(K\), and let \(L=K(\sqrt{-u})\) be the corresponding totally imaginary quadratic extension of \(K\). Let\par
\par
\DFormula{\omega :G\to \lbrace \pm 1\rbrace }
\par
be the character of \(G=\operatorname{Gal}(K^{ab}/K)\) corresponding to L/K. By composing \(\omega \) with the “Artin” map j of (2.23) we obtain a character\par
\par
\DFormula{\widehat{K}^{\ast}\to \lbrace \pm 1\rbrace ,}
\par
again denoted by \(\omega \).\par
\DAnchor{48}{274}
(7.1) Theorem. For suitable \(N\geq 1\), there exists a weight-one modular form \(F\) over \(\mathbb{C}\) on \(\Gamma _{00}(N)\) whose \(q\)-expansion at the cusp determined by each \(\alpha \in \widehat{K}^{\ast}\) is the series\par
\par
\(\mathcal{N}𝔞\cdot \omega (\alpha )\sum _{x,y\in 𝔞}\) \(q^{x^{2}+uy^{2}}\),\par
\par
where \(𝔞=(\alpha )\) is the ideal generated by \(\alpha \).\par
\par
Remark. It is well known that the series \(\sum q^{x^{2}+uy^{2}}\) is the standard \(q\)-expansion of a weight-one modular form, cf. e.g. [7]. Our interest is in calculating the “other” \(q\)-expansions of this form. We do this by analyzing the action (by “slashing”) of \(\operatorname{SL}(2,\widehat{K})\) on this form, in terms of the Weil representation.\par
\par
Proof of (7.1). Let \(\psi :A\to \mathbb{C}^{\ast}\) be the additive character of §3. We recall that the choice of \(\psi \) as an additive character on \(A\), trivial of \(K\), determines a certain representation, the Weil representation, of \(\operatorname{SL}(2,A)\) on the Schwartz space \(\mathcal{S}(L_{A})\) of the adelization \(L_{A}=L\otimes _{K}\) \(A\) of \(L\). (See [14, 35, 41].) We denote this representation by \(r\). This representation is derived from local Weil representations \(r_{v}\) of the various groups \(\operatorname{SL}(2,K_{v})\) on the corresponding Schwartz spaces \(\mathcal{S}(L_{v})\). This local representations may in turn be described by certain explicit formulas giving the actions of\par
\par
\DFormula{\begin{pmatrix}0&1\\-1&0\end{pmatrix}}
\par
and of matrices of the form\par
\par
\DFormula{\begin{pmatrix}1&x\\0&1\end{pmatrix},\,\,\,\,\begin{pmatrix}a&0\\0&a^{-1}\end{pmatrix}}
\par
cf. [35].\par
\par
For \(M\in \mathcal{S}(L_{A})\) one defines\par
\par
\DFormula{\Theta (M)=\sum _{x\in L}M(x).}
\par
This function is left invariant under \(\operatorname{SL}(2,K)\): for \(s\in \operatorname{SL}(2,K)\) we have\par
\par
\DFormula{\Theta (r(s)M)=\Theta (M).}
\par
For fixed \(M\), the function \(s\mapsto \Theta (r(s)M)\) is then a continuous function on \(\operatorname{SL}(2,A)\) which is invariant on the left by \(\operatorname{SL}(2,K)\) [35, §2].\par
\par
We shall make a particular choice for the infinite component of \(M\) (cf. [8, 2.35]). Let n denote the norm from \(L\) to \(K\). Let \(L_{\infty }=L\otimes _{\mathbb{Q}}\) \(\mathbb{R}\). For \(x\in L_{\infty }\), set\par
\par
\DFormula{M_{\infty }(x)=\operatorname{exp}\lbrace -2\pi \cdot \operatorname{tr}_{K\otimes \mathbb{R}/\mathbb{R}}(n(x))\rbrace .}
\par
Then \(M_{\infty }\) is the product of functions \(M_{v}\) on the completions \(L_{v}=L\otimes _{K}\) \(\mathbb{R}\) of \(L\) at the infinite primes of \(K\) (we use \(v:K\hookrightarrow \mathbb{R}\) to make the tensor product), each given by the same formula\par
\par
\DFormula{M_{v}(x)=e^{-2\pi n(x)}.}
\par
If we view \(L_{v}\) as the complex field \(\mathbb{C}\), the function \(M_{v}\) is simply\par
\par
\DFormula{z\mapsto e^{-2\pi \vert z\vert ^{2}}.}
\par
For \(\theta \in \mathbb{R}\), let \(\rho (\theta )=\begin{pmatrix}\cos\theta&-\sin\theta\\\sin\theta&\cos\theta\end{pmatrix}\in \operatorname{SL}(2,\mathbb{R})\). If v is an infinite place of \(K\), we write \(\rho _{v}(\theta )\) for \(\rho (\theta )\) if we think of \(\operatorname{SL}(2,\mathbb{R})\) as \(\operatorname{SL}(2,K_{v})\). One has the formula\par
\DAnchor{49}{275}
\DFormula{(7.2)\,\,\,\,r_{v}(\rho _{v}(\theta ))\cdot M_{v}=e^{-i\theta }M_{v},}
\par
cf. [36, Lemma 1.2].\par
\par
(7.3) On the other hand, let \(a,b\in \mathbb{R}\), \(b>0\). The explicit formulas which define \(r_{v}\) show that\par
\par
\DFormula{r_{v}(\begin{pmatrix}\sqrt b&a\sqrt{b^{-1}}\\0&\sqrt{b^{-1}}\end{pmatrix})\cdot M_{v}}
\par
is the function\par
\par
\(x\mapsto \sqrt{b}\cdot e^{2\pi i((a+bi)n(x))}\).\par
\par
For \(\tau \) in the upper half plane\par
\par
\DFormula{\mathfrak{H}=\lbrace \tau \in K\otimes \mathbb{C}\,\vert \,\operatorname{Im}\,\tau \gg 0\rbrace ,}
\par
and for \(g=\begin{pmatrix}a&b\\c&d\end{pmatrix}\in \operatorname{SL}(2,K\otimes \mathbb{R})\), we write \(g\cdot \tau \) for the quotient\par
\par
\DFormula{(a\tau +b)/(c\tau +d)\in \mathfrak{H}.}
\par
Writing \(\mathcal{N}\) for the norm map \(K\otimes \mathbb{C}\to \mathbb{C}\), we let \(j(g,\tau )\) be the “factor of automorphy”\par
\par
\DFormula{\mathcal{N}(c\tau +d).}
\par
\subsection*{One verifies easily the formula}
\par
(7.4)    \(j(hg,\tau )=j(h,g\cdot \tau )\cdot j(g,\tau )\)\par
\par
\DFormula{\text{for}\,g,h\in \operatorname{SL}(2,K\otimes \mathbb{R}).}
\par
We now consider the finite part \(L_{f}=L\otimes _{K}\) \(A_{f}\) of \(L_{A}\), and its Schwartz space \(\mathcal{S}(L_{f})\). \(A\) “finite” Weil representation gives an action of \(\operatorname{SL}(2,A_{f})\) on \(\mathcal{S}(L_{f})\), which we again denote \(r\). The “multiplication” map\par
\par
\DFormula{\operatorname{SL}(2,A_{f})\times \mathcal{S}(L_{f})\to \mathcal{S}(L_{f}),}
\par
is continuous, so that the stabilizer in \(\operatorname{SL}(2,A_{f})\) of each \(\varepsilon \in \mathcal{S}(L_{f})\) is open. For \(\varepsilon \in \mathcal{S}(L_{f})\), we let \(M_{\varepsilon }=\varepsilon \otimes M_{\infty }\) be the product of \(\varepsilon \) with the function \(M_{\infty }\) defined above. Let\par
\par
\DFormula{\varphi _{\varepsilon }:\operatorname{SL}(2,A)\to \mathbb{C}}
\par
be the corresponding function \(s\mapsto \Theta (r(s)\cdot M_{\varepsilon })\) on \(\operatorname{SL}(2,A)\).\par
\par
For \(g\in \operatorname{SL}(2,K\otimes \mathbb{R})\), we write simply \(\varphi _{\varepsilon }(g)\) for the value of \(\varphi _{\varepsilon }\) on the image \(1\times g\) of g in \(\operatorname{SL}(2,A)\).\par
\par
For each \(\varepsilon \), we define a function \(F_{\varepsilon }\) on \(\mathfrak{H}\) by the series\par
\par
\DFormula{\sum _{x\in L}\,\varepsilon (x)q^{n(x)}.}
\par
(7.5) Proposition. We have \(F_{\varepsilon }(\tau )=\varphi _{\varepsilon }(g)j(g,i)\) for all \(g\in \operatorname{SL}(2,K\otimes \mathbb{R})\) such that \(g\cdot i=\tau \).\par
\DAnchor{50}{276}
Proof. If g is an element such that \(g\cdot i=\tau \), then any other element is obtained by multiplying g on the right by rotations \(\rho _{v}(\theta )\in \operatorname{SL}(2,K_{v})\hookrightarrow \operatorname{SL}(2,K\otimes \mathbb{R})\). Since \(j(\rho _{v}(\theta ),i)=e^{i\theta }\), we see by (7.2) and by (7.4) that the assertion is true for one g if and only if it is true for each g. We write \(\tau =a+bi\), with \(a,b\in K\otimes \mathbb{R}\). Since \(\tau \in \mathfrak{H}\), b is totally positive. Thus b has a unique totally positive square root \(\sqrt{b}\). It is clear that the matrix\par
\par
\DFormula{g=\begin{pmatrix}\sqrt b&a\sqrt{b^{-1}}\\0&\sqrt{b^{-1}}\end{pmatrix}}
\par
is such that \(g\cdot i=\tau \). For this choice of g however, the formula \(F_{\varepsilon }(\tau )=\varphi _{\varepsilon }(g)j(g,i)\) is an immediate consequence of the definition of \(\varphi _{\varepsilon }(g)\) and (7.3).\par
\par
We now recall the “slashing” operator (with the weight \(k\) taken to be 1)\par
\par
\DFormula{(F\vert h)(\tau )=F(h\cdot \tau )j(h,\tau )^{-1},\,\,\,\,h\in \operatorname{SL}(2,K\otimes \mathbb{R}).}
\par
(7.6) Proposition. For \(h\in \operatorname{SL}(2,K\otimes \mathbb{R})\), we have\par
\par
\((F_{\varepsilon }\vert h)(\tau )=\varphi _{\varepsilon }(hg)j(g\cdot i)\)\par
\par
for each g such that \(g\cdot i=\tau \).\par
\par
Proof. This formula follows immediately from (7.4) and (7.5).\par
\par
For \(A\in \operatorname{SL}(2,A)\), we let \(A_{f}\) (resp. \(A_{\infty })\) denote the image of \(A\) in \(\operatorname{SL}(2,A_{f})\) (resp. \(\operatorname{SL}(2,K\otimes \mathbb{R}))\). Regarding \(\operatorname{SL}(2,A_{f})\) and \(\operatorname{SL}(2,K\otimes \mathbb{R})\) as subgroups of \(\operatorname{SL}(2,A)\), we have \(A=A_{f}\cdot A_{\infty }\).\par
\par
(7.7) Theorem. For all \(A\in \operatorname{SL}(2,K)\), we have\par
\par
\DFormula{F_{\varepsilon }\vert A=F_{r(A_{f})^{-1}\cdot \varepsilon }.}
\par
Proof. Let \(\tau \in \mathfrak{H}\), and choose \(g\in \operatorname{SL}(2,K\otimes \mathbb{R})\) so that \(g\cdot i=\tau \). Since \(\vert \) represents the action of a matrix in \(\operatorname{SL}(2,K\otimes \mathbb{R})\), it is \((F_{\varepsilon }\vert A_{\infty })(\tau )\) which we must compute. By (7.6), this number is\par
\par
\DFormula{\varphi _{\varepsilon }(A_{\infty }g)j(g,i)=\varphi _{\varepsilon }(A_{f}^{-1}g)j(g,i),}
\par
with the equality resulting from the left-invariance of \(\varphi _{\varepsilon }\) under \(\operatorname{SL}(2,K)\). Since \(g\in \operatorname{SL}(2,K\otimes \mathbb{R})\), \(A_{f}\) and g commute in \(\operatorname{SL}(2,A)\). Commuting the two, and recalling the definition of \(\varphi _{\varepsilon }\), we find that\par
\par
\DFormula{\varphi _{\varepsilon }(A_{f}^{-1}g)=\varphi _{r(A_{f}^{-1})\cdot \varepsilon }(g).}
\par
The asserted equality now follows from (7.5).\par
\par
Since the stabilizer of \(\varepsilon \) in \(\operatorname{SL}(2,A_{f})\) is open, (7.7) implies that \(F_{\varepsilon }\) is invariant under some congruence subgroup of \(\operatorname{SL}(2,\mathcal{O})\). Hence, by strong approximation, \(F_{\varepsilon }\vert A\) may be defined for all \(A\in \operatorname{SL}(2,A_{f})\). By continuity, we see that the definition thus obtained is the obvious one: for all \(A\in \operatorname{SL}(2,A_{f})\), we have\par
\par
\DFormula{(7.8)\,\,\,\,F_{\varepsilon }\vert A=F_{r(A^{-1})\cdot \varepsilon }.}
\par
We now make a specific choice of \(\varepsilon \). Let \(R\subset L\) be the lattice \(\mathcal{O}+\mathcal{O}u\), and let \(\varepsilon \) be the characteristic function of the subset \(R\otimes _{\mathbb{Z}}\) \(\widehat{\mathbb{Z}}=\widehat{O}+\widehat{O}u\) of \(L_{f}\). Then\par
\par
\(F_{\varepsilon }=\sum _{x,y\in \mathcal{O}}\) \(q^{x^{2}+uy^{2}}\).\par
\DAnchor{51}{277}
To show that the holomorphic function \(F=F_{\varepsilon }\) is (the \(q\)-expansion of) a modular form of weight 1 on \(\Gamma _{00}(N)\) over \(\mathbb{C}\), we must show that \(F\) is invariant under the group of matrices \(\begin{pmatrix}a&b\\c&d\end{pmatrix}\) with \(a,d\in 1+N\mathcal{O}\), \(b\in 𝔇^{-1}\), \(c\in N𝔇\). Any such matrix is the product of a transvection \(\begin{pmatrix}1&x\\0&1\end{pmatrix}\) \((x\in 𝔇^{-1})\) and a matrix which satisfies \(b\in N𝔇^{-1}\) in addition to the other requirements. Since \(F\) is clearly invariant under the transvections, it is enough to show that \(F\) is invariant under\par
\par
\DFormula{\Gamma (N)=\lbrace \begin{pmatrix}a&b\\c&d\end{pmatrix}\,\vert \,a,d\in 1+N\mathcal{O},\,b\in N𝔇^{-1},\,c\in N𝔇\rbrace .}
\par
What is required, for (7.1), is merely that \(F\) be invariant under \(\Gamma (N)\) for some \(N\). But this is simply a restatement of the fact that \(F_{\varepsilon }\) is invariant under some congruence subgroup of \(\operatorname{SL}(2,\mathcal{O})\), which we have already established.\par
\par
To conclude the proof of (7.1), we note by (5.8) that the \(q\)-expansion of \(F\) at the cusp corresponding to \(\alpha \) is \(F\vert A\), where \(A=\begin{pmatrix}\alpha&0\\0&\alpha^{-1}\end{pmatrix}\in \operatorname{SL}(2,A_{f})\). By (7.8), we may rewrite this \(F_{r(A)^{-1}\cdot \varepsilon }\). Now \(A^{-1}\) is such that \(r(A^{-1})\) is given by a particularly simple formula [35, (2.1)]. We find that \(r(A)^{-1}\varepsilon \) is the function on \(L_{f}\)\par
\par
\DFormula{x\mapsto \omega (\alpha )^{-1}\vert \vert \alpha \vert \vert ^{-1}\varepsilon (\alpha ^{-1}x).}
\par
Now \(\omega (\alpha )^{-1}=\omega (\alpha )\), and \(\vert \vert \alpha \vert \vert ^{-1}=\mathcal{N}𝔞\). Finally, \(\varepsilon (\alpha ^{-1}x)\) is the characteristic function of \(\text{â}+\text{â}u\subset L_{f}\). Hence\par
\par
\(F_{\varepsilon (\alpha ^{-1}x)}=\sum _{x,y\in 𝔞}\) \(q^{x^{2}+uy^{2}}\),\par
\par
and\par
\par
\(F\vert A=\omega (\alpha )\mathcal{N}𝔞\sum _{x,y\in 𝔞}\) \(q^{x^{2}+uy^{2}}\),\par
\par
as desired.\par
\par
\subsection*{§8. Congruences for L-values}
\par
For \(\varepsilon :I\to \mathbb{C}\) a Schwartz function, we have defined \(L(s,\varepsilon )\) as the sum\par
\par
\DFormula{\sum _{𝔞\in I_{0}}\,\varepsilon (𝔞)\mathcal{N}𝔞^{-s}}
\par
in §3. By writing \(L(s,\varepsilon )\) in terms of partial zeta functions, we see that Siegel’s rationality theorem [37, 38] is equivalent to:\par
\par
(8.1) Theorem. If \(\varepsilon \) takes values in \(\mathbb{Q}\), then for each integer \(k\geq 1\) we have\par
\par
\DFormula{L(1-k,\varepsilon )\in \mathbb{Q}.}
\par
In the following discussion, we will admit (8.1) as known. The reader will observe that it is in fact a consequence of (0.1), in view of the calculations with Eisenstein series below.\par
\DAnchor{52}{278}
As in §0, Theorem (8.1) allows us to define, by linearity, values\par
\par
\DFormula{L(1-k,\varepsilon )\in V}
\par
for each \(\varepsilon \) with values in a \(\mathbb{Q}\)-vector space \(V\). We may view each association\par
\par
\DFormula{\varepsilon \mapsto L(1-k,\varepsilon )\,\,\,\,(k\geq 1)}
\par
as a \(\mathbb{Q}\)-valued distribution \(T_{k}\) on \(I\). Each \(T_{k}\) has an obvious homogeneity property with respect to substitutions\par
\par
\DFormula{\varepsilon (x)\mapsto \varepsilon (𝔞x)}
\par
with \(𝔞\) an ideal of \(K\). In particular, \(T_{1}\) is invariant under such substitutions.\par
\par
If \(V\) is now a \(\widehat{\mathbb{Q}}\)-module and \(c\in G\), we define\par
\par
\(\Delta _{c}(1-k,\varepsilon )=L(1-k,\varepsilon )-\mathcal{N}c^{k}L(1-k,\varepsilon _{c})\),\par
\par
where \(\varepsilon _{c}\) is the function \(x\mapsto \varepsilon (cx)\), cf. §0. For each c and \(k\), the map\par
\par
\DFormula{\mu _{c,k}:\,\varepsilon \mapsto \Delta _{c}(1-k,\varepsilon )}
\par
is a distribution on \(I\) with values in \(\widehat{\mathbb{Q}}\). Let us denote \(\mu _{c,1}\) simply by \(\mu _{c}\).\par
\par
(8.2) Main Theorem. For each \(c\in G\), \(\mu _{c}\) is a measure on \(I\) with values in \(\widehat{\mathbb{Z}}\). If \(\varepsilon \) is a Schwartz function on \(I\) and \(k\) a positive integer, we have\par
\par
\DFormula{\int \,\varepsilon \cdot \mathcal{N}^{k-1}\,d\mu _{c}=\Delta _{c}(1-k,\varepsilon ).}
\par
[Note that in the above equation, \(\mathcal{N}\) is the norm function \(I\to \widehat{\mathbb{Q}}.]\)\par
\par
The principal purpose of this § is to derive (8.2). Along the way, we will prove certain further assertions concerning the integrals against \(\mu _{c}\) of odd functions. These are summarized in (8.11), (8.12). Also, we wish to point out how (0.4) is a consequence of (8.2).\par
\par
(8.3) Proposition. Let \(\varepsilon \) be a Schwartz function on \(I\) with values in a \(\mathbb{Q}\)-vector space \(V\). Suppose that \(k\geq 1\) and that c is an element of the subgroup \(\Sigma ^{\pm }\) of \(G\), cf. (2.24). Then we have\par
\par
\(L(1-k,\varepsilon _{c})=\mathcal{N}c^{k}L(1-k,\varepsilon )\).\par
\par
Proof. We easily reduce to the case where \(\varepsilon \) is a complex-valued function which has parity \((a_{v})\) for some set of integers \(a_{v}=0,1\). For such \(\varepsilon \), the proposition amounts to the vanishing of \(L(1-k,\varepsilon )\) whenever we have \((-1)^{k}\neq (-1)^{a_{v}}\) for some v. This is the vanishing given by (3.24).\par
\par
(8.4) Theorem (The “Eisenstein congruences”). Let \(\varepsilon _{1},\ldots ,\varepsilon _{k},\ldots \) be Schwartz functions on \(I\) with values in \(\widehat{\mathbb{Q}}\). Let \(𝔅\) be an ideal of \(K\). Suppose that only finitely many of the \(\varepsilon _{k}\) are non-zero and that \(\varepsilon _{k}\) has parity \((-1)^{k}\) for each \(k\). Set\par
\par
\DFormula{\varphi =\sum _{k\geq 1}\,\varepsilon _{k}\mathcal{N}^{k-1}:\,I\to \widehat{\mathbb{Q}}.}
\DAnchor{53}{279}
Suppose for each totally positive element \(\mu \) of \(K\) that the finite sum\par
\par
\DFormula{\sum _{𝔞\subset 𝔅,\,𝔞\in I_{0}}\,\varphi (\mu 𝔞^{-1})}
\par
lies in \(\widehat{\mathbb{Z}}\). Then for each \(c\in G\), we have\par
\par
\DFormula{(8.5)\,\,\,\,\Delta _{c}(0,\widetilde{\varepsilon }_{1})+\sum _{k\geq 1}\Delta _{c}(1-k,\varepsilon _{k})\in 2^{r}\cdot \widehat{\mathbb{Z}},}
\par
where \(\widetilde{\varepsilon }_{1}\) is the modification of \(\varepsilon _{1}\) with respect to \(𝔅\) (§6).\par
\par
Proof. For a given collection of functions \((\varepsilon _{k})\), the hypothesis and the conclusion of the theorem depend only on the strict ideal class of \(𝔅\). Also, if \(𝔡\) is an ideal of \(K\), the theorem for \((\varepsilon _{k})\) and for \(𝔅\) is equivalent to the theorem for the ideal \(𝔅𝔡\) and the collection of functions\par
\par
\DFormula{x\mapsto \varepsilon _{k}(x𝔡)\mathcal{N}𝔡^{k-1}.}
\par
It thus suffices to treat the case where the \(\varepsilon _{k}\) are supported on \(A\). At the same time, we may as well assume that the ideal \(𝔅\) is equal to the ideal \(\mathcal{O}\). (This latter assumption will prove to be convenient because the modular forms in §0 were made with polarization module c implicitly taken to be \(\mathcal{O}.)\) Furthermore, we observe that it suffices to treat the case where the \(\varepsilon _{k}\) are \(\mathbb{Q}\)-valued. Indeed, it is clear that the hypothesis to (8.4) does not change if we replace the \(\varepsilon _{k}\) by functions which are congruent to them modulo \(\widehat{\mathbb{Z}}\). As for the conclusion of (8.4), we may specify in advance a positive integer \(D\) such that the conclusion is unchanged if we replace the \(\varepsilon _{k}\) by functions congruent to them mod \(D\) \(\widehat{\mathbb{Z}}\). (The point is that the L-values in question are a priori given as linear combinations of certain rational numbers --- values of partial zeta functions --- which may be specified as soon as one knows a common modulus of definition for the \(\varepsilon _{k}.)\)\par
\par
So let functions \(\varepsilon _{k}\) be given as in (8.4), under the further assumption that they are \(\mathbb{Q}\)-valued and supported on \(A\). Assume that the ideal \(𝔅\) is equal to \(\mathcal{O}\). For each \(k\geq 1\), let \(F_{k}\) be the Eisenstein series \(G_{k,\varepsilon }\) of (6.1). From (8.1) and equation (6.2), we see that the \(F_{k}\) each satisfy the rationality condition of (0.1). For each \(\alpha \in \widehat{K}^{\ast}\), let\par
\par
\DFormula{S(\alpha )=\sum _{k\geq 1}\,\mathcal{N}\alpha ^{-k}\cdot F_{k,\alpha };}
\par
this is a priori a formal series with coefficients in \(\widehat{\mathbb{Q}}\).\par
\par
Now the non-constant coefficients of S(1) are rational integers and thus in particular are in \(\widehat{\mathbb{Z}}\). Hence, by (0.3), applied for each prime \(p\), the difference between the constant coefficients of S(1) and \(S(\alpha )\) lies again in \(\widehat{\mathbb{Z}}\), for \(\alpha \in \widehat{K}^{\ast}\). Given \(\alpha \), set \(c=j(\alpha )=(\alpha )\cdot \alpha ^{-1}\). Looking at (6.2), and using the equation\par
\par
\DFormula{\mathcal{N}c=\mathcal{N}(\alpha )\mathcal{N}\alpha ^{-1},}
\par
we see that the conclusion of (8.4) holds for c. But j is surjective; hence the conclusion holds for each \(c\in G\).\par
\DAnchor{54}{280}
[In fact, it is clear that the conclusion of (8.4) depends a priori only on the image of c in \(G_{𝔣}\), where \(𝔣\) is a suitable conductor. Hence all that is needed is that \(G_{𝔣}\) is a quotient of \(\widehat{K}^{\ast}.]\)\par
\par
Although (8.4) and (8.2) are quite similar, there are two obstacles which prevent us from deducing (8.2) as an immediate consequence of (8.4). First of all, there is a parity assumption on the \(\varepsilon _{k}\) in (8.4). Secondly, (8.4) contains a perturbing term, involving the modification of \(\varepsilon _{1}\). These difficulties may account for the complexity of the calculations that follow.\par
\par
We shall suppose below without comment that \(\varepsilon \) is a Schwartz function on \(I\) with values in \(\widehat{\mathbb{Z}}\).\par
\par
(8.6) Proposition. If \(K\) has some unit of norm \(-1\) (e.g., if \(r\) is odd), and if \(\varepsilon \) is an odd function, then \(\Delta _{c}(0,\varepsilon )\in 2^{r}\) \(\widehat{\mathbb{Z}}\).\par
\par
Proof. By (8.4), it suffices to show that \(\widetilde{\varepsilon }=0\). However, the unit hypothesis shows that there is an element \(\sigma \) of \(\Sigma ^{\pm }\), not in \(\Sigma \), whose image in \(G_{1}\) is trivial. It follows that the restriction to the set of elements of \(I\) of the form \((x\cdot 0)\) of any odd function on \(I\) is identically zero. Hence, by definition, \(\widetilde{\varepsilon }=0\) in this case.\par
\par
We now consider the case where no unit of \(K\) has norm \(-1\). In \(G_{1}\), the Frobenius group \(\Sigma _{1}^{\pm }\) contains as a subgroup of index 2 the group \(\Sigma _{1}\) generated by products of pairs of Frobenius elements. We let \(𝔛=G_{1}/\Sigma _{1}\) and let \(\sigma \) be the real Frobenius element of \(𝔛\). The order of \(\sigma \) is 2. (We may view \(𝔛\) as the Galois group over \(K\) of the largest abelian extension of \(K\) which is unramified at the finite places and either totally real or else a “CM field.” The fact that \(K\) has no unit of norm \(-1\) means that this extension is in fact a CM field. In this optic, \(\sigma \) is the canonical “complex conjugation” of the field.)\par
\par
(8.7) Lemma. Let \(\varepsilon \) be a function \(𝔛\to \widehat{\mathbb{Z}}\) which is odd \([\varepsilon (\sigma x)=-\varepsilon (x)]\) and which satisfies\par
\par
\DFormula{\sum _{x\,\operatorname{mod}\,\lbrace 1,\sigma \rbrace }\,\varepsilon (x)\equiv 0\,(\operatorname{mod}\,2).}
\par
Then \(\varepsilon \) is the sum of functions of the form \(h(x)-h(ax)\), where \(h:𝔛\to \widehat{\mathbb{Z}}\) is odd and where \(a\in 𝔛\).\par
\par
Proof. For each \(a\in 𝔛\), define an odd function \(h_{a}:𝔛\to \widehat{\mathbb{Z}}\) by:\par
\par
\(h_{a}(x)=\lbrace \,1\,\text{if}\,x=a;\,-1\,\text{if}\,x=\sigma a;\,0\,otherwise.\,\rbrace \)\par
\par
By subtracting from \(\varepsilon \) multiples of the functions \(h_{a}(x)-h_{a}(ax)\), we may assume that \(\varepsilon \) is supported on \(\lbrace 1,\sigma \rbrace \). The hypothesis on \(\varepsilon \) then means that \(\varepsilon \) takes even values. The equation\par
\par
\DFormula{\varepsilon (x)=\varepsilon (x)/2-\varepsilon (\sigma x)/2}
\par
exhibits \(\varepsilon \) as a function of the form \(h(x)-h(ax)\).\par
\par
Suppose now that \(\varphi :I\to \widehat{\mathbb{Z}}\) is a function which is odd (or, more generally, which has parity \((a_{v})\) for some collection of integers \(a_{v}=0\) or 1). We will define\par
\DAnchor{55}{281}
an invariant \(\delta (\varphi )\) in \(\mathbb{Z}/2\mathbb{Z}\). The restriction of \(\varphi \) to the set of elements of the form \((x\cdot 0)\) in \(I\) may be viewed as a function\par
\par
\DFormula{\tau :G_{1}\to \widehat{\mathbb{Z}}.}
\par
Composing \(\tau \) with the natural map \(\widehat{\mathbb{Z}}\to \mathbb{Z}/2\mathbb{Z}\), we obtain a function \(\overline{\tau }\) on \(G_{1}\) which is even, i.e., \(\Sigma _{1}^{\pm }\)-invariant. We set\par
\par
\DFormula{\delta (\varphi )=\sum \overline{\tau }(g),}
\par
the summation running over elements g of \(G_{1}\), taken modulo \(\Sigma _{1}^{\pm }\).\par
\par
(8.8) Proposition. Suppose that \(\varepsilon \) is an odd function and that \(\delta (\varepsilon )=0\). Then\par
\par
\DFormula{\Delta _{c}(0,\varepsilon )\in 2^{r}\,\widehat{\mathbb{Z}}.}
\par
Proof. Because of (8.6), we may assume that no unit of \(K\) has norm \(-1\). The function \(\tau :G_{1}\to \widehat{\mathbb{Z}}\) constructed as above for \(\varepsilon \) is \(\Sigma _{1}\)-invariant and consequently may be viewed as a function on \(𝔛\). By (8.7), we may find odd locally constant functions \(H:A\to \widehat{\mathbb{Z}}\) and ideals \(𝔞\) of \(K\) such that the difference \(\varepsilon '\) between \(\varepsilon \) and the sum of the functions\par
\par
\DFormula{H(x)-H(𝔞x)}
\par
vanishes on \(𝔛\). By (8.4), we get then\par
\par
\DFormula{\Delta _{c}(0,\varepsilon ')\in 2^{r}\,\widehat{\mathbb{Z}}.}
\par
On the other hand, we have \(\Delta _{c}(0,H(x))=\Delta _{c}(0,H(𝔞x))\) for each \(H\), so that\par
\par
\DFormula{\Delta _{c}(0,\varepsilon )=\Delta _{c}(0,\varepsilon ').}
\par
This proves what is wanted.\par
\par
(8.9) Corollary. For each odd \(\varepsilon \), we have \(\Delta _{c}(0,\varepsilon )\in 2^{r-1}\widehat{\mathbb{Z}}\).\par
\par
We now introduce the notion of an exceptional field. We say that \(K\) is exceptional if the following two conditions are satisfied:\par
\par
i) All units of \(K\) have norm \(+1\),\par
\par
ii) There are units of \(K\) of all signatures compatible with (i).\par
\par
When \(K\) is exceptional, we have \(\Sigma _{1}=\lbrace 1\rbrace \), and \(\Sigma _{1}^{\pm }\) consists of two elements. \(A\) field is exceptional if its strict Hilbert class field is a CM field.\par
\par
(8.10) Proposition. Suppose that \(K\) is not exceptional. Then for each odd \(\varepsilon \) we have\par
\par
\DFormula{\Delta _{c}(0,\varepsilon )\in 2^{r}\,\widehat{\mathbb{Z}}.}
\par
Proof. In view of (8.6), we may again suppose that \(K\) satisfies (i). The hypothesis that \(K\) is not exceptional then means that there is a principal ideal \((\alpha )\) generated by an integer \(\alpha \) of positive norm, whose strict ideal class is non-trivial. Let \(\eta \) be the locally constant compactly supported function on \(I\) whose value on an ideal \(𝔞\) is given by\par
\par
\DFormula{\eta(\mathfrak a)=\begin{cases}0&\text{if }\mathfrak a\not\subset(\alpha);\\\frac12\operatorname{sgn}\mathfrak a&\text{if }\mathfrak a\subset(\alpha).\end{cases}}
\DAnchor{56}{282}
Here\par
\par
\DFormula{\operatorname{sgn}\mathfrak a=\begin{cases}0&\text{if }\mathfrak a\text{ is non-principal};\\+1&\text{if }\mathfrak a=(\beta)\text{ with }\mathcal N\beta>0;\\-1&\text{if }\mathfrak a=(\beta)\text{ with }\mathcal N\beta<0.\end{cases}}
\par
Note that sgn is well defined precisely because there is no unit of negative norm in \(K\).\par
\par
Let us check that \(\eta \) satisfies the hypothesis to (8.4) with \(𝔅=\mathcal{O}\), i.e.,\par
\par
\DFormula{\sum _{𝔞\subset \mathcal{O}}\eta (\mu 𝔞^{-1})=\sum _{(\mu )\subset 𝔞\subset (\alpha )}\eta (𝔞)\in \widehat{\mathbb{Z}},}
\par
for all \(\mu \gg 0\). We observe that \(𝔞\mapsto \alpha 𝔞^{-1}\mu \) gives an involution on the set of \(𝔞\) in the second sum, and that \(\eta (𝔞)=\eta (\alpha 𝔞^{-1}\mu )\). Hence \(\eta \) does indeed satisfy the hypothesis if there are no fixed points of the involution on which \(\eta \) takes non-zero values. But if \(𝔞=\alpha 𝔞^{-1}\mu \) and if \(𝔞\) is principal, then \((\alpha \mu )\) is the square of a principal ideal and hence in the trivial strict ideal class of \(K\). This is contrary to the choice of \(\alpha \) and the fact that \(\mu \gg 0\). Hence \(\Delta _{c}(0,\eta )+\Delta _{c}(0,\widetilde{\eta })\in 2^{r}\) \(\widehat{\mathbb{Z}}\), by (8.4).\par
\par
However \(\widetilde{\eta }\) is obviously the “same” function as \(\eta \), except that its support is \(\mathcal{O}\) rather than \((\alpha )\). More precisely, we have \(\eta (x)=\widetilde{\eta }((\alpha )x)\). It follows that \(\Delta _{c}(0,\eta )=\Delta _{c}(0,\widetilde{\eta })\). Thus\par
\par
\DFormula{\Delta _{c}(0,2\widetilde{\eta })\in 2^{r}\,\widehat{\mathbb{Z}}.}
\par
The function \(\varepsilon =2\widetilde{\eta }\) is odd and \(\widehat{\mathbb{Z}}\)-valued (it is just the function sgn with support on \(A)\). It satisfies \(\delta (\varepsilon )=1\) because, on the \((x\cdot 0)\), \(\varepsilon \) vanishes on all elements except those with wide class 1; in the sum defining \(\delta (\varepsilon )\) there is only one term. This gives (8.10).\par
\par
Proof of (8.2). For \(c\in G\), we must first show that \(\mu _{c}\) is a measure. This means showing that \(\Delta _{c}(0,\varepsilon )\in \widehat{\mathbb{Z}}\) for all \(\varepsilon \). Given \(\varepsilon \), we may define an odd function \(\varepsilon ^{-}\) by the formula\par
\par
\DFormula{\varepsilon ^{-}(x)=2^{-r}\sum _{\sigma }\,\mathcal{N}\sigma \cdot \varepsilon (\sigma x)\in 2^{-r}\widehat{\mathbb{Z}},}
\par
in which the summation runs over the group \(\Sigma ^{\pm }\) (whose cardinality is \(2^{r})\). It is apparent from (8.3) that \(\Delta _{c}(0,\varepsilon )=\Delta _{c}(0,\varepsilon ^{-})\). By (8.8) and (8.10), it suffices to show that \(\delta (2^{r}\varepsilon ^{-})=0\) whenever \(K\) is exceptional.\par
\par
Under this hypothesis on \(K\), the restriction of \(\varepsilon \) to the \((x\cdot 0)\) is already invariant under \(\Sigma \), since the image of this group in \(G_{1}\) is trivial. Hence the restriction of \(\varepsilon ^{-}\) to the \((x\cdot 0)\) takes values in \(2^{-1}\widehat{\mathbb{Z}}\). Since \(r>1\), the restriction of \(2^{r}\varepsilon ^{-}\) to the \((x\cdot 0)\) thus takes even values. In particular, \(\delta (2^{r}\varepsilon ^{-})=0\), as required. Thus \(\mu _{c}\) is indeed a measure with values in \(\widehat{\mathbb{Z}}\).\par
\par
\subsection*{It remains now to show that}
\par
\DFormula{\int \varepsilon \cdot \mathcal{N}^{k-1}d\mu _{c}=\Delta _{c}(1-k,\varepsilon )}
\par
for \(k>1\), and all Schwartz functions \(\varepsilon \). We define \(\varepsilon ^{-}\) as above, and we similarly define \(\varepsilon ^{+}\) by omitting the factor \(\mathcal{N}\sigma \) in the sum defining \(\varepsilon ^{-}\). Let \(\varepsilon ^{\dagger}\) be either \(\varepsilon ^{-}\)\par
\DAnchor{57}{283}
or \(\varepsilon ^{+}\) according as \(k\) is odd or even, so that \(\varepsilon ^{\dagger}\) has the same parity as \(k\). By (8.3), we have\par
\par
\DFormula{\Delta _{c}(1-k,\varepsilon )=\Delta _{c}(1-k,\varepsilon ^{\dagger}),}
\par
and\par
\par
\DFormula{\varepsilon \cdot \mathcal{N}^{k-1}\mu _{c}=\varepsilon ^{\dagger}\cdot \mathcal{N}^{k-1}\mu _{c}.}
\par
Hence there is no loss of generality in supposing that \(\varepsilon \) has parity \((-1)^{k}\), which we now do. Making a change of variables, we shall suppose also that \(\varepsilon \) is supported on \(A\).\par
\par
Let m be a positive integer. Choose a locally constant function\par
\par
\DFormula{\eta :A\to \widehat{\mathbb{Z}}}
\par
which is odd and congruent modulo m to the norm function \(\mathcal{N}\) on \(A\). By (8.4), we obtain from\par
\par
\(\varepsilon \mathcal{N}^{k-1}\equiv \varepsilon \eta ^{k-1}\) (mod \(m\widehat{\mathbb{Z}})\)\par
\par
the congruence\par
\par
\(\Delta _{c}(1-k,\varepsilon )\equiv \Delta _{c}(0,\varepsilon \eta ^{k-1})+\Delta _{c}(0,\widetilde{\varepsilon \eta ^{k-1}})\) mod \(m\widehat{\mathbb{Z}}\).\par
\par
Now \(\mathcal{N}\) vanishes on the \((x\cdot 0)\), so that \(\eta \) takes values on these elements which are divisible by m. Hence \(\widetilde{\varepsilon \eta ^{k-1}}\equiv 0\) mod m, so that its integral against \(\mu _{c}\) is again divisible by m. Thus we have simply\par
\par
\DFormula{\Delta _{c}(1-k,\varepsilon )\equiv \Delta _{c}(0,\varepsilon \eta ^{k-1})\,\operatorname{mod}\,m.}
\par
On the other hand, because \(\mu _{c}\) is a measure, the right hand member of this congruence is congruent mod m to the integral against \(\mu _{c}\) of \(\varepsilon \mathcal{N}^{k-1}\). Hence the two numbers we wish to prove equal are congruent mod m. But m was an arbitrary positive integer, so that the proof of (8.2) is complete.\par
\par
We now summarize what we can say about the integral of an odd function. Let \(\varphi :I\to \widehat{\mathbb{Z}}\) be odd, continuous, and compactly supported.\par
\par
(8.11) Theorem. We have \(\int \varphi d\mu _{c}\in 2^{r-1}\widehat{\mathbb{Z}}\). Moreover, this integral lies in \(2^{r}\widehat{\mathbb{Z}}\) unless \(K\) is exceptional and \(\delta (\varphi )=1\).\par
\par
Proof. Choose an odd Schwartz function \(\varepsilon :I\to \widehat{\mathbb{Z}}\) which is congruent to \(\varphi \) mod \(2^{r}\). Then \(\delta (\varphi )=\delta (\varepsilon )\) and the integrals against \(\mu _{c}\) of \(\varepsilon \) and \(\varphi \) are congruent mod \(2^{r}\). Hence our assertions follow from (8.8), (8.9), and (8.10).\par
\par
Suppose now that \(K\) is exceptional. For \(c\in G\) and \(\varphi \) with \(\delta (\varphi )=1\), the quantity\par
\par
\DFormula{2^{1-r}\cdot \int \varphi d\mu _{c}\,(\operatorname{mod}\,2)}
\par
is an element \(\xi (c)\) of \(\mathbb{Z}/2\mathbb{Z}\) which, by (8.11), is independent of \(\varphi \). To compute it, we can take \(\varphi \) to be locally constant. The equation\par
\par
\(\Delta _{cc'}(0,\varphi )=\Delta _{c}(0,\varphi )+\Delta _{c'}(0,\mathcal{N}c\cdot \varphi _{c})\)\par
\DAnchor{58}{284}
then shows that \(\xi \) is a homomorphism \(G\to \mathbb{Z}/2\mathbb{Z}\). Its kernel corresponds, then, to some extension of \(K\) of degree 1 or 2 which, by (8.3) is easily seen to be totally real. After looking at a small number of numerical examples, Lenstra predicted that this extension would prove to be the field \(M\) obtained by adjoining to \(K\) the square roots of all positive units. (That this field is a quadratic extension follows from the fact that \(K\) is exceptional.)\par
\par
His prediction was correct.\par
\par
(8.12) Theorem. If \(K\) is exceptional and \(\delta (\varphi )=1\), then we have\par
\par
\(\int \varphi d\mu _{c}\in 2^{r}\widehat{\mathbb{Z}}\)\par
\par
if and only if the image of c in \(\operatorname{Gal}(M/K)\) is trivial.\par
\par
Proof. Let \(\varepsilon \) be the “sign” function with support on \(A\), as in the proof of (8.10). We have already noted in that proof that \(\varepsilon \) is an odd function with \(\delta (\varepsilon )=1\). Also, with \(𝔅=\mathcal{O}\), we have \(\varepsilon =\widetilde{\varepsilon }\).\par
\par
Let \(E\) be the Eisenstein series \(G_{1,\varepsilon }\), made with \(𝔅=\mathcal{O}\). Next, choose u to be a totally positive unit which is not a square in \(K\), so that\par
\par
\DFormula{M=K(\sqrt{u}).}
\par
Let \(F\) be the theta series of (7.1), made with this choice of u. The two forms \(F\) and \(E\) satisfy the rationality property of (0.1), as we see by looking at the \(q\)-expansions of \(F\) and \(E\) at the standard cusp. (For the constant term of \(F\), we need also (8.1).) An elementary argument shows that the non-constant terms of the standard \(q\)-expansions of \(E\) and \(1/2\) \(F\) are congruent modulo 2. In other words, the form \(E/2-F/4\) has a standard \(q\)-expansion whose non-constant terms all lie in \(\widehat{\mathbb{Z}}\). Arguing now as in the proof of (8.4), we find for all \(c\in G\) a congruence\par
\par
\DFormula{[2^{-r}\Delta _{c}(0,\varepsilon )+2^{-r}\Delta _{c}(0,\widetilde{\varepsilon })]/2\equiv [1-\mathcal{N}c\omega (c)]/4\,\operatorname{mod}\,\widehat{\mathbb{Z}},}
\par
with \(\omega \) as in (7.1). On the other hand, we have already remarked that \(\varepsilon \) and \(\widetilde{\varepsilon }\) are equal. Hence the integral in (8.12) lies in \(2^{r}\widehat{\mathbb{Z}}\) if and only if\par
\par
\DFormula{\mathcal{N}c\omega (c)\equiv 1\,(\operatorname{mod}\,4).}
\par
Now, mod 4, \(\mathcal{N}\) is the quadratic character of \(G\) corresponding to the extension\par
\par
\DFormula{K(\sqrt{-1})/K,}
\par
whereas \(\omega \) corresponds to the quadratic extension\par
\par
\DFormula{K(\sqrt{-u})/K.}
\par
Hence the product of these two quadratic characters corresponds to the extension M/K, as desired.\par
\DAnchor{59}{285}
Proof of (0.4). Let \((\varepsilon _{k})\) be a sequence of functions as in (0.4). We may view these as functions on \(I\), defined mod \(𝔣\) and supported on \(A\). From this point of view, the \(\varepsilon _{k}\) furthermore vanish on integral ideals which are not prime to \(𝔣\). Hence if we form the sum \(\sum \varepsilon _{k}\mathcal{N}^{k-1}\) as a function \(I\to \mathbb{Q}_{p}\), the hypothesis to (0.4) implies that this function is in fact \(\mathbb{Z}_{p}\)-valued. (Note that \(\mathcal{N}\) is a priori \(\widehat{\mathbb{Z}}\)-valued, while \(\varepsilon _{k}\) is \(\mathbb{Q}_{p}\)-valued, so we may consider the product \(\varepsilon _{k}\mathcal{N}^{k-1}\) as a number in \(\mathbb{Q}_{p}.)\) For each \(c\in G\), the integral against \(\mu _{c}\) of this sum is again an element of \(\mathbb{Z}_{p}\). Using the second part of (8.2), we get then\par
\par
\DFormula{\sum _{k\geq 1}\Delta _{c}(1-k,\varepsilon _{k})\in \mathbb{Z}_{p},}
\par
as needed for (0.4).\par
\par
We must observe that the group called \(G\) in §0 is not the full group \(\operatorname{Gal}(K^{ab}/K)\) as in the succeeding §§, but rather a quotient of it. The explanation is that for \(\varepsilon \) a function mod \(𝔣\) with values in \(\mathbb{Q}_{p}\), the quantities \(\Delta _{c}(1-k,\varepsilon )\in \mathbb{Q}_{p}\) depend only on the image of c in the quotient group used in §0.\par
\par
References\par
\par
1. Barsky, \(D\).: Fonctions zêta \(p\)-adiques d’une classe de rayon des corps totalement réels. Groupe d’étude d’analyse ultramétrique 1977--78; errata, 1978--79\par
\par
2. Cassou-Noguès, \(P\).: Valeurs aux entiers négatifs des fonctions zêta et fonctions zêta \(p\)-adiques. Invent. Math. 51, 29--59 (1979)\par
\par
3. Coates, \(J\).: \(p\)-adic L-functions and Iwasawa’s theory. In: Algebraic Number Fields (A. Fröhlich, ed.). London-New York-San Francisco: Academic Press 1977\par
\par
4. Deligne, \(P\).: Variétés abéliennes ordinaires sur un corps fini. Invent. Math. 8, 238--243 (1969)\par
\par
5. Deligne, \(P\).: Letters to J-P. Serre: October, 1972 and December, 1973\par
\par
6. Deligne, \(P\)., Rapoport, \(M\).: Les schémas de modules de courbes elliptiques. Lecture Notes in Math. 349, 143--316 (1973)\par
\par
7. Eichler, \(M\).: On theta functions of real algebraic number fields. Acta Arith. 33, 269--292 (1977)\par
\par
8. Gelbart, \(S\).: Weil’s representation and the spectrum of the metaplectic group. Lecture Notes in Math. 530, 1976\par
\par
9. Hecke, \(E\).: Mathematische Werke. Göttingen: Vandenhoeck und Ruprecht 1959\par
\par
10. Hirzebruch, \(F\).: Hilbert modular surfaces. Enseignement Math. 19, 183--281 (1973)\par
\par
11. Hurwitz, \(A\).: Mathematische Werke. Basel: Emil Birkhäuser 1932\par
\par
12. Igusa, \(J\).: On the algebraic theory of elliptic modular functions. \(J\). Math. Soc. Japan 20, 96--106 (1968)\par
\par
13. Iwasawa, \(K\).: Lectures on \(p\)-adic L-functions. Princeton: Princeton University Press 1972\par
\par
14. Jacquet, \(H\)., Langlands, R.P.: Automorphic forms on \(\operatorname{GL}(2)\). Lectures Notes in Math 114 (1970)\par
\par
15. Katz, \(N\).: \(p\)-adic properties of modular schemes and modular forms. Lecture Notes in Math. 350, 70--190 (1973)\par
\par
16. Katz, \(N\).: The Eisenstein measure and \(p\)-adic interpolation. Amer. \(J\). Math. 99, 238--311 (1977)\par
\par
17. Katz, \(N\).: \(p\)-adic L-functions via moduli of elliptic curves. Proc. Symp. Pure Math. 29, 479--506 (1975)\par
\par
18. Katz, \(N\).: \(p\)-adic interpolation of real analytic Eisenstein series. Ann. of Math. 104, 459--571 (1976)\par
\par
19. Katz, \(N\).: \(p\)-adic L-functions for CM fields. Invent. Math. 49, 199--297 (1978)\par
\par
20. Koblitz, \(N\).: \(p\)-adic numbers, \(p\)-adic analysis, and Zeta-Functions. New York-Heidelberg-Berlin: Springer-Verlag 1977\par
\par
21. Kubert, \(D\)., Lang, \(S\).: Units in the modular function field. Lecture Notes in Math. 601, 247--275 (1977)\par
\DAnchor{60}{286}
22. Lang, \(S\).: Algebraic Number Theory. Reading, Mass.: Addison-Wesley 1970\par
\par
23. Lang, \(S\).: Introduction to Modular Forms. Berlin-Heidelberg-New York: Springer-Verlag 1976\par
\par
24. Lang, \(S\).: Cyclotomic Fields. Berlin-Heidelberg-New York: Springer-Verlag 1978\par
\par
25. Mazur, \(B\).: Analyse \(p\)-adique. Unpublished manuscript, 1973\par
\par
26. Mazur, \(B\)., Swinnerton-Dyer, H.P.F.: Arithmetic of Weil Curves. Invent. Math. 25, 1--61 (1974)\par
\par
27. Queen, \(\mathbb{C}\).: The existence of \(p\)-adic L-functions. Number Theory and Algebra, pp. 263--288. New York: Academic Press 1977\par
\par
28. Rapoport, \(M\).: Compactifications de l’espace de modules de Hilbert-Blumenthal. Compositio Math. 36, 255--335 (1978)\par
\par
29. Ribet, \(K\).: \(p\)-adic interpolation via Hilbert modular forms. Proc. Symp. Pure Math. 29, 581--592 (1975)\par
\par
30. Ribet, \(K\).: Report on \(p\)-adic L-functions over totally real fields. Astérisque 61, 177--192 (1979)\par
\par
31. Sato, \(M\)., Shintani, \(T\).: On zeta functions associated with prehomogeneous vector spaces. Ann. of Math. 100, 131--170 (1974)\par
\par
32. Serre, J-P.: Abelian \(l\)-adic representations and elliptic curves. New York: Benjamin 1968\par
\par
33. Serre, J-P.: Formes modulaires et fonctions zêta \(p\)-adiques. Lecture Notes in Math. 350, 191--268 (1973)\par
\par
34. Serre, J-P.: Sur le résidu de la fonction zêta \(p\)-adique d’un corps de nombres. C.R. Acad. Sci. Paris 287, Série \(A\), 183--188 (1978)\par
\par
35. Shimizu, \(H\).: Theta series and automorphic forms on \(\operatorname{GL}_{2}\). \(J\). Math. Soc. Japan 24, 638--683 (1972). (Correction: \(J\). Math. Soc. Japan 26, 374--376 (1974))\par
\par
36. Shintani, \(T\).: On construction of holomorphic cusp forms of half integral weight. Nagoya Math. \(J\). 58, 83--126 (1975)\par
\par
37. Shintani, \(T\).: On evaluation of zeta functions of totally real algebraic number fields at non-positive integers. \(J\). Fac. Sci. Univ. Tokyo Sect. IA 23, 393--417 (1976)\par
\par
38. Siegel, C.L.: Über die Fourierschen Koeffizienten von Modulformen. Gött. Nach. 3, 15--56 (1970)\par
\par
39. Tate, \(J\).: Fourier analysis in number fields and Hecke’s zeta-functions. Algebraic Number Theory (J.W.S. Cassels and \(A\). Fröhlich, eds.) London-New York: Academic Press 1967\par
\par
40. Weil, \(A\).: Fonction zêta et distributions. Séminaire Bourbaki 312, Juin, 1966 \((=\mathbb{C}ollected\) Papers [1966])\par
\par
41. Weil, \(A\).: Sur certains groupes d’opérateurs unitaires. Acta Math. 111, 143--211 (1964) \((=\mathbb{C}ollected\) Papers [1964b])\par
\par
42. Whittaker, E.T., Watson, G.N.: \(A\) Course of Modern Analysis. Fourth edition, reprinted. Cambridge: Cambridge University Press 1935\par
\par
Received January 7, 1980\par
\end{document}
